% ============================================================================
%  L'algorithme de résolution de PASAPAS (Cast3M) — note de lecture du source
%  Compilation : xelatex ou lualatex (deux passes), ou tectonic.
%  Polices : TeX Gyre Pagella (texte et maths), DejaVu Sans Mono.
% ============================================================================
\documentclass[11pt,a4paper]{article}

\usepackage[a4paper,margin=2.2cm,headheight=14pt]{geometry}
\usepackage{fontspec}
\usepackage{amsmath,amssymb}
\usepackage{unicode-math}
\setmainfont{texgyrepagella}[Extension=.otf,UprightFont=*-regular,BoldFont=*-bold,ItalicFont=*-italic,BoldItalicFont=*-bolditalic]
\setmathfont{texgyrepagella-math.otf}
\setsansfont{texgyreheros}[Extension=.otf,UprightFont=*-regular,BoldFont=*-bold,ItalicFont=*-italic,BoldItalicFont=*-bolditalic,Scale=0.92]
\setmonofont{DejaVuSansMono}[Extension=.ttf,BoldFont=DejaVuSansMono-Bold.ttf,Scale=0.82]
\usepackage[french]{babel}
\usepackage{xcolor}
\usepackage{booktabs,longtable,tabularx,array}
\usepackage{float}
\usepackage{needspace}
\usepackage{enumitem}
\usepackage{tikz}
\usetikzlibrary{arrows.meta,positioning,calc,decorations.pathreplacing}
\usepackage[most]{tcolorbox}
\usepackage[noend]{algpseudocode}
\usepackage{fancyhdr}
\usepackage[hidelinks]{hyperref}
\usepackage{bookmark}

\numberwithin{equation}{section}
\setlength{\parskip}{0.45em}
\setlength{\parindent}{0pt}
\setlength{\emergencystretch}{3em}
\setlist{itemsep=0.2em,topsep=0.3em}

% ---------------------------------------------------------------- couleurs
\definecolor{algframe}{HTML}{4A5D73}
\definecolor{algtitle}{HTML}{E4EAF1}
\definecolor{glohead}{HTML}{EEF1E6}
\definecolor{notecol}{HTML}{8A6D1E}
\definecolor{risk}{HTML}{9B2C2C}
\definecolor{know}{HTML}{8A5A00}
\definecolor{detail}{HTML}{4A5D73}
\definecolor{ther}{HTML}{B5532A}
\definecolor{meca}{HTML}{2F5E8C}

% ---------------------------------------------------------------- en-têtes
\pagestyle{fancy}
\fancyhf{}
\fancyhead[L]{\small\sffamily PASAPAS — lecture du source}
\fancyhead[R]{\small\sffamily\nouppercase{\leftmark}}
\fancyfoot[C]{\small\thepage}
\renewcommand{\headrulewidth}{0.3pt}

% ---------------------------------------------------------------- macros
\newcommand{\src}[1]{\texttt{\detokenize{#1}}}          % nom Gibiane
\newcommand{\pr}[1]{\textsf{#1}}                        % nom de procédure
\newcommand{\op}[1]{\textsf{\small #1}}                 % opérateur Gibiane
\newcommand{\lig}[1]{l.~#1}
\newcommand{\vc}[1]{\symbf{#1}}                         % vecteur / matrice
\newcommand{\eps}{\varepsilon}
\newcommand{\Deps}{\Delta\symbfit{\varepsilon}}
\newcommand{\sig}{\symbfit{\sigma}}
\newcommand{\vg}[1]{\symbfit{#1}}                       % grec en gras
\newcommand{\norminf}[1]{\lVert #1\rVert_{\infty}}
\newcommand{\algoref}[1]{\hyperref[algo:#1]{Algo~#1}}
\newcommand{\annref}[1]{\hyperref[annexe:#1]{Annexe~#1}}
\newcommand{\Sortir}{\textbf{sortir}}
\newcommand{\Continuer}{\textbf{continuer}}
\newcommand{\Retourner}{\textbf{retourner}}

% ---------------------------------------------------------------- pseudo-code
\algrenewcommand\algorithmicfor{\textbf{pour}}
\algrenewcommand\algorithmicforall{\textbf{pour chaque}}
\algrenewcommand\algorithmicdo{\textbf{faire}}
\algrenewcommand\algorithmicwhile{\textbf{tant que}}
\algrenewcommand\algorithmicif{\textbf{si}}
\algrenewcommand\algorithmicthen{\textbf{alors}}
\algrenewcommand\algorithmicelse{\textbf{sinon}}
\algrenewcommand\algorithmicloop{\textbf{répéter}}
\algrenewcommand\algorithmicreturn{\textbf{retourner}}
\algrenewcommand\algorithmicrequire{\textbf{Entrées :}}
\algrenewcommand\algorithmicensure{\textbf{Sorties :}}
\algrenewcommand{\algorithmiccomment}[1]{\hfill{\footnotesize\color{gray}\,#1}}
\algrenewcommand\alglinenumber[1]{\scriptsize\color{gray}#1}
% bloc « étape » : une ligne dont le contenu est indenté dessous, sans ligne de fin
\algdef{SE}[ETAPE]{Etape}{EndEtape}[1]{#1}{\algorithmicend}
\algnotext{EndEtape}

% Boîte d'algorithme : \begin{algobox}{2L}{TRANSLIN : thermique linéaire}{lignes de TRANSLIN}
\newtcolorbox{algobox}[3]{%
  breakable, enhanced, colback=white, colframe=algframe, boxrule=0.6pt,
  arc=1mm, left=4pt, right=4pt, top=3pt, bottom=3pt,
  colbacktitle=algtitle, coltitle=black, fonttitle=\sffamily\small,
  title={\textbf{Algo~#1}\enspace #2\hfill{\footnotesize\itshape #3}},
  before upper={\phantomsection\label{algo:#1}\small}}
\newtcolorbox{annexbox}[3]{%
  breakable, enhanced, colback=white, colframe=algframe, boxrule=0.6pt,
  arc=1mm, left=4pt, right=4pt, top=3pt, bottom=3pt,
  colbacktitle=algtitle, coltitle=black, fonttitle=\sffamily\small,
  title={\textbf{Annexe~#1}\enspace #2\hfill{\footnotesize\itshape #3}},
  before upper={\phantomsection\label{annexe:#1}\small}}

% Encadré de remarque
\newtcolorbox{remarque}[1][Remarque]{%
  breakable, enhanced, colback=notecol!6, colframe=notecol!70, boxrule=0.5pt,
  arc=1mm, left=6pt, right=6pt, top=4pt, bottom=4pt,
  fonttitle=\sffamily\small\bfseries, coltitle=notecol!90!black,
  attach title to upper={\enspace}, title={#1.}}

% Glossaire : trois colonnes, coupé sur plusieurs pages si besoin
\newenvironment{glossaire}[1][Glossaire]{%
  \par\smallskip\needspace{7\baselineskip}\noindent{\sffamily\bfseries\small #1}\par\nopagebreak
  \small
  \renewcommand{\arraystretch}{1.18}
  \begin{longtable}{@{}>{\raggedright\arraybackslash}p{0.17\linewidth}
                     >{\raggedright\arraybackslash}p{0.22\linewidth}
                     >{\raggedright\arraybackslash}p{0.555\linewidth}@{}}
  \toprule
  \textbf{Symbole} & \textbf{Nom dans le source} & \textbf{Contenu et usage}\\
  \midrule\endhead
  \bottomrule\endfoot}
  {\end{longtable}\addtocounter{table}{-1}\normalsize\vspace{-0.4em}}

% Constat
\newcommand{\constat}[3]{%
  \par\smallskip\noindent{\sffamily\small\bfseries\color{#1}#2}\enspace{\bfseries #3}\par\nopagebreak}

\title{\sffamily\bfseries L'algorithme de résolution de PASAPAS\\[0.3em]
  \large\mdseries Cast3M — lecture commentée du source Gibiane et du Fortran des opérateurs}
\author{}
\date{16 octobre 2026}

\begin{document}
\maketitle
\thispagestyle{empty}

\begin{abstract}
\noindent Cette note décrit ce que fait réellement la procédure \pr{PASAPAS} de Cast3M, telle qu'elle est écrite.
Elle couvre la boucle en temps, le couplage thermique-mécanique, les trois procédures thermiques et le pas de mécanique \pr{UNPAS} : itérations d'équilibre, opérateur du Newton, accélération de convergence, intégration du comportement, configurations en grands déplacements, critères de convergence et de non-convergence.
Elle traite aussi les options avancées (non local, méthode asymptotique numérique, pilotage, contact, rayonnement, rigidité augmentée, dynamique).
Chaque algorithme est écrit en pseudo-code, précédé d'un glossaire des variables du source et suivi d'explications et des équations qu'il résout.
Toutes les affirmations viennent de la lecture du source ; aucune n'a été confirmée par un calcul.
\end{abstract}

\setcounter{tocdepth}{2}
\tableofcontents
\clearpage
% ============================================================================
\section{Objet, sources et conventions}
% ============================================================================

\subsection{Objet et méthode}

\pr{PASAPAS} est la procédure générale de Cast3M pour les calculs incrémentaux : elle avance dans le temps et résout, à chaque instant, un problème mécanique non linéaire, un problème thermique, ou les deux.
Elle est écrite en Gibiane, le langage de commandes de Cast3M, et s'appuie sur une vingtaine de sous-procédures et sur des opérateurs écrits en Fortran (Esope).

Cette note décrit le code tel qu'il est écrit, et non tel que la documentation le présente.
Les procédures Gibiane ont été téléchargées sur le site de Cast3M, aux versions du tableau~\ref{tab:versions}.
Les opérateurs cités ont été vérifiés dans leur source Fortran, que le même site distribue :
\begin{itemize}
\item procédures : \url{https://www-cast3m.cea.fr/download.php?nom=pasapas&type=procedure} ;
\item Fortran : \url{https://www-cast3m.cea.fr/download.php?nom=act3&type=source}, en changeant le nom.
\end{itemize}
Les notices et les supports de formation n'ont servi qu'à orienter la lecture : ils peuvent être faux, et ils le sont sur au moins deux points (KTAN et CONF, § \ref{sec:ktan} et \ref{sec:conf}).
Les numéros de ligne cités renvoient à ces versions.
Rien de ce qui suit n'a été confirmé par un calcul : les constats du § \ref{sec:constats} méritent chacun un jeu de données d'essai.

\begin{table}[ht]
\centering\small
\caption{Procédures lues, avec la date de version portée en tête de fichier.}
\label{tab:versions}
\begin{tabularx}{\linewidth}{@{}>{\raggedright\arraybackslash}p{0.38\linewidth}l>{\raggedright\arraybackslash}X@{}}
\toprule
Procédure & Version & Rôle\\
\midrule
\pr{PASAPAS}  & 26/06/09 & boucle en temps, couplage\\
\pr{UNPAS}    & 26/09/09 & un pas de mécanique (3\,843 lignes)\\
\pr{TRANSNON} & 26/03/20 & un pas de thermique non linéaire\\
\pr{TRANSLIN} & 24/09/11 & un pas de thermique linéaire\\
\pr{DUPONT2}  & 24/09/11 & un pas de thermique, schéma de Dupont\\
\pr{PAS\_DEFA} & 26/08/05 & valeurs par défaut\\
\pr{PAS\_MODL} & 26/05/28 & classification du modèle et du matériau\\
\pr{PAS\_INIT} & 26/10/15 & configurations initiales\\
\pr{PAS\_ETAT} & 26/04/14 & état aux instants intermédiaires\\
\pr{PAS\_RESU} & 25/03/07 & transfert et sauvegarde des résultats\\
\pr{CORMAN}   & 13/02/05 & méthode asymptotique numérique\\
\pr{UNILATER} & 26/04/21 & conditions unilatérales (contact)\\
\pr{PAS\_RAYO}, \pr{HRAYO}, \pr{HRCAV} & 25/05/21 & rayonnement (\pr{HRAYO} 14/04/17, \pr{HRCAV} 12/06/07)\\
\pr{AUTOPILO}, \pr{PILOINDI}, \pr{PAS\_HELM}, \pr{DEFO\_IMP} & — & pilotage, Helmholtz, contrôle de déformation\\
\bottomrule
\end{tabularx}
\end{table}

\subsection{Lire le pseudo-code}

Chaque algorithme est écrit en pseudo-code, dans un cadre intitulé « Algo~$n$ » (ou « Annexe~X » pour les options).
Les règles de lecture sont les suivantes.

\begin{itemize}
\item \textbf{Une ligne plus indentée s'exécute à l'intérieur de la ligne au-dessus} (corps d'une boucle ou d'un test). Il n'y a pas de « fin pour » ni de « fin si » : seule l'indentation porte l'imbrication.
\item \textbf{Deux lignes alignées s'exécutent l'une après l'autre.}
\item \textbf{« Algo~$n$ » ou « Annexe~X » dans une ligne est un appel} : le détail de cette étape est écrit ailleurs, sous ce nom.
\item $a \gets b$ est une affectation ; \Sortir{} quitte la boucle la plus intérieure ; \Retourner{} quitte l'algorithme.
\item Le commentaire grisé en bout de ligne donne les lignes du source, dans \pr{UNPAS} sauf mention contraire, ou le nom de la boucle Gibiane (\src{BEXTERN}, \src{ETIQ}…).
\end{itemize}

Les schémas ne contiennent aucune flèche entre boucles : une boucle contenue dans une autre est dessinée à l'intérieur de son cadre.

\subsection{Notations}

Les vecteurs et matrices discrets sont en gras : $\vc u$ (déplacements nodaux), $\vg\lambda$ (multiplicateurs de Lagrange), $\vc A$ (matrice des liaisons), $\vc K$ (rigidité), $\vc M$ (masse), $\vc C$ (amortissement, ou capacité thermique selon le contexte), $\vc F$ (forces).
Les champs aux points d'intégration sont aussi en gras : $\sig$ (contraintes), $\vg\eps$ (déformations), $\vc V$ (variables internes).
L'indice $n$ désigne le début du pas et $n+1$ sa fin ; l'exposant $(i)$ désigne l'itéré numéro $i$.
$\Delta$ désigne un incrément cumulé depuis le début du pas et $\delta$ la correction apportée par une itération.
$\norminf{\cdot}$ est le maximum des valeurs absolues sur tous les nœuds (ou points de Gauss) et toutes les composantes retenues.
« GD » abrège « grands déplacements » et « HPP » « hypothèse des petites perturbations ».

Les équations sont de deux natures.
Celles qui traduisent une ligne du source sont accompagnées de sa référence.
Les autres sont des rappels classiques, donnés pour expliquer ce que le code calcule ; le texte le signale.

\subsection{Trois règles de Gibiane utiles à la lecture}
\label{sec:gibiane}

\paragraph{Les expressions s'évaluent de gauche à droite, sans priorité d'opérateurs.}
C'est la règle la plus importante pour traduire le source en formules, et elle est établie par le Fortran de l'interpréteur :
\begin{itemize}
\item \texttt{ANASYN} découpe l'instruction en « phrases élémentaires » : il cherche la première parenthèse fermante et remonte à l'ouvrante correspondante. La parenthèse la plus intérieure est donc évaluée d'abord.
\item \texttt{PILOT} lit ensuite les opérateurs un à un, dans l'ordre où ils apparaissent dans la phrase.
\item Chaque opérateur prend ses opérandes dans une pile (\texttt{LIRABJ}) : le premier objet disponible du type voulu, en lisant de gauche à droite.
\item Le résultat est réinséré en tête de la partie non lue de la pile (\texttt{ECPI}), où l'opérateur suivant le trouve en premier.
\end{itemize}
Ainsi \texttt{1 + 2 * 3} vaut $(1+2)\times 3 = 9$, et l'expression \texttt{psc + 1.D0 ** 0.5D0} de \pr{CORMAN} vaut $(\mathrm{psc}+1)^{1/2}$.
Toutes les formules de cette note ont été dérivées en respectant cette règle.

\paragraph{La portée des noms est dynamique.}
À l'appel d'une procédure, \texttt{INIPRO} initialise chaque nom qu'elle emploie avec l'objet de même nom des niveaux appelants, s'il en existe un.
Une procédure peut donc lire une variable de son appelant sans la recevoir en argument (§ \ref{sec:translin}).

\paragraph{Lire un indice absent d'une table est une erreur.}
\texttt{REMPIL} lève l'erreur~535 quand on lit \texttt{T.'X'} et que l'indice \texttt{'X'} n'existe pas (§ \ref{sec:constats}).

% ============================================================================
\section{Vue d'ensemble}
\label{sec:vue}
% ============================================================================

\begin{glossaire}[Glossaire — objets partagés par toutes les procédures]
— & \src{PRECED} & Table passée à \pr{PASAPAS} par l'utilisateur : modèle, matériau, chargement, options, et résultats des appels précédents. Toutes les sous-procédures la reçoivent.\\
— & \src{WTAB} = \src{PRECED.'WTABLE'} & Table de travail, recréée par \pr{PAS\_DEFA} à chaque appel. Elle porte les options après défauts, les configurations, la raideur mémorisée, les compteurs.\\
— & \src{CONTINUATION} (\src{conti}) & État convergé au début du pas : temps $t_n$, $\vc u_n$, $\sig_n$, $\vc V_n$, $T_n$\ldots\\
— & \src{ESTIMATION} (\src{estim}) & État de fin de pas en construction ($t_{n+1}$, $\vc u_{n+1}$\ldots). \pr{PAS\_RESU} le recopie dans \src{CONTINUATION} en fin de pas.\\
$\{t_k\}$ & \src{TEMPS_CALCULES} & Instants à atteindre, dans l'ordre.\\
— & \src{TEMPS_SAUVES} & Instants dont les résultats sont sauvegardés dans les tables de sortie.\\
— & \src{MODELE}, \src{CARACTERISTIQUES} & Modèle (maillage, formulation, loi) et matériau ; ils peuvent évoluer dans le temps (\src{CHAR_MODE}).\\
— & \src{CHARGEMENT} & Chargements tabulés en temps (\op{TIRE} les évalue à un instant).\\
— & \src{FOR0}, \src{GE0_DEB}, \src{FOR} & Configurations : de référence, de début de pas, de fin de pas (§ \ref{sec:configs}).\\
\end{glossaire}

\subsection{Le problème résolu}

\paragraph{Mécanique.}
Après discrétisation par éléments finis, on cherche à chaque instant $t_{n+1}$ les déplacements $\vc u_{n+1}$ et les multiplicateurs $\vg\lambda_{n+1}$ tels que
\begin{equation}
\label{eq:equilibre}
\vc F^{int}(\vc u_{n+1}) + \vc A^{\mathsf T}\vg\lambda_{n+1} = \vc F^{ext}(t_{n+1}) + \vc F^{suiv}(\vc u_{n+1}),
\qquad
\vc A\,\vc u_{n+1} = \vc u^{imp}(t_{n+1}),
\end{equation}
où $\vc F^{int} = \int_\Omega \vc B^{\mathsf T}\sig\,\mathrm d\Omega$ (opérateur \op{BSIGMA}), $\sig$ étant obtenue en intégrant la loi de comportement sur le pas, $\vc F^{suiv}$ regroupe les chargements suiveurs (pression en GD, procédure utilisateur \src{CHARMECA}), et $\vc A\vc u = \vc u^{imp}$ représente les déplacements imposés et les relations.
Les blocages sont traités par multiplicateurs de Lagrange : $-\vc A^{\mathsf T}\vg\lambda$ est la réaction.
En dynamique, on ajoute $\vc M\ddot{\vc u} + \vc C\dot{\vc u}$ au premier membre ; en consolidation, un terme de couplage avec la pression de pore.

\paragraph{Thermique.}
On cherche les températures nodales $\vc T$ telles que
\begin{equation}
\label{eq:thermique}
\vc C(\vc T)\,\dot{\vc T} + \vc K(\vc T)\,\vc T = \vc F(\vc T, t),
\end{equation}
où $\vc C$ est la capacité, $\vc K$ regroupe conduction, convection, advection et rayonnement linéarisé, et $\vc F$ les flux, sources, températures extérieures de convection et le second membre du rayonnement ; les températures imposées passent aussi par des multiplicateurs.

\paragraph{Couplage.}
La thermique agit sur la mécanique par le matériau, évalué à la température courante, et par la déformation thermique $\vg\eps^{th} = \op{EPTH}(T) - \op{EPTH}(T_{ref})$.
La mécanique n'agit sur la thermique que par la géométrie en GD et par les procédures utilisateur (§ \ref{sec:couplage}).

\subsection{Structure générale}

La figure~\ref{fig:imbrication} montre l'imbrication des boucles : chaque cadre est une boucle, contenue dans celle qui l'entoure.
Dans une passe, la thermique (cadre de gauche) s'exécute entièrement avant la mécanique (cadre de droite).

\begin{figure}[ht]
\centering
\begin{tikzpicture}[x=1cm,y=1cm,
  bx/.style={draw=algframe,rounded corners=2pt,line width=0.5pt},
  lb/.style={anchor=north west,font=\sffamily\scriptsize,inner sep=2pt,align=left}]
  \draw[bx,fill=algtitle!35] (0,0) rectangle (16,10.2);
  \node[lb,text width=15.6cm] at (0.05,10.15) {\textbf{PASAPAS} — \algoref{1}};
  \draw[bx,fill=white] (0.3,0.3) rectangle (15.7,9.55);
  \node[lb,text width=15cm] at (0.35,9.5) {pour chaque instant $T_I$ de \src{TEMPS_CALCULES} — boucle \src{BEXTERN}};
  \draw[bx,fill=algtitle!20] (0.6,0.6) rectangle (15.4,8.85);
  \node[lb,text width=14.5cm] at (0.65,8.8) {pour chaque passe $k$, une seule par défaut — boucle \src{BO_BOTH}};
  % thermique
  \draw[bx,draw=ther,fill=ther!6] (0.9,0.9) rectangle (6.7,8.1);
  \node[lb,text width=5.6cm] at (0.95,8.05) {\textbf{\textcircled{\tiny 1} thermique de $t$ à $T_I$}\\ \algoref{2}, \algoref{2L} ou \algoref{2D}};
  \draw[bx,draw=ther,fill=white] (1.2,1.2) rectangle (6.4,6.95);
  \node[lb,text width=5cm] at (1.25,6.9) {pour chaque sous-pas thermique\\ boucle \src{BOUC_SO} (si \src{PAS_MAX})};
  \draw[bx,draw=ther,fill=ther!6] (1.5,1.5) rectangle (6.1,5.8);
  \node[lb,text width=4.5cm] at (1.55,5.75) {répéter les itérations — \src{BOU_IT1}\\ (\pr{TRANSNON} seulement)};
  \node[font=\sffamily\scriptsize,align=center,text width=4.2cm] at (3.8,3.4) {$\vc S\,\Delta\vc T^{(i)} = \vc R^{(i-1)}$\\[2pt] critère sur la variation\\ de $\Delta\vc T$};
  % mécanique
  \draw[bx,draw=meca,fill=meca!6] (7,0.9) rectangle (15.1,8.1);
  \node[lb,text width=7.9cm] at (7.05,8.05) {\textbf{\textcircled{\tiny 2} mécanique de $t$ à $T_I$} — \pr{UNPAS}, \algoref{3}};
  \draw[bx,draw=meca,fill=white] (7.3,1.2) rectangle (14.8,7.35);
  \node[lb,text width=7.3cm] at (7.35,7.3) {pour chaque sous-pas $s$, 100 au plus — \src{BONOCONV} (convergence forcée)};
  \draw[bx,draw=meca,fill=meca!6] (7.6,1.5) rectangle (14.5,6.6);
  \node[lb,text width=6.8cm] at (7.65,6.55) {répéter les itérations d'équilibre $i$ — \src{ETIQ}, \algoref{4} ; tests : \algoref{6}};
  \draw[bx,draw=meca,fill=white] (7.9,1.8) rectangle (14.2,5.85);
  \node[lb,text width=6.2cm] at (7.95,5.8) {pour chaque sous-incrément $k$ — \src{sousinc}, \algoref{5}};
  \draw[bx,draw=meca,fill=meca!6] (8.2,2.1) rectangle (13.9,5.1);
  \node[lb,text width=5.6cm] at (8.25,5.05) {répéter les sous-itérations non locales — \src{BIHL}, \annref{A} (si \src{NON_LOCAL})};
  \node[font=\sffamily\scriptsize,align=center,text width=5cm] at (11.05,3.2) {loi de comportement \op{COMP}\\ puis régularisation de $\vc V$};
\end{tikzpicture}
\caption{Imbrication des boucles. Un cadre contenu dans un autre est une boucle intérieure. Dans une passe, le bloc \textcircled{\tiny 1} s'exécute avant le bloc \textcircled{\tiny 2}.}
\label{fig:imbrication}
\end{figure}

\begin{algobox}{0}{Imbrication des boucles}{PASAPAS · TRANSNON · UNPAS}
\begin{algorithmic}[1]
\Etape{\pr{PASAPAS}} \Comment{\algoref{1}}
\ForAll{instant $T_I$} \Comment{boucle \src{BEXTERN}}
  \ForAll{passe $k$, une seule par défaut} \Comment{boucle \src{BO_BOTH}}
    \Etape{thermique de $t$ à $T_I$} \Comment{\algoref{2}, \algoref{2L} ou \algoref{2D}}
      \ForAll{sous-pas thermique} \Comment{\src{BOUC_SO}, si \src{PAS_MAX}}
        \State répéter les itérations \Comment{\src{BOU_IT1}}
      \EndFor
    \EndEtape
    \Etape{puis mécanique de $t$ à $T_I$} \Comment{\algoref{3}, \pr{UNPAS}}
      \ForAll{sous-pas $s$, 100 au plus} \Comment{\src{BONOCONV}, convergence forcée}
        \Etape{répéter les itérations $i$} \Comment{\src{ETIQ}, \algoref{4}}
          \ForAll{sous-incrément $k$} \Comment{\src{sousinc}, \algoref{5}}
            \State répéter les sous-itérations non locales \Comment{\src{BIHL}, \annref{A}}
          \EndFor
        \EndEtape
      \EndFor
    \EndEtape
  \EndFor
\EndFor
\EndEtape
\end{algorithmic}
\end{algobox}

\subsection{L'essentiel en six points}

\begin{enumerate}
\item \textbf{Boucle en temps.} On parcourt une liste d'instants imposés (\src{TEMPS_CALCULES}). Si l'on ne donne que \src{TEMPS_SAUVES}, le pas s'adapte au nombre d'itérations : divisé par $\sqrt2$ au-delà de 20, multiplié par $\sqrt2$ en deçà de 7 (§ \ref{sec:temps}).
\item \textbf{Couplage étagé.} La thermique avance de $t_n$ à $t_{n+1}$, puis la mécanique utilise $T_{n+1}$. Il n'y a qu'une passe par défaut (§ \ref{sec:couplage}).
\item \textbf{Newton modifié, accéléré.} L'opérateur des itérations est par défaut la rigidité élastique (plus $\vc K_\sigma$ en GD), gardée d'un pas à l'autre quand c'est possible. L'accélération \op{ACT3} intervient une itération sur deux. La tangente n'est utilisée qu'avec \src{K_TANGENT}, et ce n'est pas la tangente algorithmique (§ \ref{sec:operateur}).
\item \textbf{Deux critères à $10^{-4}$.} Le résidu d'équilibre, en norme max et rapporté à une force de référence fixée en début de pas, et la variation de l'incrément de déformation entre deux itérations ; un troisième critère porte sur les moments (§ \ref{sec:criteres}).
\item \textbf{Échec.} La non-convergence est détectée quand le critère n'a pas baissé de 1\,\% en 10 itérations, ou après \src{MAXITERATION} itérations. Par défaut, on repart alors de l'itéré non convergé comme nouvel état initial (convergence forcée), jusqu'à 100 sous-pas (§ \ref{sec:pas}).
\item \textbf{Configurations.} En HPP, tout se fait sur la configuration initiale. En GD : début de pas pour la rigidité initiale et $\vc F^{int}_n$, mi-incrément pour $\Deps$ et la loi de comportement, fin de pas recalculée à chaque itération pour les forces internes (§ \ref{sec:configs}).
\end{enumerate}
% ============================================================================
\section{La boucle en temps et le couplage}
\label{sec:temps}
% ============================================================================

\begin{glossaire}
$t$ & \src{TEMP0} & Instant de début du pas (lu dans \src{conti.'TEMPS'}).\\
$T_I$ & \src{TI} & Instant visé par le pas en cours.\\
$k_I$ & \src{ICHG} (\src{WTAB.'ICHG'}) & Rang de l'instant visé dans \src{TEMPS_CALCULES}. On le décrémente pour « rejouer » une cible non atteinte.\\
$\Delta t$ & \src{WTAB.'DT'} & Durée du pas, $T_I - t$.\\
$\Delta t_{prec}$ & \src{WTAB.'DTPREC'} & Durée du pas précédent, utilisée en dynamique et par le prédicteur.\\
$\Delta t_{ref}$ & \src{DT_AVANT} & Pas de référence de \src{PAS_AJUSTE}, ajusté d'un pas à l'autre.\\
$T_I^{\,cible}$ & \src{TIV} & Cible de \src{PAS_AJUSTE} : l'instant de \src{TEMPS_CALCULES} qu'on cherche à atteindre.\\
$N_{it}$ & \src{NOMBRE_ITERATIONS} & Somme des itérations des sous-pas convergés du pas précédent (\pr{UNPAS} \lig{3420}).\\
— & \src{IREREDU} & Faux juste après un pas non convergé : interdit alors d'agrandir le pas.\\
— & \src{PASFINAL} & Compte les pas consécutifs qui atteignent directement la cible.\\
$N_{both}$ & \src{BOU_METH} (\src{NB_BOTH}) & Nombre maximal de passes thermique-mécanique : 1 par défaut, 10 avec \src{CONVERGENCE_MEC_THE}.\\
$k$ & \src{I_BOTH} & Numéro de la passe.\\
$\vc T_k$ & \src{THER_COURANT} & Température obtenue par la thermique de la passe $k$ ; \src{THER_COURANTP} garde celle de la passe précédente.\\
$c_{coh}$ & \src{CRITERE_COHERENCE} & Seuil de cohérence thermique-mécanique, \src{PRECISION} par défaut.\\
— & \src{CONV_THMEC} & Vrai quand on peut quitter la boucle des passes ; vrai par défaut.\\
$\vc T_{n-1}$ & \src{TET0} & Température « deux niveaux avant », pour \pr{DUPONT2}.\\
$\vc T_n$, $\vc T_{n+1}$ & \src{TET1}, \src{TET2} & Températures de début et de fin de pas vues par la mécanique.\\
$\vg\eps^{th}_n$, $\vg\eps^{th}_{n+1}$ & \src{ETHER1}, \src{ETHER2} & Déformations thermiques de début et de fin de pas, comptées depuis $T_{ref}$.\\
$\vg\eps^{th}_{ref}$ & \src{ETREF} & $\op{EPTH}(T_{ref})$, calculée une fois.\\
— & \src{DEFOR1}, \src{DEFOR2} & Déformations imposées (\src{DEFI}) en début et fin de pas.\\
— & \src{FREA1}, \src{FOPL} & Dynamique : partie du second membre héritée du pas précédent, et déséquilibre de fin de pas (§ \ref{sec:dynamique}).\\
$\alpha$ & \src{AUTOCOEF} & Fraction du pas effectivement faite : dernier coefficient $\alpha$ (\src{COEPI}) d'\pr{UNPAS} (\lig{3808}).\\
— & \src{WTAB.'CONV'} & Vrai si le dernier pas mécanique a convergé.\\
— & \src{REEQUILIBRAGE} & Option : autorise un pas de durée nulle pour rééquilibrer.\\
— & \src{REA_GEOM}, \src{RECARI} & Drapeaux : recalculer la géométrie du rayonnement ; recalculer la rigidité à chaque pas.\\
\end{glossaire}

\subsection{Initialisation}

\pr{PASAPAS} commence par deux procédures.
\pr{PAS\_DEFA} crée une table de travail neuve (\src{TWORK = 'TABLE'}, \lig{102} de \pr{PAS\_DEFA}), y range les valeurs par défaut puis les options de l'utilisateur, et ne recopie de l'appel précédent qu'une liste fermée d'objets : la configuration de référence, le prédicteur, les données du pilotage et de la dynamique, la dernière tangente, \src{NOMBRE_ITERATIONS}, le modèle et le matériau (\pr{PAS\_DEFA} \lig{611-660}).
Tout le reste repart de zéro à chaque appel ; c'est le cas de la rigidité mémorisée \src{RRRR} et des données de \pr{DUPONT2} (\src{TET0}, \src{VIEUXPAS}).
\pr{PAS\_INIT} crée les tables de résultats et fixe les configurations (\pr{PAS\_INIT} \lig{402-411}) :
\begin{equation}
\mathrm{FOR0} = \text{configuration au premier appel},\qquad
\mathrm{GE0\_DEB} =
\begin{cases}
\mathrm{FOR0} & \text{en HPP},\\
\mathrm{FOR0} + \vc u_0 & \text{en GD, FEFP.}
\end{cases}
\end{equation}

Viennent ensuite le matériau à l'instant initial, $\vc D_0 = \pr{PAS\_MATE}(\pr{PAS\_ETAT}(t_0))$, et la déformation thermique initiale
\begin{equation}
\vg\eps^{th}_{ref} = \op{EPTH}\big(\vc D_{ref}, T_{ref}\big),\qquad
\vg\eps^{th}_0 = \op{EPTH}\big(\vc D_0, T_0\big) - \vg\eps^{th}_{ref}\qquad(\text{\lig{50-76}}),
\end{equation}
où $\vc D_{ref}$ est le matériau évalué à la température de référence $T_{ref}$.
La déformation thermique est donc toujours comptée depuis $T_{ref}$, et le coefficient de dilatation est pris au matériau de l'instant considéré.

\subsection{Choix de l'instant suivant}

À chaque tour de la boucle \src{BEXTERN}, on prend l'instant suivant $T_I$ de \src{TEMPS_CALCULES}.
On le saute s'il est antérieur à $t$, s'il coïncide avec l'instant initial, ou s'il vaut $t$ sans l'option \src{REEQUILIBRAGE} (\lig{105-119}).

\paragraph{Pas adaptatif (\src{PAS_AJUSTE}).}
Quand l'utilisateur ne donne que \src{TEMPS_SAUVES}, \pr{PAS\_DEFA} active \src{PAS_AJUSTE} et prend ces instants comme cibles (\pr{PAS\_DEFA} \lig{803-806}).
Le pas suit alors le nombre d'itérations du pas précédent (\lig{123-158}) :
\begin{equation}
\Delta t_{ref} \gets
\begin{cases}
\Delta t_{ref}/\sqrt2 & \text{si le pas précédent a convergé et } N_{it} > 20,\\
\Delta t_{ref}\,\sqrt2 & \text{si le pas précédent a convergé, } N_{it} < 7 \text{ et le pas d'avant n'a pas échoué},\\
\Delta t_{ref} & \text{sinon, s'il a convergé.}
\end{cases}
\end{equation}
Après un pas non convergé, on fait un pas quasi nul, $\Delta t = 10^{-10}\,\Delta t_{ref}$, puis on repart avec $\Delta t_{ref}/2$.
Ce pas quasi nul rééquilibre vraisemblablement l'état non convergé que \pr{PAS\_RESU} a recopié (voir plus bas) ; c'est une interprétation, le source ne le commente pas.
L'instant visé devient $T_I = t + 1{,}0000000001\,\Delta t$ s'il ne dépasse pas la cible, et la cible sinon.
Quand la cible est atteinte directement deux pas de suite, le pas de référence devient le pas réellement fait, $T_I - t$.

$N_{it}$ ne compte que les itérations des sous-pas \emph{convergés} : un pas sauvé par la convergence forcée, dont le dernier sous-pas converge vite, peut voir son pas augmenter (constat~\ref{c:ajuste}, § \ref{sec:constats}).

\paragraph{Durée du pas.}
$\Delta t = T_I - t$. En dynamique et en consolidation, un pas de durée nulle est sauté, et $\Delta t$ reprend la valeur du pas précédent (\lig{153-158}).

\paragraph{Modèle évolutif (\src{CHAR_MODE}).}
Si le modèle change avec le temps, on prend le modèle (et éventuellement le matériau) à $T_I$, on relance \pr{PAS\_MODL}, et l'on initialise température et déformations des zones ajoutées (\lig{165-336}).

\subsection{La boucle des passes et l'appel des physiques}
\label{sec:passes}

Chaque instant est traité par une boucle de passes (\src{BO_BOTH}).
Dans une passe, on résout d'abord la thermique, puis la mécanique.
Par défaut il n'y a qu'une passe : c'est un couplage \emph{étagé}.

\paragraph{Thermique.}
Elle n'est appelée que si $T_I \neq t$ (\lig{375}).
La géométrie est celle de début de pas (\src{GE0_DEB}) à la première passe ; aux passes suivantes en GD, c'est la configuration de fin de pas (\src{FOR}), et l'on met \src{RECARI} et \src{REA_GEOM} à vrai (\lig{361-372}).
Si \src{PAS_MAX} est donné et que $t + \src{PAS_MAX} < T_I$, l'intervalle est découpé en sous-pas thermiques de longueur \src{PAS_MAX} ; sinon il n'y a qu'un sous-pas.
La procédure appelée dépend de \src{PROCEDURE_THERMIQUE} : \pr{TRANSNON} (\src{NONLINEAIRE}, défaut, \algoref{2}), \pr{TRANSLIN} (\src{LINEAIRE}, \algoref{2L}) ou \pr{DUPONT2} (\src{DUPONT}, \algoref{2D}) (\lig{402-412}).

\paragraph{Mécanique.}
On se place sur \src{GE0_DEB} et l'on appelle \pr{UNPAS} (\algoref{3}).
Si le pas n'a pas convergé et que la convergence forcée est désactivée, \pr{PASAPAS} s'arrête sur l'erreur~997 (\lig{440-448}).

\paragraph{Cohérence (\src{CONVERGENCE_MEC_THE}).}
Avec cette option, on recommence les passes jusqu'à ce que la température ne bouge plus.
À partir de la deuxième passe, on calcule, inconnue par inconnue (composante $c$ : \src{T}, concentrations…),
\begin{equation}
e_{coh} = \max_c\;\frac{\max_{\text{nœuds}}\big|T_{k,c} - T_{k-1,c}\big|}{\max_{\text{nœuds}}\big|T_{k,c}\big|}
\qquad(\text{\lig{459-494}}),
\end{equation}
le dénominateur étant remplacé par 1 s'il est négligeable devant le numérateur.
On sort si $e_{coh} < c_{coh}$ ; à la dernière passe autorisée, on s'arrête sur l'erreur~999.

\subsection{Fin du pas}

\begin{itemize}
\item \textbf{\pr{DUPONT2}} : \src{TET0} reçoit la température de fin de pas et \src{VIEUXPAS} la durée du premier sous-pas (\lig{509-514}). Ce choix fait dégénérer le schéma de Dupont d'un pas à l'autre (§ \ref{sec:dupont}).
\item \textbf{Mécanique} : $\vc T_n \gets \vc T_{n+1}$, $\vg\eps^{th}_n \gets \vg\eps^{th}_{n+1}$, $\mathrm{DEFOR1} \gets \mathrm{DEFOR2}$ ; en dynamique, $\mathrm{FREA1} \gets \mathrm{FOPL}$ (\lig{527-550}). En GD, ou en FEFP réactualisé, la configuration de début du pas suivant devient la configuration de fin : $\mathrm{GE0\_DEB} \gets \mathrm{FOR}$ (\lig{553-557}).
\item \textbf{Temps atteint} : en pilotage, ou si le pas n'a pas convergé, $t \gets t + \alpha\,(T_I - t)$, sinon $t \gets T_I$ (\lig{561-569}). Hors pilotage, $\alpha$ vaut \src{RELAXATION_NONCONV}, donc 1 par défaut.
\item \textbf{Rigidité mémorisée} : elle est oubliée si \src{RECARI} est vrai (\lig{589-591}), et dans tous les cas en fin d'appel (\lig{644-645}).
\item \textbf{\pr{PAS\_RESU}} recopie \emph{toujours} \src{ESTIMATION} dans \src{CONTINUATION}, même si le pas n'a pas convergé, puis sauvegarde les résultats si le pas a convergé et que l'instant figure dans \src{TEMPS_SAUVES} (\pr{PAS\_RESU} \lig{19-42}).
\item \textbf{Rejeu} : en pilotage, si la cible n'est pas atteinte, on rejoue le même instant (erreur~1000 après \src{AUTOPAS} essais, \lig{601-617}) ; avec \src{PAS_AJUSTE}, on rejoue la cible tant qu'elle n'est pas atteinte ou si le pas n'a pas convergé (\lig{621-625}).
\end{itemize}

\begin{algobox}{1}{PASAPAS : boucle en temps et couplage}{lignes de PASAPAS}
\begin{algorithmic}[1]
\Require modèle, matériau, chargement, options (table \src{PRECED})
\State \pr{PAS\_DEFA}, \pr{PAS\_INIT} : défauts ; $\mathrm{FOR0} \gets$ configuration courante \Comment{\lig{23} ; \pr{PAS\_INIT} \lig{402-411}}
\State $\mathrm{GE0\_DEB} \gets \mathrm{FOR0}$, ou $\mathrm{FOR0} + \vc u_0$ en GD et en FEFP
\State matériau à $t_0$ ; $\vg\eps^{th}_{ref} \gets \op{EPTH}(T_{ref})$ ; $\vg\eps^{th}_n \gets \op{EPTH}(T_0) - \vg\eps^{th}_{ref}$ \Comment{\lig{50-76}}
\ForAll{instant $T_I$ de \src{TEMPS_CALCULES}} \Comment{\src{BEXTERN}, \lig{101}}
  \State \textbf{si} $T_I < t$, ou $T_I = t_0$, ou $T_I = t$ sans \src{REEQUILIBRAGE} : \Continuer \Comment{\lig{105-119}}
  \If{\src{PAS_AJUSTE}} \Comment{\lig{123-158}}
    \If{le pas précédent a convergé}
      \State \textbf{si} $N_{it} > 20$ : $\Delta t \gets \Delta t/\sqrt2$
      \State \textbf{si} $N_{it} < 7$, sauf juste après un échec : $\Delta t \gets \Delta t\,\sqrt2$
    \Else
      \State pas quasi nul ($10^{-10}\,\Delta t$), puis repartir avec $\Delta t/2$
    \EndIf
    \State $T_I \gets t + \Delta t$, sans dépasser la cible
  \EndIf
  \State $\Delta t \gets T_I - t$ ; si $\Delta t = 0$ en dynamique ou en consolidation : \Continuer \Comment{\lig{153-158}}
  \State \textbf{si} le modèle évolue (\src{CHAR_MODE}) : modèle à $T_I$ ; $T$ et $\vg\eps$ des zones ajoutées \Comment{\lig{165-336}}
  \For{$k = 1 \dots N_{both}$} \Comment{\src{BO_BOTH}, \lig{349-499}}
    \State \textbf{si} \src{CONVERGENCE_MEC_THE} : $\src{CHARGEMENT} \gets \src{CHARGEMENT_I}$
    \If{le modèle a une thermique et $T_I \neq t$} \Comment{\lig{360, 375}}
      \State géométrie $\gets$ \src{GE0_DEB} si $k = 1$ ; \src{FOR} si $k \ge 2$ en GD \Comment{\lig{361-372}}
      \State $\vc T_{n+1} \gets$ \algoref{2} (\pr{TRANSNON}), \algoref{2L} (\pr{TRANSLIN}) ou \algoref{2D} (\pr{DUPONT2}) \Comment{\lig{402-412}}
      \State \pr{REEV\_THE} \Comment{procédure utilisateur}
    \EndIf
    \If{le modèle a une mécanique}
      \State géométrie $\gets$ \src{GE0_DEB} \Comment{\lig{437}}
      \State appeler \algoref{3} ; si échec sans convergence forcée : erreur 997 \Comment{\lig{440-448}}
      \State \pr{REEV\_MEC} \Comment{procédure utilisateur}
    \EndIf
    \If{\src{CONVERGENCE_MEC_THE}}
      \State \textbf{si} $k \ge 2$ et $e_{coh} < c_{coh}$ : \Sortir \Comment{\lig{459-494}}
      \State \textbf{si} $k = N_{both}$ : erreur 999
    \Else
      \State \Sortir \Comment{une seule passe}
    \EndIf
  \EndFor
  \State \textbf{si} \pr{DUPONT2} : $\src{TET0} \gets \vc T_{n+1}$ ; $\src{VIEUXPAS} \gets \Delta t$ \Comment{\lig{509-514}}
  \State $\vc T_n \gets \vc T_{n+1}$ ; $\vg\eps^{th}_n \gets \vg\eps^{th}_{n+1}$ ; $\mathrm{DEFOR1} \gets \mathrm{DEFOR2}$ ; $\mathrm{FREA1} \gets \mathrm{FOPL}$ \Comment{\lig{527-550}}
  \State \textbf{si} GD, ou FEFP réactualisé : $\mathrm{GE0\_DEB} \gets \mathrm{FOR}$ \Comment{\lig{553-557}}
  \State \textbf{si} pilotage ou non convergé : $t \gets t + \alpha\,(T_I - t)$ ; sinon $t \gets T_I$ \Comment{\lig{561-569}}
  \State \pr{PAS\_RESU} : $\src{CONTINUATION} \gets \src{ESTIMATION}$ ; sauvegarde si convergé \Comment{\lig{598}}
  \State \textbf{si} pilotage et cible non atteinte : rejouer la cible ; erreur 1000 après \src{AUTOPAS} essais \Comment{\lig{601-617}}
  \State \textbf{si} \src{PAS_AJUSTE} et ($T_I$ avant la cible ou non convergé) : rejouer la cible \Comment{\lig{621-625}}
\EndFor
\end{algorithmic}
\end{algobox}

\subsection{Le couplage thermique-mécanique}
\label{sec:couplage}

\paragraph{De la thermique vers la mécanique.}
La mécanique reçoit la température de fin de pas $\vc T_{n+1}$, projetée par \op{PROI} si les maillages thermique et mécanique diffèrent (\pr{UNPAS} \lig{352-375}).
Elle l'utilise de trois façons.
\begin{enumerate}
\item \textbf{Matériau à $T_I$.} \pr{PAS\_ETAT} lit $\vc T_{n+1}$ dans \src{ESTIMATION}, et \pr{PAS\_MATE} évalue le matériau. Si le matériau dépend d'autres variables que le seul coefficient de dilatation, \pr{PAS\_MODL} pose \src{MATVAR}, et la rigidité est recalculée à chaque pas (\pr{PAS\_MODL} \lig{345-350}).
\item \textbf{Pseudo-chargement thermique.} L'incrément de déformation thermique
\begin{equation}
\Delta\vg\eps^{th} = \vg\eps^{th}_{n+1} - \vg\eps^{th}_n = \op{EPTH}(\vc D_{n+1}, T_{n+1}) - \op{EPTH}(\vc D_n, T_n)
\end{equation}
entre au second membre sous forme de pseudo-force $\vc B^{\mathsf T}\vc D_{n+1}\,\Delta\vg\eps^{th}$, et il est retiré de $\Deps$ avant l'intégration de la loi (\algoref{3}, \algoref{5}).
\item \textbf{Température dans la loi.} \op{COMP} reçoit, à chaque sous-incrément $k$ sur $N$, la température interpolée linéairement $T(t_k) = T_n + \tfrac kN (T_{n+1} - T_n)$.
\end{enumerate}

\paragraph{De la mécanique vers la thermique.}
L'influence est plus faible.
\begin{itemize}
\item \textbf{Géométrie} : configuration de début de pas à la passe~1, configuration de fin de pas aux passes suivantes, en GD seulement.
\item \textbf{Champs mécaniques} : \pr{TRANSNON} les lit dans \src{CONTINUATION}, donc au début du pas (\pr{TRANSNON} \lig{312-344}), et lève une erreur si le matériau thermique en dépend alors que \src{RELAXATION_THETA} n'est pas nul (\lig{635-640}). \pr{DUPONT2} refuse ce cas dans tous les cas (§ \ref{sec:dupont}).
\item \textbf{Procédures utilisateur} : \pr{REEV\_THE}, \pr{REEV\_MEC}, \pr{CHARTHER}, \pr{PARATHER}.
\end{itemize}

\begin{remarque}[Conséquence]
En HPP et sans procédure utilisateur, la deuxième passe de \src{CONVERGENCE_MEC_THE} redonne exactement la même température : la cohérence est atteinte d'office, et le calcul coûte deux fois plus.
Cette option coupe en outre la convergence forcée et le pilotage, et elle est incompatible avec \src{PAS_AJUSTE} (\pr{PAS\_DEFA} \lig{985-999}).
Avec cette option, la table \src{WTAB.MAT1}, qui porte le matériau de « début de pas », reçoit dès la passe~2 le matériau de fin de pas (\pr{UNPAS} \lig{283}, \lig{3792}) ; cela touche la correction « matériau variable » de \algoref{3}, $\vc F^{int}_n$ des éléments de structure et l'état initial passé à \op{COMP}.
\end{remarque}
% ============================================================================
\section{La thermique}
\label{sec:thermique}
% ============================================================================

\pr{PASAPAS} dispose de trois procédures pour avancer la thermique de $t$ à $T_I$.
\src{PROCEDURE_THERMIQUE} les choisit : \src{NONLINEAIRE} (défaut) appelle \pr{TRANSNON}, \src{LINEAIRE} appelle \pr{TRANSLIN}, \src{DUPONT} appelle \pr{DUPONT2}.
\pr{PAS\_DEFA} impose \pr{TRANSNON} en cas de changement de phase, de solubilité, de diffusion ou de modèle THM (\pr{PAS\_DEFA} \lig{1203-1206}), mais pas en présence de sources de chaleur, que seul \pr{TRANSNON} sait traiter.

\subsection{Le problème discret et le schéma en temps}

Les trois procédures intègrent l'équation~\eqref{eq:thermique} sur un sous-pas $[t_0, t_1]$ de durée $\Delta t$.
\pr{TRANSNON} et \pr{TRANSLIN} utilisent un $\theta$-schéma.
Avec $\vc T_0$ la température de début de sous-pas, $\vc T_1$ celle de fin, et
\begin{equation}
\vc T_\theta = \theta\,\vc T_1 + (1-\theta)\,\vc T_0,\qquad t_\theta = \theta\,t_1 + (1-\theta)\,t_0,
\end{equation}
il s'agit de trouver $\vc T_1$ tel que
\begin{equation}
\label{eq:theta}
\vc C(\vc T_\theta)\,\frac{\vc T_1 - \vc T_0}{\Delta t} + \vc K(\vc T_\theta)\,\vc T_\theta = \vc F(\vc T_\theta, t_\theta).
\end{equation}
Dans \pr{PASAPAS}, $\theta = \src{RELAXATION_THETA} = 1$ (\pr{PAS\_DEFA} \lig{517}) : c'est le schéma d'Euler implicite, alors que l'en-tête de \pr{TRANSNON} annonce $\theta = 0{,}5$.
$\theta$ est de toute façon forcé à 1 en régime stationnaire et en présence d'advection (\pr{TRANSNON} \lig{428-429}).

\subsection{TRANSNON : point fixe à matrice constante}
\label{sec:transnon}

\begin{glossaire}[Glossaire — TRANSNON]
$\vc T_0$ & \src{U0} & Température de début de sous-pas (\src{CONTI.'TEMPERATURES'}, plus les concentrations en diffusion), décalée de $T_{ref} = 273{,}15$ si \src{CELSIUS} : \pr{TRANSNON} travaille alors en kelvins.\\
$\vc T^{(i)}$ & \src{U1} & Itéré courant, $\vc T_0 + \Delta\vc T^{(i)}$.\\
$\Delta\vc T^{(i)}$, $\Delta\vc T^{(i-1)}$ & \src{DU1}, \src{DU0} & Incrément total depuis le début du sous-pas, à l'itération courante et à la précédente.\\
$\theta$ & \src{LAMBDA1} & \src{RELAXATION_THETA} ; \src{I_LAMBD1} repère les cas $\theta = 0$, $\tfrac12$, $1$.\\
$\vc T_\theta$ & \src{U_BOU1} & $\theta\,\vc T^{(i)} + (1-\theta)\,\vc T_0$ ; on range $\vc T_\theta - T_{ref}$ dans \src{ETAB.'THER_COURANT'} (\src{TH_COUR}), que lisent le matériau, le rayonnement et \pr{PARATHER}.\\
$t_0$, $t_1$, $t_\theta$, $\Delta t$ & \src{TEM0}, \src{TEM1}, \src{TEMM}, \src{T_PAS} & Bornes et instant intermédiaire du sous-pas.\\
$\vc K$ & \src{KDETO} & Conduction, convection, advection, diffusion et conductivité constante, assemblées avec le matériau à $\vc T_0$.\\
— & \src{BCONS}, \src{BB} & Parties de $\vc K$ qui ne dépendent pas de la température ; \src{BB} y ajoute, à chaque itération, les parties qui en dépendent, recalculées à $\vc T_\theta$, et le rayonnement : c'est la matrice du résidu.\\
$\vc C/\Delta t$ & \src{CDET} & Capacité divisée par le pas.\\
$\vc S$ & \src{AA} & Matrice d'itération : $\theta\,\vc K(\vc T_0) + \vc C(\vc T_0)/\Delta t$, plus $\theta\,\vc K_r$ (rayonnement), $\theta$ fois la matrice de \pr{CHARTHER}, les blocages et les conditions de changement de phase. Elle est gardée tout le sous-pas.\\
— & \src{CCONS}, \src{CC} & Partie de $\vc S$ sans blocages ; \src{CC} lui retire $\vc C(\vc T_\theta)/\Delta t$ quand la capacité dépend de $T$.\\
$\vc F$ & \src{SECON} & Second membre indépendant de l'itération : températures imposées (à $t_1$), flux et sources (à $t_\theta$), convection si elle ne dépend pas de $T$.\\
$\vc R^{(i)}$ & \src{MAT_CHPO} & Second membre de l'itération $i$.\\
— & \src{MAT_BLOC} & Blocages thermiques (multiplicateurs).\\
$\epsilon_T$ & \src{ZPREC1} & \src{PRECISION} ($10^{-4}$) : seuil du critère relatif.\\
— & \src{TTOL} & Tolérance absolue sur la température ; si elle est donnée, elle remplace le critère relatif.\\
$i$ & \src{ITE_IT} & Numéro d'itération ; limite : \src{MAXITERATION} de la mécanique.\\
— & \src{ACFP}$k$, \src{ACFEP}$k$, \src{CORREC} & Mémoire et correction de l'accélération \op{ACT3} (§ \ref{sec:act3}).\\
\end{glossaire}

\paragraph{Le résidu.}
On écrit l'équation~\eqref{eq:theta} sous forme d'un résidu en l'inconnue $\Delta\vc T = \vc T_1 - \vc T_0$ :
\begin{equation}
\label{eq:residu-th}
\vc r(\Delta\vc T) = \vc F(\vc T_\theta, t_\theta) - \vc K(\vc T_\theta)\,\vc T_\theta - \vc C(\vc T_\theta)\,\frac{\Delta\vc T}{\Delta t},
\qquad \vc T_\theta = \vc T_0 + \theta\,\Delta\vc T .
\end{equation}

\paragraph{L'itération.}
La matrice d'itération est assemblée une fois par sous-pas, à la température de début :
\begin{equation}
\vc S = \theta\,\vc K(\vc T_0) + \frac{\vc C(\vc T_0)}{\Delta t} \qquad(\text{\lig{726-982}}).
\end{equation}
L'itération est un point fixe sur l'incrément total :
\begin{equation}
\label{eq:pointfixe}
\vc S\,\Delta\vc T^{(i)} = \vc R^{(i-1)},
\qquad
\vc R^{(i-1)} = \vc r\big(\Delta\vc T^{(i-1)}\big) + \vc S\,\Delta\vc T^{(i-1)},
\end{equation}
c'est-à-dire $\Delta\vc T^{(i)} = \Delta\vc T^{(i-1)} + \vc S^{-1}\,\vc r(\Delta\vc T^{(i-1)})$.
On part de $\Delta\vc T^{(0)} = 0$, donc $\vc R^{(0)} = \vc F - \vc K(\vc T_0)\,\vc T_0$ (\lig{1026-1027}).
Le source assemble $\vc R^{(i)}$ à partir des pièces du glossaire (\lig{1220-1468}) :
\begin{equation}
\vc R^{(i)} = \vc F + \big[\vc S - \vc C(\vc T_\theta)/\Delta t\big]\,\Delta\vc T^{(i)} - \vc K(\vc T_\theta)\,\big(\vc T_0 + \theta\,\Delta\vc T^{(i)}\big),
\end{equation}
qui est bien $\vc r + \vc S\,\Delta\vc T^{(i)}$.
Si $\vc K$ et $\vc C$ ne dépendent pas de la température, $\vc S$ est la jacobienne exacte et la deuxième itération retrouve la première.
Sinon, $\vc S$ ignore à la fois la variation de $\vc K$ et $\vc C$ pendant le pas et les termes $\partial\vc K/\partial T\cdot\vc T$ : c'est un Newton modifié, de convergence linéaire.

\paragraph{Le matériau.}
À chaque itération, le matériau est réévalué à $\vc T_\theta$ et à $t_\theta$ (\op{VARI}, et \pr{PARATHER} s'il y en a une), et les parties de $\vc K$ et $\vc C$ qui dépendent de la température sont réassemblées pour le résidu seulement.
Le rayonnement est relinéarisé à $\vc T_\theta$ pour le résidu (\annref{E}).

\paragraph{Le critère.}
Il porte sur la variation de l'incrément entre deux itérations, inconnue par inconnue :
\begin{equation}
\label{eq:critere-th}
e^{(i)} = \max_c\;\frac{\lVert \Delta\vc T^{(i)} - \Delta\vc T^{(i-1)}\rVert_{\infty,c}}{\tfrac12\big(\lVert \vc T_0\rVert_{\infty,c} + \lVert \vc T^{(i)}\rVert_{\infty,c}\big)} \le \epsilon_T
\qquad(\text{\lig{1187-1217}}).
\end{equation}
Avec \src{TTOL}, on exige à la place $|\Delta T^{(i)} - \Delta T^{(i-1)}| \le \src{TTOL}$ en tout nœud.
Comme $\Delta\vc T^{(i)} - \Delta\vc T^{(i-1)} = \vc S^{-1}\vc r^{(i-1)}$, c'est le résidu mesuré en température par l'opérateur figé.
En Celsius, les normes portent sur des kelvins, ce qui rend le critère relatif plus lâche qu'en Celsius pur.
Au-delà de \src{MAXITERATION} itérations (celui de la mécanique, puisque \pr{TRANSNON} lit la même table de travail), le calcul s'arrête sur « LE CALCUL N A PAS CONVERGE » (\lig{1059-1065}).

\paragraph{L'accélération.}
Aux itérations paires à partir de la quatrième (la sixième en cas de changement de phase), on retranche du second membre la correction d'\op{ACT3} (\lig{1084-1089}).
Les champs duaux sont les incréments du second membre, $\vc R^{(i)} - \tilde{\vc R}^{(i-1)}$, qui valent les résidus $\vc r^{(i)}$ ; ils ne sont pas pondérés, contrairement à la mécanique (§ \ref{sec:act3}).

\begin{algobox}{2}{TRANSNON : un pas de thermique}{lignes de TRANSNON}
\begin{algorithmic}[1]
\Require $\vc T_n$ (\src{CONTINUATION}), intervalle $[t, T_I]$
\ForAll{sous-pas thermique $[t_0, t_1]$} \Comment{\src{BOUC_SO}, \lig{572} ; découpé si \src{PAS_MAX}}
  \State $t_\theta \gets \theta\,t_1 + (1-\theta)\,t_0$ \Comment{$\theta = \src{RELAXATION_THETA} = 1$}
  \State matériau $\gets f(\vc T_0, t_\theta$, champs mécaniques de début de pas$)$ \Comment{\lig{581-700}}
  \State $\vc S \gets \theta\,\vc K(\vc T_0) + \vc C(\vc T_0)/\Delta t$ + blocages + $\theta\,\vc K_r(\vc T_0)$ (\annref{E}) \Comment{\lig{726-982} ; gardée tout le sous-pas}
  \State $\vc F \gets$ flux, convection, sources, $T$ imposées, rayonnement, \pr{CHARTHER} \Comment{\lig{884-1031}}
  \State $\vc R \gets \vc F - \vc K(\vc T_0)\,\vc T_0$
  \For{$i = 1, 2, \dots$} \Comment{\src{BOU_IT1}, \lig{1052}}
    \State \textbf{si} $i > \src{MAXITERATION}$ : « LE CALCUL N A PAS CONVERGE » ; erreur \Comment{\lig{1059-1065}}
    \State \textbf{si} $i$ pair et $i > 3$ : $\vc R \gets \vc R - \op{ACT3}$(4 derniers résidus) \Comment{\lig{1084-1089}}
    \State $\Delta\vc T^{(i)} \gets \op{RESO}(\vc S, \vc R)$ \Comment{incrément total}
    \State $e \gets$ équation \eqref{eq:critere-th} \Comment{\lig{1187-1217}}
    \State \textbf{si} $e \le \src{PRECISION}$ (avec \src{TTOL} : toutes les variations $\le \src{TTOL}$) : \Sortir
    \State $\vc T_\theta \gets \vc T_0 + \theta\,\Delta\vc T^{(i)}$ ; matériau $\gets f(\vc T_\theta)$
    \State $\vc R \gets \vc F(\vc T_\theta) - \vc K(\vc T_\theta)\,\vc T_\theta - \vc C(\vc T_\theta)\,\Delta\vc T^{(i)}/\Delta t + \vc S\,\Delta\vc T^{(i)}$ \Comment{\lig{1220-1468}}
  \EndFor
  \State $\vc T_0 \gets \vc T_0 + \Delta\vc T^{(i)}$
\EndFor
\State $\src{ESTIMATION.TEMPERATURES} \gets \vc T_0$ \Comment{\lig{1524}}
\end{algorithmic}
\end{algobox}

\subsection{TRANSLIN : une résolution par sous-pas}
\label{sec:translin}

\begin{glossaire}[Glossaire — TRANSLIN (en plus de celui de TRANSNON)]
$\vc K$ & \src{MAT_COND} & \op{CONDU} sur tous les sous-modèles thermiques (\src{MOD_THE}), donc conduction \emph{et} convection, plus \src{CONDUCTIVITE_CONSTANTE}. Le matériau n'est pas réévalué en température.\\
$\vc C$ & \src{MAT_CAPA} & \op{CAPACITE}, plus \src{CAPACITE_CONSTANTE} ; absente si \src{C} ou \src{RHO} manque.\\
$\vc S$ & \src{MAT_RIGI} & $\theta\,\vc K + \vc C/\Delta t$ + blocages.\\
— & \src{B1}, \src{A1} & $(\theta - 1)\,\vc K + \vc C/\Delta t$ et $\vc C/\Delta t$.\\
$\vc F$ & \src{CH_FL}, \src{CH_TI}, \src{CHAL_TE} & Flux \src{Q} à $t_m$, températures imposées à $t_1$, convection $\op{CONV}(T_e(t_m))$.\\
$\lambda_2$, $t_m$ & \src{LAMBDA2}, \src{TEMM} & Instant d'évaluation des chargements, $t_m = \lambda_2\,t_1 + (1-\lambda_2)\,t_0$. \src{SOUS-RELAXATION} n'est jamais transmis par \pr{PAS\_DEFA}, donc $\lambda_2 = 1$.\\
— & \src{ICAPA} & Vrai si le matériau a une capacité ; sinon $\theta = \lambda_2 = 1$ (stationnaire).\\
\end{glossaire}

\pr{TRANSLIN} applique le $\theta$-schéma~\eqref{eq:theta} à un problème supposé linéaire.
Avec $\vc K$ et $\vc C$ constantes, l'équation~\eqref{eq:theta} se réécrit
\begin{equation}
\big(\theta\,\vc K + \vc C/\Delta t\big)\,\vc T_1 = \big((\theta - 1)\,\vc K + \vc C/\Delta t\big)\,\vc T_0 + \vc F(t_m)
\qquad(\text{\lig{203-215}}),
\end{equation}
qu'une seule résolution suffit à résoudre.
Il n'y a ni itération, ni critère de convergence.
Les matrices sont réassemblées à chaque sous-pas, mais avec le matériau tel qu'il est donné : il n'est pas instancié en fonction de la température.
La matrice de convection n'apparaît pas explicitement (sa ligne est commentée) parce que \op{CONDU} sur \src{MOD_THE}, qui contient tous les sous-modèles thermiques, l'inclut déjà.
\pr{TRANSLIN} n'appelle ni \pr{PAS\_RAYO}, ni le modèle de sources, ni le changement de phase, ni \pr{CHARTHER}, ni \pr{PARATHER} : ces données sont ignorées sans message.

Les lignes 87-88 lisent \src{WTAB.'CAPACONST'} et \src{WTAB.'CONDCONST'} alors que \pr{TRANSLIN} n'a défini que \src{ETAB}.
Ce n'est pas une erreur : par la portée dynamique (§ \ref{sec:gibiane}), \src{WTAB} est celui de \pr{PASAPAS}, qui désigne la même table.

\begin{algobox}{2L}{TRANSLIN : thermique linéaire}{lignes de TRANSLIN · 24/09/11}
\begin{algorithmic}[1]
\Require $\vc T_n$ (\src{CONTINUATION}), sous-pas \src{TEM_CALC}
\State $\theta \gets \src{RELAXATION_THETA}$ (1) ; $\lambda_2 \gets \src{SOUS-RELAXATION}$, jamais transmis, donc 1 \Comment{\lig{72-81}}
\State \textbf{si} pas de capacité (\src{C} ou \src{RHO} absent) : $\theta \gets 1$ ; $\lambda_2 \gets 1$ \Comment{\lig{87-95}, stationnaire}
\ForAll{sous-pas $[t_0, t_1]$} \Comment{\src{BOUC_SO}, \lig{143}}
  \State $t_m \gets \lambda_2\,t_1 + (1-\lambda_2)\,t_0$ \Comment{\lig{148}}
  \State $\vc K \gets \op{CONDU}(\src{MOD_THE}, \src{MAT_THE})$ [+ conductivité constante] \Comment{\lig{172-175} ; matériau sans \op{VARI}}
  \State $\vc C \gets \op{CAPA}(\src{MOD_THE}, \src{MAT_THE})$ [+ capacité constante] \Comment{\lig{192-199}}
  \State $\vc F \gets$ flux $Q(t_m)$ + $T$ imposées$(t_1)$ + $\op{CONV}(T_e(t_m))$ \Comment{\lig{160-168, 184-188, 221-226}}
  \State $(\theta\,\vc K + \vc C/\Delta t$ + blocages$)\,\vc T_1 = ((\theta - 1)\,\vc K + \vc C/\Delta t)\,\vc T_0 + \vc F$ \Comment{\lig{203-215}}
  \State $\vc T_1 \gets \op{RESO}$, une seule résolution \Comment{\lig{230}}
  \State $\vc T_0 \gets \vc T_1$ \Comment{\lig{248-252}}
\EndFor
\State $\src{ESTIMATION.TEMPERATURES} \gets \vc T_1$ \Comment{\lig{270}}
\end{algorithmic}
\end{algobox}

\subsection{DUPONT2 : schéma de Dupont à trois niveaux}
\label{sec:dupont}

\begin{glossaire}[Glossaire — DUPONT2 (en plus de ceux de TRANSNON et TRANSLIN)]
$\lambda$ & \src{LAMBDA} & \src{RELAXATION_DUPONT}, $\tfrac14$ par défaut ; rendu à \pr{PASAPAS} en fin d'appel.\\
$\lambda_2$ & \src{LAMBDA2} & \src{SOUS_RELAXATION}, jamais transmis, donc $\tfrac12$ (1 sans capacité) : les chargements sont pris au milieu du pas.\\
$\vc T_{n-1}$, $\vc T_n$, $\vc T_{n+1}$ & \src{U0}, \src{U1}, \src{U2} & Les trois niveaux de température.\\
$t_{n-1}$, $t_n$, $t_{n+1}$, $t_m$ & \src{TEM0}, \src{TEM1}, \src{TEM2}, \src{TEMM} & Les instants correspondants et l'instant d'évaluation $t_m = \lambda_2\,t_{n+1} + (1-\lambda_2)\,t_n$.\\
$\Delta t$, $\Delta t_{prec}$ & \src{T_PAS}, \src{T_VPAS} & Pas courant (\src{NOUV_PAS}) et pas précédent (\src{VIEUXPAS}).\\
$\vc T^{*}$ & \src{UC_T} & Température extrapolée en $t_m$, où l'on évalue le matériau et le rayonnement.\\
— & \src{SCA1}, \src{SCA2}, \src{SCA3} & Poids de la conduction : $\tfrac12 + \lambda$, $2\lambda - \tfrac12$, $-\lambda$.\\
— & \src{INIT} & Vaut 1 s'il faut démarrer par \pr{TRANSNON}, 0 si les deux niveaux précédents sont disponibles.\\
— & \src{TET0} & Niveau $\vc T_{n-1}$ transmis par \pr{PASAPAS} ; l'entier 0 par défaut.\\
\end{glossaire}

\paragraph{Le schéma.}
\pr{DUPONT2} évalue la conduction sur trois niveaux de température, avec des poids qui somment à 1, et la capacité sur les deux derniers :
\begin{equation}
\label{eq:dupont}
\frac{\vc C(\vc T^{*})}{\Delta t}\big(\vc T_{n+1} - \vc T_n\big)
+ \vc K(\vc T^{*})\Big[\big(\tfrac12 + \lambda\big)\vc T_{n+1} + \big(\tfrac12 - 2\lambda\big)\vc T_n + \lambda\,\vc T_{n-1}\Big]
= \vc F(t_m).
\end{equation}
Le source assemble exactement ce système (\lig{408-431}) :
\begin{equation}
\Big(\frac{\vc C}{\Delta t} + \big(\tfrac12+\lambda\big)\vc K\Big)\vc T_{n+1}
= \Big(\frac{\vc C}{\Delta t} + \big(2\lambda - \tfrac12\big)\vc K\Big)\vc T_n - \lambda\,\vc K\,\vc T_{n-1} + \vc F .
\end{equation}
Pour des pas égaux, la combinaison $\big(\tfrac12+\lambda\big)t_{n+1} + \big(\tfrac12-2\lambda\big)t_n + \lambda\,t_{n-1}$ vaut $t_n + \Delta t/2$, quel que soit $\lambda$ : la conduction est centrée au milieu du pas, comme la capacité et le chargement ($\lambda_2 = \tfrac12$).
Avec $\lambda = \tfrac14$, le poids de $\vc T_n$ s'annule et le crochet devient $\tfrac34\,\vc T_{n+1} + \tfrac14\,\vc T_{n-1}$.

\paragraph{La non-linéarité.}
Elle est traitée sans itération : le matériau est évalué une seule fois, à la température extrapolée linéairement en $t_m$ à partir des deux niveaux connus (\op{COLI}, \lig{238-242}),
\begin{equation}
\vc T^{*} = \frac{t_m - t_{n-1}}{t_n - t_{n-1}}\,\vc T_n + \frac{t_n - t_m}{t_n - t_{n-1}}\,\vc T_{n-1}
\;=\; \tfrac32\,\vc T_n - \tfrac12\,\vc T_{n-1}\quad\text{pour des pas égaux}.
\end{equation}
Le rayonnement est linéarisé une fois, à $\vc T^{*}$ (\annref{E}), et \pr{CHARTHER} est appelée à $t_m$.
Si le matériau thermique dépend de champs mécaniques, \pr{DUPONT2} s'arrête en erreur dès que $\lambda_2 \neq 0$, donc toujours (\lig{280-287}).

\paragraph{Le démarrage.}
Le schéma a besoin de deux niveaux.
Si \src{TET0} est un champ et que le pas n'a pas changé, $|\Delta t - \Delta t_{prec}|/\Delta t \le 10^{-5}$, on prend $\vc T_{n-1} = \src{TET0}$ et $\vc T_n = \vc T$ de \src{CONTINUATION} (\lig{188-197}).
Sinon, on fait le premier sous-pas avec \pr{TRANSNON} ; s'il n'y a qu'un sous-pas, on rend ce résultat, sinon on enchaîne les sous-pas suivants avec Dupont (\lig{199-222}).
Comme \src{VIEUXPAS} repart à 0 à chaque appel de \pr{PASAPAS} (\pr{PAS\_DEFA} \lig{572}), le premier pas de chaque appel passe toujours par \pr{TRANSNON}.

\begin{figure}[ht]
\centering
\begin{tikzpicture}[x=1.5cm,y=1cm,>={Stealth[length=2mm]}]
  \draw[->] (-0.4,0) -- (3.6,0) node[right] {$t$};
  \draw[->] (-0.4,0) -- (-0.4,3.3) node[above] {$T$};
  \foreach \x/\l in {0/{t_{n-1}},1/{t_n},1.5/{t_m},2/{t_{n+1}}} \draw (\x,0.07) -- (\x,-0.07) node[below] {$\l$};
  \draw[dashed,gray] (0,0.9) -- (1.5,2.25);
  \fill (0,0.9) circle (1.6pt) node[above left] {$\vc T_{n-1}$};
  \fill (1,1.8) circle (1.6pt) node[above left] {$\vc T_n$};
  \draw[thick,ther] (1.5,2.25) circle (2.2pt) node[above] {$\vc T^{*}$};
  \fill[meca] (2,2.5) circle (1.6pt) node[right] {$\vc T_{n+1}$ (inconnue)};
  % cas dégénéré
  \begin{scope}[xshift=7.5cm]
    \draw[->] (-0.4,0) -- (3.0,0) node[right] {$t$};
    \draw[->] (-0.4,0) -- (-0.4,3.3) node[above] {$T$};
    \foreach \x/\l in {1/{t_n},1.5/{t_m},2/{t_{n+1}}} \draw (\x,0.07) -- (\x,-0.07) node[below] {$\l$};
    \draw[dashed,gray] (0,1.8) -- (1.5,1.8);
    \fill (1,1.8) circle (1.6pt) node[above left] {$\vc T_{n-1} = \vc T_n$};
    \draw[thick,ther] (1.5,1.8) circle (2.2pt) node[below=3pt] {$\vc T^{*} = \vc T_n$};
    \fill[meca] (2,2.5) circle (1.6pt) node[right] {$\vc T_{n+1}$};
  \end{scope}
\end{tikzpicture}
\caption{À gauche, l'extrapolation de \pr{DUPONT2} entre deux sous-pas d'un même pas. À droite, ce qu'il reçoit d'un pas \pr{PASAPAS} au suivant : \src{TET0} et \src{CONTINUATION} portent la même température, l'extrapolation ne voit plus rien venir.}
\label{fig:dupont}
\end{figure}

\paragraph{La dégénérescence d'un pas à l'autre.}
En fin de pas, \pr{PASAPAS} range \src{ESTIMATION.TEMPERATURES} dans \src{TET0} (\lig{511}).
\pr{DUPONT2} vient d'y écrire $\vc T_{n+1}$ (\lig{485}), et \pr{PAS\_RESU} ne le recopie dans \src{CONTINUATION} qu'ensuite (\lig{598}).
Au pas suivant, \pr{DUPONT2} lit donc $\vc T_{n-1} = \vc T_n$ (figure~\ref{fig:dupont}, à droite).
L'extrapolation donne $\vc T^{*} = \vc T_n$, et le schéma~\eqref{eq:dupont} devient
\begin{equation}
\frac{\vc C(\vc T_n)}{\Delta t}\big(\vc T_{n+1} - \vc T_n\big) + \vc K(\vc T_n)\Big[\big(\tfrac12+\lambda\big)\vc T_{n+1} + \big(\tfrac12-\lambda\big)\vc T_n\Big] = \vc F(t_m),
\end{equation}
c'est-à-dire un $\theta$-schéma à deux niveaux avec $\theta = \tfrac12 + \lambda = \tfrac34$, et un matériau figé au début du pas.
Le vrai schéma à trois niveaux ne s'applique qu'entre les sous-pas d'un même pas (\src{PAS_MAX}).
Après une reprise, \pr{PAS\_INIT} range dans \src{TET0} la température d'indice $N_{sor} - 2$, deux instants avant le dernier (\pr{PAS\_INIT} \lig{84}) ; mais \src{VIEUXPAS} étant nul, ce \src{TET0} n'est pas utilisé.

\begin{algobox}{2D}{DUPONT2 : schéma de Dupont à trois niveaux}{lignes de DUPONT2 · 24/09/11}
\begin{algorithmic}[1]
\Require $\vc T_n$ (\src{CONTINUATION}), \src{TET0}, $\Delta t = \src{NOUV_PAS}$, $\Delta t_{prec} = \src{VIEUXPAS}$
\State $\lambda \gets \src{RELAXATION_DUPONT}$ (0,25) ; $\lambda_2 \gets \src{SOUS_RELAXATION}$, jamais transmis, donc $\tfrac12$ \Comment{\lig{92-104}}
\State \textbf{si} pas de capacité : $\lambda_2 \gets 1$ \Comment{\lig{156-164}}
\If{\src{TET0} est un champ et $|\Delta t - \Delta t_{prec}|/\Delta t \le 10^{-5}$} \Comment{\lig{188-197}}
  \State $\vc T_{n-1} \gets \src{TET0}$ ; $\vc T_n \gets \vc T_n$ \Comment{niveaux $t_{n-1}$ et $t_n$}
\Else \Comment{premier pas, ou $\Delta t$ modifié}
  \State $\vc T_{n+1} \gets$ \algoref{2} sur le premier sous-pas \Comment{\pr{TRANSNON}, \lig{199-207}}
  \State \textbf{si} un seul sous-pas : \Retourner{} $\vc T_{n+1}$ \Comment{\lig{209-213}}
  \State $\vc T_{n-1} \gets \vc T_n$ ; $\vc T_n \gets \vc T_{n+1}$
\EndIf
\ForAll{sous-pas restant $[t_n, t_{n+1}]$} \Comment{\src{BOUC_SO}, \lig{233}}
  \State $t_m \gets \lambda_2\,t_{n+1} + (1-\lambda_2)\,t_n$
  \State $\vc T^{*} \gets$ extrapolation linéaire en $t_m$ depuis $(t_{n-1}, \vc T_{n-1})$ et $(t_n, \vc T_n)$ \Comment{\op{COLI}, \lig{238-242}}
  \State matériau $\gets f(\vc T^{*}, t_m)$ ; dépendance mécanique : erreur si $\lambda_2 \neq 0$ \Comment{\lig{252-325}}
  \State $\vc K \gets$ conduction + rayonnement à $\vc T^{*}$ (\annref{E}) + \pr{CHARTHER} ; $\vc C \gets$ capacité \Comment{\lig{327-406}}
  \State $\vc F \gets$ flux $Q(t_m)$ + $T$ imposées$(t_{n+1})$ + convection $T_e(t_m)$ + seconds membres \Comment{\lig{310-406}}
  \State $(\vc C/\Delta t + (\tfrac12+\lambda)\vc K)\,\vc T_{n+1} = (\vc C/\Delta t + (2\lambda-\tfrac12)\vc K)\,\vc T_n - \lambda\,\vc K\,\vc T_{n-1} + \vc F$ \Comment{\lig{408-431}}
  \State $\vc T_{n-1} \gets \vc T_n$ ; $\vc T_n \gets \vc T_{n+1}$ \Comment{\lig{455-461}}
\EndFor
\State $\src{ESTIMATION.TEMPERATURES} \gets \vc T_{n+1}$ ; \src{RELAXATION_DUPONT} rendu à \pr{PASAPAS} \Comment{\lig{476-486}}
\end{algorithmic}
\end{algobox}

\subsection{Comparaison des trois procédures}

Le tableau~\ref{tab:thermique} résume ce qui distingue les trois procédures.

\begin{table}[H]
\centering\small
\caption{Les trois procédures thermiques.}
\label{tab:thermique}
\begin{tabularx}{\linewidth}{@{}>{\raggedright\arraybackslash}p{0.2\linewidth}>{\raggedright\arraybackslash}X>{\raggedright\arraybackslash}X>{\raggedright\arraybackslash}X@{}}
\toprule
 & \pr{TRANSNON} & \pr{TRANSLIN} & \pr{DUPONT2}\\
\midrule
Schéma en temps & $\theta$-schéma, $\theta = 1$ & $\theta$-schéma, $\theta = 1$ & Dupont, $\lambda = \tfrac14$ ($\theta = \tfrac34$ d'un pas à l'autre)\\
Matériau & réévalué à chaque itération, à $\vc T_\theta$ & constant & évalué une fois, à $\vc T^{*}$ extrapolée\\
Résolutions par sous-pas & jusqu'à \src{MAXITERATION} & 1 & 1 ; \pr{TRANSNON} au 1\textsuperscript{er} pas de chaque appel et si $\Delta t$ change\\
Contrôle de convergence & variation de $\Delta\vc T$ $\le$ \src{PRECISION} & aucun & aucun\\
Chargements évalués à & $t_\theta$ ($= t_1$) ; $T$ imposées à $t_1$ & $t_m$ ($= t_1$) ; $T$ imposées à $t_1$ & mi-pas ; $T$ imposées à $t_{n+1}$\\
Rayonnement & oui : matrice à $\vc T_0$, résidu à $\vc T_\theta$ & non & oui, une fois à $\vc T^{*}$\\
Sources (\src{SOURCE}) & oui & non & non\\
Changement de phase, diffusion & oui & non & non\\
\pr{CHARTHER} / \pr{PARATHER} & oui / oui & non / non & oui / oui\\
\bottomrule
\end{tabularx}
\end{table}
% ============================================================================
\section{Un pas de mécanique : UNPAS}
\label{sec:pas}
% ============================================================================

\pr{UNPAS} fait avancer la mécanique de $t$ à $T_I$.
Il prépare le pas (matériau de fin de pas, chargement, pseudo-chargements, premier second membre), puis lance une boucle de sous-pas, \src{BONOCONV}, qui ne compte qu'un tour si les itérations d'équilibre convergent.
En cas d'échec, il repart de l'itéré non convergé : c'est la « convergence forcée ».

\begin{glossaire}
$\vc u_n$, $\vg\lambda_n$ & \src{ZDEP0} & Déplacements de début de (sous-)pas, avec les multiplicateurs (composantes \src{LX}). Après un échec, \src{ZDEP0} prend l'itéré non convergé \src{ZDETOT}.\\
$\sig_n$ & \src{ZSIG0}, \src{ZSIG1} & Contraintes de début de pas, rangées sur FOR0 ; \src{ZSIG1} est leur transport sur GEOM1.\\
$\vg\eps_n$, $\vc V_n$, $\vg\eps^{in}_n$ & \src{ZDEF0}, \src{ZVAR0}, \src{ZDEI0} & Déformations, variables internes, déformations inélastiques de début de pas.\\
$\vc D_n$ & \src{ZMAT1} (\src{WTAB.'MAT1'}) & Matériau de début de pas.\\
$\vc D_{n+1}$ & \src{ZMAT20}, \src{ZMAT21}, \src{ZMAT22} & Matériau à $T_I$ (\pr{PAS\_MATE} de \pr{PAS\_ETAT} à $T_I$), rangé sur FOR0, transporté sur GEOM1, puis sur GEOM2.\\
$T_n$, $T_{n+1}$ & \src{ZTEMP1}, \src{ZTEMP2} & Températures de début et de fin de pas vues par la mécanique.\\
$\vc F^{ext}_{n+1}$ & \src{ZFEXT2}, \src{ZFEXT} & Chargement à $T_I$ : \src{MECA}, \src{FORC}, pression convertie en forces en HPP, \src{DIMP}, \src{DINC}.\\
— & \src{ZFCONSTA} & $[\vc F^{ext}_{n+1} ; \vc u^{imp}_{n+1}]$ : forces et valeurs imposées des liaisons.\\
$\vc F^{suiv}$ & \src{ZFSUIV} & Forces suiveuses : pression en GD, \pr{CHARMECA}.\\
$\Deps^{*}$ & \src{DEPSI} & Pseudo-déformation du pas : thermique, \src{DEFI}, matériau variable.\\
$\vg\eps^{*}_{n+1}$ & \src{EPSIF} & Déformations imposées totales en fin de pas (thermique et \src{DEFI}).\\
$\vc F^{*}$ & \src{DFDEF} & Pseudo-force de $\Deps^{*}$ : $\vc B^{\mathsf T}\vc D_{n+1}\Deps^{*}$ ; \src{DSIGM} est la contrainte $\vc D_{n+1}\Deps^{*}$.\\
— & \src{FSIGF} & Force des déformations imposées totales, $\vc B^{\mathsf T}\vc D_{n+1}\vg\eps^{*}_{n+1}$ ; ne sert qu'à $F_{ref}$.\\
$\vc F^{int}_n$ & \src{ZF1} & $\op{BSIGMA}(\sig_n)$ sur GEOM1 (plus $\vc K^{cst}\vc u_n$ avec \src{RIGIDITE_CONSTANTE}).\\
$\vc A^{\mathsf T}\vg\lambda_n$, $\vc A\vc u_n$ & \src{FREAP}, \src{FLXINI} & Les deux parties de $\vc A\cdot[\vc u_n ; \vg\lambda_n]$ : forces de liaison et valeurs actuelles des liaisons.\\
$\vc F^{ext}_n$ & \src{FEXT0} & $\vc F^{int}_n + \vc A^{\mathsf T}\vg\lambda_n$ : chargement équilibré en début de pas.\\
$\vc R^{0}$ & \src{RESIDU} & Premier second membre.\\
— & \src{ZDFORC}, \src{ZDFINI} & Incrément de chargement $[\vc F^{ext}_{n+1} - \vc F^{ext}_n ; \Delta\vc u^{imp}]$, et le même plus $\vc F^{*}$.\\
$\vc K$ & \src{ZRAID} & Opérateur des itérations (§ \ref{sec:operateur}).\\
$(\delta\vc u^{0}, \vg\lambda^{0})$ & \src{ZDEP1} & Première solution, \op{RESO} de $\vc K$ et $\vc R^{0}$.\\
$F_{ref}$, $M_{ref}$ & \src{XDENO}, \src{XDENOM} & Forces et moments de référence du critère (§ \ref{sec:fref}).\\
$\vc F^{NL}$ & \src{ZFNONL} & Part non linéaire des forces internes du pas précédent : prédicteur.\\
$a$ & \src{AMPL} & Coefficient du prédicteur ; \src{AMPLT} $= \Delta t/\Delta t_{prec}$.\\
— & \src{INIT1} & Vrai si le prédicteur a été appliqué ; le test de convergence est alors sauté à l'itération 1.\\
$s$ & \src{knoconv} & Numéro du sous-pas ; \src{WTAB.'ISOUSPAS'} compte les sous-pas ratés du pas, \src{ISOUSPPP} ceux du pas précédent.\\
$\alpha$ & \src{COEPI} & Fraction de l'intervalle restant visée par le sous-pas : 1, ou \src{RELAXATION_NONCONV} après une non-convergence détectée.\\
$t_{déb}$, $\Delta t$ & \src{TDEBUT}, \src{PASDT} & Début du sous-pas et durée restante jusqu'à $T_I$ ; \src{ZDT} $= \alpha\,\Delta t$.\\
— & \src{ZICONV} & Vrai si le sous-pas a convergé.\\
— & \src{IVISCO}, \src{IVIDOM}, \src{IVIEXT}, \src{IDYN}, \src{ISOL} & Lois visqueuses, viscodommage, loi visqueuse externe, dynamique, consolidation.\\
— & \src{HPP_EPS} & Prédicteur HPP en cours (option \src{PREDICTEUR 'HPP'}).\\
\end{glossaire}

\subsection{Le problème du pas}

On cherche $\Delta\vc u = \vc u_{n+1} - \vc u_n$ et $\vg\lambda_{n+1}$ vérifiant~\eqref{eq:equilibre} à $T_I$.
\pr{UNPAS} ne linéarise pas autour de l'itéré courant à chaque itération : il résout une suite de systèmes linéaires avec un opérateur $\vc K$ en général fixe (§ \ref{sec:operateur}), dont le second membre est le déséquilibre courant.
Cette section décrit ce qui précède les itérations : l'état de départ, les chargements, et le premier second membre.

\subsection{Matériau et chargement de fin de pas}

Le matériau de fin de pas $\vc D_{n+1}$ est évalué à $T_I$ : \pr{PAS\_ETAT} construit l'état (température \src{ESTIMATION}, temps, champs) et \pr{PAS\_MATE} évalue les propriétés (\lig{283-297}).
Le chargement $\vc F^{ext}_{n+1}$ est lu à $T_I$ par \op{TIRE} (\lig{301-347}).
En HPP, une pression est convertie une fois en forces nodales, par $\op{BSIGMA}$ sur le modèle de pression ; en GD, elle devient une force suiveuse, recalculée sur la configuration courante à chaque itération.

\subsection{Pseudo-déformations et pseudo-forces}

Trois phénomènes produisent une déformation que la structure doit accommoder sans qu'elle provienne d'un déplacement : la dilatation thermique, les déformations imposées (\src{DEFI}), et la variation du matériau pendant le pas.
\pr{UNPAS} les regroupe dans une pseudo-déformation (\lig{987-1046}) :
\begin{equation}
\label{eq:depsstar}
\Deps^{*} = \underbrace{\big[\vg\eps^{th}_{n+1} - \vg\eps^{th}_n\big]}_{\text{thermique}}
+ \underbrace{\big[\vg\eps^{DEFI}_{n+1} - \vg\eps^{DEFI}_n\big]}_{\text{imposée}}
+ \underbrace{\big[\vc D_{n+1}^{-1} - \vc D_n^{-1}\big]\sig_n}_{\text{matériau variable}} .
\end{equation}
Le dernier terme mérite une explication.
Le source le calcule par deux appels à \op{ELAS} sur $\sig_n$ (déformation élastique avec le matériau de fin, puis de début).
Pour un comportement élastique linéaire, l'incrément de contrainte de \algoref{5} vaut $\Delta\sig = \vc D_{n+1}\,(\Deps - \Deps^{*})$.
À déformation constante ($\Deps = 0$) et sans dilatation, il donne
\begin{equation}
\sig_{n+1} = \sig_n - \vc D_{n+1}\big(\vc D_{n+1}^{-1} - \vc D_n^{-1}\big)\sig_n = \vc D_{n+1}\,\vc D_n^{-1}\,\sig_n = \vc D_{n+1}\,\vg\eps^{el}_n :
\end{equation}
la déformation élastique est conservée, et la contrainte suit la variation du module.
La loi élastique reste donc totale, $\sig = \vc D(T)\,(\vg\eps - \vg\eps^{*})$, au lieu de devenir hypo-élastique.
Pour un comportement non linéaire, $\Deps^{*}$ entier est retiré de l'incrément de déformation transmis à \op{COMP} (\lig{2443-2447}) ; ce que fait ensuite \op{COMP} de la différence entre matériau de début et de fin dépend de la loi, et n'a pas été lu.

La pseudo-déformation devient une pseudo-force, assemblée sur la configuration de début de pas (\lig{1060-1073}) :
\begin{equation}
\vc F^{*} = \int_\Omega \vc B^{\mathsf T}\,\vc D_{n+1}\,\Deps^{*}\,\mathrm d\Omega .
\end{equation}

\subsection{Le premier second membre}

Le chargement de début de pas est reconstruit à partir de l'état équilibré, $\vc F^{ext}_n = \vc F^{int}_n + \vc A^{\mathsf T}\vg\lambda_n$, ce qui évite de relire le chargement à $t$ (\lig{1158-1168}).
Le premier second membre est (\lig{1129-1204})
\begin{equation}
\label{eq:R0}
\vc R^{0} = \begin{Bmatrix}\vc F^{ext}_{n+1} + \vc F^{suiv} + \vc F^{*} + \vc F^{rea} - \vc F^{int}(\sig_n)\\[2pt] \vc u^{imp}_{n+1} - \vc A\,\vc u_n\end{Bmatrix},
\end{equation}
où $\vc F^{rea}$ n'existe qu'en dynamique (§ \ref{sec:dynamique}).
La première ligne ne contient pas les réactions : comme les multiplicateurs sont résolus « en total » (§ \ref{sec:systeme}), la solution $\vg\lambda^{0}$ est directement la réaction cherchée.
La seconde ligne est l'incrément des valeurs imposées.

\subsection{Le prédicteur}
\label{sec:predicteur}

À la fin de chaque pas, \pr{UNPAS} garde la part non linéaire des forces internes (\lig{3607-3628}) :
\begin{equation}
\vc F^{NL}_n = \vc K\,\Delta\vc u_n - \big(\vc F^{int}_{n+1} - \vc F^{int}_n\big) - \alpha\,\vc F^{*}_n .
\end{equation}
C'est la différence entre ce qu'aurait donné l'opérateur des itérations et ce que la loi a réellement produit : elle est nulle pour un pas thermo-élastique.
Après un sous-pas raté, elle s'ajoute à celle des sous-pas précédents.

Au pas suivant, on ajoute $a\,\vc F^{NL}$ au premier second membre (\lig{1393-1443}), en pariant que la non-linéarité se répète.
Le coefficient $a$ est le rapport des incréments de chargement des deux pas, projetés sur la somme des deux premières solutions $\vc z = \delta\vc u^{0}_{n} + \delta\vc u^{0}_{n+1}$ :
\begin{equation}
a = \min\!\left(\frac{\Delta\vc F_{n+1}\cdot\vc z}{\Delta\vc F_n\cdot\vc z},\ \frac{\Delta t}{\Delta t_{prec}}\right).
\end{equation}
Si l'incrément de chargement est négligeable (fluage, thermique), ou petit devant $\vc F^{NL}$ (projection de $\vc F^{NL}$ plus de 200 fois plus grande), on prend $a = \Delta t/\Delta t_{prec}$, ou $a = 0$ si, dans ce cas, le rapport des chargements est inférieur à $-0{,}05$ (décharge).
Le prédicteur ne s'applique que si $0 < a < 20$, si le pas précédent a convergé sans sous-pas, s'il était non linéaire ($\norminf{\vc F^{NL}} > \src{PRECISION}\cdot F_{ref}$), et si le modèle n'a pas changé.
Quand il s'applique, le test de convergence est sauté à l'itération~1.
En pilotage, on l'applique avec $a = 1$ quand on rejoue une cible.

\subsection{Sous-pas et convergence forcée}

La boucle \src{BONOCONV} enchaîne les sous-pas.
Le premier vise $T_I$ avec $\alpha = 1$.
Si les itérations (\algoref{4}) échouent et que la convergence forcée est désactivée, \pr{UNPAS} rend un échec, et \pr{PASAPAS} s'arrête sur l'erreur~997.
Sinon, l'état non convergé devient l'état de départ d'un nouveau sous-pas (\lig{3660-3721}) :
\begin{equation}
t_{déb} \gets t_{déb} + \alpha\,(T_I - t_{déb}),\qquad \Delta t \gets T_I - t_{déb},\qquad
(\vc u_n, \sig_n, \vc V_n, \vg\eps_n, T_n, \mathrm{GEOM1}) \gets \text{itéré non convergé}.
\end{equation}
Avec $\alpha = 1$ (défaut de \src{RELAXATION_NONCONV}), il ne reste alors plus de temps à intégrer : le sous-pas suivant ne fait que rééquilibrer l'état à $T_I$.
Pour les lois visqueuses, en dynamique et en consolidation, un intervalle nul est une erreur (1155, \lig{3672-3674}).
Avec $\alpha < 1$, la cible du sous-pas est ramenée à $t_{déb} + \alpha\,\Delta t$ dès que la non-convergence est détectée (§ \ref{sec:criteres}) ; une fois cette cible atteinte, \pr{UNPAS} déclare « non convergence due à la relaxation » et repart de là : c'est un vrai sous-découpage du pas.
Au-delà de \src{MAXSOUSPAS} (100) sous-pas ratés, \pr{PASAPAS} s'arrête sur l'erreur~996 (\lig{3718-3721}).
Un état non convergé n'est donc jamais accepté comme résultat.

\begin{algobox}{3}{UNPAS : préparation et sous-pas}{lignes d'UNPAS}
\begin{algorithmic}[1]
\Require état de début ($\vc u_n$, $\vg\lambda_n$, $\sig_n$, $\vc V_n$, $\vg\eps_n$), $T_n$, $T_{n+1}$, $t$, $T_I$
\State $\vc D_{n+1} \gets$ matériau à $T_I$ (\pr{PAS\_ETAT}, \pr{PAS\_MATE}) \Comment{\lig{283-297}}
\State $\vc F^{ext} \gets$ chargement à $T_I$ ; pression convertie une fois en HPP \Comment{\lig{301-347}}
\State $\Deps^{*} \gets [\op{EPTH}(T_{n+1}) - \op{EPTH}(T_n)] + [\mathrm{DEFI}_{n+1} - \mathrm{DEFI}_n] + [\vc D_{n+1}^{-1} - \vc D_n^{-1}]\,\sig_n$ \Comment{\lig{987-1046}}
\State $\vc F^{*} \gets \vc B^{\mathsf T}\vc D_{n+1}\Deps^{*}$ \Comment{pseudo-force, \lig{1060-1073}}
\State \textbf{si} dynamique : $\vc F^{rea} \gets \vc M\vc a_n + \vc C\vc v_n + (4/\Delta t)\,\vc M\vc v_n$ \Comment{\lig{407-471}}
\State $t_{déb} \gets t$
\For{$s = 1 \dots \src{MAXSOUSPAS}$} \Comment{\src{BONOCONV}, \lig{649}}
  \State $\vc K \gets$ opérateur initial (tableau~\ref{tab:operateur}) \Comment{\lig{709-966}}
  \State $\vc R^{0} \gets [\vc F^{ext} + \vc F^{suiv} + \vc F^{*} + \vc F^{rea} - \vc F^{int}(\sig_n) ;\ \vc u^{imp} - \vc A\vc u_n]$ \Comment{\lig{1129-1204}}
  \State $(\delta\vc u^{0}, \vg\lambda^{0}) \gets \op{RESO}(\vc K, \vc R^{0})$ \Comment{\lig{1252-1322}}
  \State \textbf{si} $s = 1$ : $F_{ref}, M_{ref} \gets$ références du critère (§ \ref{sec:fref}) \Comment{\lig{1334-1355}}
  \State \textbf{si} le prédicteur s'applique : $\vc R^{0} \gets \vc R^{0} + a\,\vc F^{NL}$ \Comment{\lig{1393-1443}}
  \State (état, convergé) $\gets$ \algoref{4}
  \State $\vc F^{NL} \gets \vc K\,\Delta\vc u - \Delta\vc F^{int} - \alpha\,\vc F^{*}$, cumulée d'un sous-pas à l'autre après un échec \Comment{\lig{3607-3628}}
  \State \textbf{si} convergé : \Sortir
  \State \textbf{si} pas de \src{CONVERGENCE_FORCEE} : \Retourner{} échec \Comment{erreur 997 dans \pr{PASAPAS}}
  \State $t_{déb} \gets t_{déb} + \alpha\,(T_I - t_{déb})$ ; $\Delta t \gets T_I - t_{déb}$ \Comment{$\alpha$ = \src{COEPI}, \lig{3670-3671}}
  \State \textbf{si} $\Delta t \approx 0$ et (loi visqueuse ou dynamique ou consolidation) : erreur 1155 \Comment{\lig{3672-3674}}
  \State $\vc u_n, \sig_n, \vc V_n, \vg\eps_n, T_n, \mathrm{GEOM1} \gets$ itéré non convergé \Comment{\lig{3676-3704}}
  \State matériau de début $\gets$ matériau à $T_I$ \Comment{\lig{700-707}}
  \State \textbf{si} $s = \src{MAXSOUSPAS}$ : erreur 996 \Comment{\lig{3718-3721}}
\EndFor
\State $\src{ESTIMATION} \gets \vc u$, réactions, $\sig$, $\vg\eps$, $\vc V$, $\vg\eps^{in}$ ; $\vc v$ et $\vc a$ en dynamique \Comment{\lig{3733-3828}}
\end{algorithmic}
\end{algobox}

\subsection{Dynamique et consolidation}
\label{sec:dynamique}

\begin{glossaire}[Glossaire — dynamique]
$\vc M$, $\vc C$ & \src{WTAB.'MASSE'}, \src{WTAB.'AMORTISSEMENT'} & Matrices de masse et d'amortissement.\\
$\vc v_n$, $\vc a_n$ & \src{conti.'VITESSES'}, \src{conti.'ACCELERATIONS'} & Vitesses et accélérations de début de pas.\\
$\vc F^{rea}$ & \src{FREA1} & $\vc M\vc a_n + \vc C\vc v_n$, puis $+ (4/\Delta t)\,\vc M\vc v_n$ : partie du second membre qui ne dépend que du pas précédent.\\
— & \src{FOPL} & Déséquilibre statique de fin de pas, $\vc F^{ext} + \vc F^{suiv} - \vc F^{int}$, soit $\vc M\vc a_{n+1} + \vc C\vc v_{n+1}$ ; devient \src{FREA1} au pas suivant.\\
$1/\Delta t$ & \src{UNSURH} & Inverse du pas.\\
— & \src{RECAOP} & Faut-il reconstruire l'opérateur dynamique ? Vrai au premier appel et jamais remis à faux.\\
$\vc v_{n+1}$, $\vc a_{n+1}$ & \src{VITI}, \src{ACCEI} & Vitesses et accélérations de fin de pas.\\
$\theta_c$ & \src{TETA} & Paramètre du schéma de consolidation, 1.\\
\end{glossaire}

\paragraph{Le schéma de Newmark.}
La dynamique utilise le schéma de l'accélération moyenne ($\gamma = \tfrac12$, $\beta = \tfrac14$) :
\begin{equation}
\Delta\vc u = \Delta t\,\vc v_n + \frac{\Delta t^2}{4}\big(\vc a_n + \vc a_{n+1}\big),\qquad
\vc v_{n+1} = \vc v_n + \frac{\Delta t}{2}\big(\vc a_n + \vc a_{n+1}\big).
\end{equation}
On en tire $\vc a_{n+1} = \tfrac{4}{\Delta t^2}\Delta\vc u - \tfrac4{\Delta t}\vc v_n - \vc a_n$ et $\vc v_{n+1} = \tfrac2{\Delta t}\Delta\vc u - \vc v_n$, qu'on reporte dans $\vc M\vc a_{n+1} + \vc C\vc v_{n+1} + \vc F^{int} = \vc F^{ext}$.
L'itération résout donc
\begin{equation}
\Big(\vc K + \tfrac{4}{\Delta t^{2}}\vc M + \tfrac{2}{\Delta t}\vc C\Big)\,\delta\vc u
= \vc F^{ext}_{n+1} + \underbrace{\vc M\vc a_n + \vc C\vc v_n + \tfrac{4}{\Delta t}\vc M\vc v_n}_{\vc F^{rea}}
- \vc F^{int} - \tfrac{4}{\Delta t^{2}}\vc M\,\Delta\vc u - \tfrac{2}{\Delta t}\vc C\,\Delta\vc u ,
\end{equation}
les deux derniers termes étant ajoutés au résidu à chaque itération (\lig{3110-3151}).
En fin de pas, $\vc v_{n+1} = 2\Delta\vc u/\Delta t - \vc v_n$ et $\vc a_{n+1} = 2(\vc v_{n+1} - \vc v_n)/\Delta t - \vc a_n$ (\lig{3534-3579}).
Au premier pas, $\vc M\vc a_n + \vc C\vc v_n$ est calculé comme $\vc F^{ext}_n - \vc F^{int}_n$ ; ensuite, c'est \src{FOPL} du pas précédent.
L'opérateur dynamique est reconstruit à chaque pas, car \src{RECAOP} vaut vrai au premier appel (\pr{PAS\_DEFA} \lig{674}) et n'est jamais remis à faux ; les tests sur $\Delta t$ et sur le matériau ne jouent qu'en reprise.

\paragraph{La consolidation.}
Elle ajoute à l'opérateur le terme de perméabilité $-\theta_c\,\Delta t\,\vc H$ ($\theta_c = 1$), et à $\vc F^{int}_n$ le flux $\Delta t\;\op{GNFL}\big((1-\theta_c)\nabla p_n + \theta_c\nabla p\big)$ (\lig{1146-1155}).
% ============================================================================
\section{Les itérations d'équilibre et l'opérateur}
\label{sec:iterations}
% ============================================================================

\begin{glossaire}
$i$ & \src{IT} & Numéro de l'itération dans le sous-pas (boucle \src{ETIQ}).\\
$\vc K$ & \src{ZRAID} & Opérateur des itérations, blocages compris.\\
$(\delta\vc u, \vg\lambda^{(i)})$ & \src{ZDEP1} & Solution de l'itération : correction des déplacements, et multiplicateurs en total.\\
$(\Delta\vc u^{(i)}, \vg\lambda^{(i)})$ & \src{ZDEPT} & Incrément cumulé depuis le début du sous-pas ; la partie \src{LX} est le multiplicateur total.\\
$\vc u^{(i)}$ & \src{ZDETOT} & Déplacement total, $\vc u_n + \Delta\vc u^{(i)}$.\\
— & \src{GEOM2} & Configuration courante en GD, $\mathrm{FOR0} + \vc u^{(i)}$.\\
$\sig^{(i)}$, $\vc V^{(i)}$ & \src{ZSIGF}, \src{ZVARF} & Contraintes et variables internes rendues par \algoref{5}.\\
$\vc F^{int}$ & \src{FEQU2} & $\op{BSIGMA}(\sig^{(i)})$ sur GEOM2, plus les termes dynamiques.\\
$\vc A^{\mathsf T}\vg\lambda^{(i)}$, $\vc A\vc u^{(i)}$ & \src{FCORF}, \src{FCORU} & Les deux parties de $\vc A\cdot[\vc u^{(i)} ; \vg\lambda^{(i)}]$ (\src{ZCLIM * ZDETOT}).\\
$\vc r^{(i)}$ & \src{RESIDC} & Déséquilibre vrai, $\vc F^{ext} + \vc F^{suiv} - \vc F^{int} - \vc A^{\mathsf T}\vg\lambda^{(i)}$ : sert au critère.\\
$\vc R^{(i)}$ & \src{RESIDU} & Second membre de l'itération suivante, $[\vc F^{ext} + \vc F^{suiv} - \vc F^{int} ; \vc u^{imp} - \vc A\vc u^{(i)}]$.\\
— & \src{ITACC} & Compteur de l'accélération : \op{ACT3} s'applique quand il est $\le 0$ et $i > 3$ ; remis à 2 après chaque accélération, à 4 pour la suspendre.\\
— & \src{ZNACCE} & Accélération autorisée ; faux avec \src{K_TANGENT} en FEFP ou SSTE.\\
$\vc F_0 \dots \vc F_3$ & \src{ACFP0} … \src{ACFP3} & Les quatre derniers résidus, pondérés par $\vc J$.\\
$\vc D_0 \dots \vc D_2$ & \src{ACFEP0} … \src{ACFEP2} & Les trois derniers résidus non pondérés, diminués de la correction précédente.\\
$\vc c$, $\vc c^{-}$ & \src{CORREC}, \src{CORRECA} & Correction d'\op{ACT3} de l'itération, et celle de l'itération précédente.\\
$\vc J$ & \src{ZJAC} & Poids nodaux $1/\sqrt{V_{nœud}}$, $V_{nœud}$ étant la somme des volumes des éléments qui touchent le nœud (\lig{583-588}).\\
$\delta\vc u^{-}$ & \src{zdep1d} & Dernière correction de déplacement (sans \src{LX}) ; sert aux garde-fous d'\op{ACT3} et à l'augmentation.\\
— & \src{URG}, \src{ITURG}, \src{IT_RECA}, \src{RECA_N} & Reconstruction de $\vc K$ en GD : demandée d'urgence, itérations depuis la dernière, itération de la dernière, nombre de reconstructions du sous-pas.\\
$p_K$ & \src{ITRCLC} & Période de reconstruction en GD, $\max(20, \src{DELTAITER})$ (\lig{124-125}).\\
— & \src{IKTAN}, \src{NITER_KTANGENT}, \src{IPERT} & Option \src{K_TANGENT}, période de recalcul de la tangente, tangente par perturbation.\\
$\vc K_t$ & \src{ZRIKTA}, \src{LASTKTAN} & Tangente de l'itération ; dernière tangente, gardée d'un pas à l'autre.\\
$\vc K_\sigma$ & \src{ZKSIG} (\src{IKSIA}) & Rigidité géométrique (\op{KSIGMA}), si \src{K_SIGMA}.\\
— & \src{ZPREK}, \src{DTTAN} & Précision et pas de temps passés à \op{KTAN} : \src{PRECISINTER} ($10^{-8}$), et $\Delta t$ pour les lois visqueuses.\\
— & \src{IRATE}, \src{PASREINI} & Itération non achevée (on repart alors du dernier état accepté) ; réinitialisation encore permise.\\
\end{glossaire}

\subsection{Le système résolu à chaque itération}
\label{sec:systeme}

À l'itération $i$, on résout le système à point-selle
\begin{equation}
\label{eq:systeme}
\begin{bmatrix}\vc K & \vc A^{\mathsf T}\\ \vc A & \vc 0\end{bmatrix}
\begin{Bmatrix}\delta\vc u\\ \vg\lambda^{(i)}\end{Bmatrix}
=
\begin{Bmatrix}\vc F^{ext}_{n+1} + \vc F^{suiv} - \vc F^{int}\big(\sig^{(i-1)}\big)\\[2pt] \vc u^{imp}_{n+1} - \vc A\,\vc u^{(i-1)}\end{Bmatrix}
\qquad(\text{\lig{1930-1982}}).
\end{equation}
Les déplacements sont résolus en correction, les multiplicateurs en total.
Pour le voir, partons de la linéarisation de~\eqref{eq:equilibre} autour de l'itéré $(\vc u^{(i-1)}, \vg\lambda^{(i-1)})$, avec l'opérateur $\vc K$ :
\begin{equation}
\vc K\,\delta\vc u + \vc A^{\mathsf T}\big(\vg\lambda^{(i)} - \vg\lambda^{(i-1)}\big) = \vc r^{(i-1)} = \vc F^{ext} + \vc F^{suiv} - \vc F^{int} - \vc A^{\mathsf T}\vg\lambda^{(i-1)} .
\end{equation}
Le terme $\vc A^{\mathsf T}\vg\lambda^{(i-1)}$ apparaît des deux côtés et s'élimine : on obtient la première ligne de~\eqref{eq:systeme}.
Le second membre ne contient donc pas les réactions, et $\vg\lambda^{(i)}$ est directement la réaction courante (\lig{2058-2069}).
Le déséquilibre vrai, qui sert au critère, est ensuite
\begin{equation}
\vc r^{(i)} = \vc F^{ext} + \vc F^{suiv} - \vc F^{int}\big(\sig^{(i)}\big) - \vc A^{\mathsf T}\vg\lambda^{(i)} .
\end{equation}

\subsection{L'algorithme des itérations}

Chaque itération enchaîne quatre étapes : éventuellement reconstruire $\vc K$, accélérer le second membre, résoudre, puis intégrer la loi et recalculer le déséquilibre.

\begin{itemize}
\item \textbf{Reconstruire $\vc K$.} En GD, on reconstruit la rigidité sur la configuration courante GEOM2 toutes les $p_K = \max(20, \src{DELTAITER})$ itérations, ou d'urgence (itération abandonnée, bascule HPP $\to$ GD, souci de comportement, contact à actualiser) (\lig{1540-1795}). Cette reconstruction perd l'amortissement et la perméabilité (constat~\ref{c:ktangent}). Plus de 20 reconstructions dans un sous-pas mettent la non-convergence à vrai (\lig{1552-1553}). Avec \src{K_TANGENT}, on recalcule la tangente toutes les \src{NITER_KTANGENT} itérations à partir de la deuxième (\lig{1799-1842}).
\item \textbf{Accélérer.} Aux itérations 4, 6, 8…, on retranche du second membre la correction d'\op{ACT3} (§ \ref{sec:act3}).
\item \textbf{Résoudre.} \op{RESO} résout~\eqref{eq:systeme} ; en présence de contact, il passe la main à la procédure \pr{UNILATER} (\annref{D}). La MAN remplace la correction de la première itération (\annref{B}), le pilotage indirect lui ajoute une composante pilotée (\annref{C}).
\item \textbf{Mettre à jour.} $\Delta\vc u \gets \Delta\vc u + \delta\vc u$, $\vc u^{(i)} = \vc u_n + \Delta\vc u$ et, en GD, $\mathrm{GEOM2} = \mathrm{FOR0} + \vc u^{(i)}$ (\lig{2070-2152}). Si le critère de l'itération précédente dépasse $10^{8}$, on abandonne l'itération une fois par sous-pas : on revient au dernier état accepté et l'on reconstruit $\vc K$ (\lig{2128-2132}).
\item \textbf{Intégrer et rééquilibrer.} \algoref{5} rend $\sig$, $\vc V$, $\vg\eps^{in}$ et $\Deps$. On calcule $\vc F^{int} = \op{BSIGMA}(\sig)$ sur GEOM2 (FOR0 en HPP), plus les termes dynamiques, et les forces suiveuses sur GEOM2 (\lig{2991-3151}). On forme $\vc r$ et $\vc R$ (\lig{3206-3225}), et \algoref{6} décide.
\end{itemize}

\begin{algobox}{4}{ETIQ : itérations d'équilibre}{lignes d'UNPAS}
\begin{algorithmic}[1]
\State $\Delta\vc u \gets 0$ ; $\vg\lambda \gets \vg\lambda_n$ ; $\Deps_{prec} \gets 0$ ; $\vc R \gets \vc R^{0}$
\For{$i = 1, 2, \dots$} \Comment{\src{ETIQ}, \lig{1483}}
  \If{GD et $i > 1$ et ($i - i_K$ multiple de $\max(20, \src{DELTAITER})$, ou urgence)} \Comment{\lig{1540-1795}}
    \State $\vc K \gets \vc K_e(\mathrm{GEOM2}) + \vc K_\sigma(\sig)$ + blocages + contact + masse \Comment{perd $\vc C$ et $\vc H$}
    \State \textbf{si} plus de 20 reconstructions dans le sous-pas : NONCONV $\gets$ vrai \Comment{\lig{1552-1553}}
  \EndIf
  \If{\src{K_TANGENT} et $i \ge 2$ et $i$ multiple de \src{NITER_KTANGENT}} \Comment{\lig{1799-1842}}
    \State $\vc K \gets \op{KTAN}(\sig, \vc V) + \vc K_\sigma$ + blocages \Comment{ou tangente SSTE, ou \op{KTAN} \src{PERT}}
  \EndIf
  \If{$i > 3$ et $\src{ITACC} \le 0$} \Comment{$i$ = 4, 6, 8…, \lig{1875-1921}}
    \State $\vc R \gets \vc R - \op{ACT3}$(4 derniers résidus) ; $\src{ITACC} \gets 2$
  \EndIf
  \State $(\delta\vc u, \vg\lambda) \gets \op{RESO}(\vc K, \vc R)$ \Comment{\lig{1930-1982} ; contact : \annref{D}}
  \State \textbf{si} MAN et $i = 1$ : $\delta\vc u \gets$ prédicteur asymptotique \Comment{\annref{B}}
  \State \textbf{si} pilotage indirect : $\delta\vc u \gets \delta\vc u + \Delta\eta\,\delta\vc u_{II}$ \Comment{\annref{C}}
  \State $\Delta\vc u \gets \Delta\vc u + \delta\vc u$ ; $\vc u \gets \vc u_n + \Delta\vc u$ ; en GD : $\mathrm{GEOM2} \gets \mathrm{FOR0} + \vc u$ \Comment{$\vg\lambda$ en total, \lig{2070-2152}}
  \State \textbf{si} pilotage automatique : normaliser $\Delta\vc u$, recalculer $\alpha$ \Comment{\annref{C}}
  \State \textbf{si} le critère précédent dépasse $10^{8}$, une seule fois : abandonner l'itération \Comment{\lig{2128-2132}}
  \State $(\sig, \vc V, \vg\eps^{in}, \Deps) \gets$ \algoref{5}
  \State $\vc F^{int} \gets \op{BSIGMA}(\sig)$ sur GEOM2 (FOR0 en HPP) $[+\,(4/\Delta t^2)\,\vc M\Delta\vc u + (2/\Delta t)\,\vc C\Delta\vc u]$ \Comment{\lig{2991-3151}}
  \State $\vc F^{suiv} \gets$ pression suiveuse et \pr{CHARMECA} sur GEOM2
  \State $\vc r \gets \vc F^{ext} + \vc F^{suiv} - \vc F^{int} - \vc A^{\mathsf T}\vg\lambda$ \Comment{\src{RESIDC}, pour le critère}
  \State $\vc R \gets [\vc F^{ext} + \vc F^{suiv} - \vc F^{int} ;\ \vc u^{imp} - \vc A\vc u]$ \Comment{\lig{3206-3225}}
  \State \textbf{selon} \algoref{6} : convergé $\to$ \Retourner{} (état, vrai) ; échec $\to$ \Retourner{} (état, faux)
  \State $\Deps_{prec} \gets \Deps$
\EndFor
\end{algorithmic}
\end{algobox}

\subsection{L'opérateur}
\label{sec:operateur}

L'opérateur des itérations n'est en général pas la tangente.
Par défaut, c'est la rigidité élastique, calculée avec le matériau de fin de pas et mémorisée dans \src{WTAB.RRRR} pour être réutilisée tant que rien n'oblige à la recalculer.
Le tableau~\ref{tab:operateur} récapitule les cas.

\begin{table}[ht]
\centering\small
\caption{L'opérateur des itérations selon les options.}
\label{tab:operateur}
\begin{tabularx}{\linewidth}{@{}>{\raggedright\arraybackslash}p{0.2\linewidth}>{\raggedright\arraybackslash}X>{\raggedright\arraybackslash}p{0.27\linewidth}>{\raggedright\arraybackslash}p{0.13\linewidth}@{}}
\toprule
Situation & Opérateur & Reconstruit & Source\\
\midrule
Défaut : HPP, matériau constant & $\vc K_e(\vc D_{n+1})$ + blocages & une fois par appel de \pr{PASAPAS} (\src{RRRR}, oubliée en fin d'appel), si la raideur ne contient pas de relations & \lig{711-771}\\
Matériau variable (hors $\alpha$ seul) & $\vc K_e(\vc D_{n+1})$ & chaque pas & \pr{PAS\_MODL} 345-350\\
Endommagement, viscodommage, céramique & raideur sécante endommagée, $\op{HOOK}(\vc V_n)$ & chaque pas & \lig{713-716}\\
Grands déplacements & $\vc K_e(\vc x_n) + \vc K_\sigma(\sig_n)$, puis $\vc K_e(\vc x^{(i)}) + \vc K_\sigma(\sig^{(i)})$ & chaque pas, puis toutes les $p_K$ itérations ou en urgence & \lig{961-966}, 1540-1795\\
\src{K_TANGENT} & itération 1 : \src{LASTKTAN}, sinon $\op{KTAN}(\sig_n, \vc V_n, \Delta t = 0)$ ; ensuite $\op{KTAN}(\sig^{(i)}, \vc V^{(i)})$ & chaque itération à partir de la 2\textsuperscript{e} & \lig{912-956}, 1799-1842\\
\src{K_TANGENT_PERT} & tangente par perturbation de l'intégration ($C_1 = 10^{-3}$, $C_2 = C_1/100$) & idem & \lig{1809-1821}\\
FEFP ; SSTE avec \src{K_TANGENT} & tangente cohérente rendue par \op{ECFE} ou \op{SSTE} & chaque itération & \lig{2179-2188}, 2457\\
Dynamique / consolidation & $+\,\tfrac4{\Delta t^2}\vc M + \tfrac2{\Delta t}\vc C$ / $-\theta_c\Delta t\,\vc H$ & chaque pas (\src{RECAOP}) & \lig{789-855}\\
Contact & + relations unilatérales \op{RFCO} & chaque sous-pas, à chaque reconstruction en GD ; contrôle à convergence & \lig{861-908}\\
Divers & \src{RIGIDITE_CONSTANTE}, rigidité augmentée, \src{ADDI_MATRICE} de \pr{CHARMECA}, advection & avec l'opérateur & \lig{737-762}, 1102-1104\\
\bottomrule
\end{tabularx}
\end{table}

\paragraph{Ce que cela implique pour la convergence.}
Avec un opérateur fixe $\vc K$ et une loi de tangente $\vc K_t$, l'erreur se propage, près de la solution, comme
\begin{equation}
\vc e^{(i)} \approx \big(\vc I - \vc K^{-1}\vc K_t\big)\,\vc e^{(i-1)} :
\end{equation}
la convergence est linéaire, d'autant plus lente que la tangente s'écarte de la rigidité élastique (plasticité étendue, adoucissement).
C'est ce qui justifie l'accélération \op{ACT3}, qui extrapole à partir des itérés précédents, et le seuil élevé d'itérations (250 par défaut).
En élasticité linéaire, $\vc K = \vc K_t$ et l'équilibre est atteint dès la première itération ; mais le critère en déformation exige au moins une seconde itération dès que le pas déforme de plus de $10^{-4}$ (§ \ref{sec:criteres}).

\paragraph{Rigidité géométrique.}
En GD, \src{K_SIGMA} est vrai par défaut avec l'hypothèse \src{QUADRATIQUE} (défaut des GD) ou si \src{STABILITE} est faux (\pr{PAS\_DEFA} \lig{1428-1439}).
$\vc K_\sigma(\sig)$ (\op{KSIGMA}) est alors ajoutée à l'opérateur initial (calculée avec $\sig_n$ sur GEOM1, sauf en prédicteur HPP) et à chaque reconstruction.
Au sens classique, c'est $\int_\Omega \vc G^{\mathsf T}\,\vg\Sigma\,\vc G\,\mathrm d\Omega$, où $\vc G$ contient les dérivées des fonctions de forme et $\vg\Sigma$ les contraintes ; elle rend compte de la raideur apportée par l'état de contrainte (flambage, membranes tendues).

\subsection{KTAN n'est pas la tangente algorithmique}
\label{sec:ktan}

L'opérateur \op{KTAN} (routines Fortran \texttt{KTANGE}, \texttt{KTANGA}, \texttt{DOHOT1}, \texttt{PENT1}) ne reçoit que l'état $(\sig, \vc V)$, le matériau, une précision \src{PREC} et un pas de temps : jamais l'incrément de déformation.
En chaque point de Gauss, il teste où se trouve la contrainte équivalente $\sigma_{eq}$ par rapport au seuil courant $\sigma_y(p)$ ($p$ : déformation plastique cumulée) :
\begin{equation}
\vc D_t =
\begin{cases}
\vc D & \text{si } \sigma_{eq} < (1 - \mathrm{PREC})\,\sigma_y(p),\\[2pt]
\vc D - \beta\,\vc s\otimes\vc s & \text{si } |\sigma_{eq} - \sigma_y(p)| \le \mathrm{PREC}\,\sigma_y(p),\\[2pt]
\text{erreur 275} & \text{si } \sigma_{eq} > (1 + \mathrm{PREC})\,\sigma_y(p),
\end{cases}
\end{equation}
où $\vc s$ est le déviateur des contraintes.
Pour le critère de von Mises, le source (\texttt{DOHOT1}) calcule
\begin{equation}
\beta = \frac{9E\,(E - E_t)}{2\,\sigma_{eq}^2\,(1+\nu)\,\big[3(E - E_t) + 2(1+\nu)E_t\big]},
\qquad E_t = \frac{H'E}{H' + E},
\end{equation}
$H'$ étant la pente de la courbe d'écrouissage en $p$ (\texttt{PENT1}).
Avec $G = E/2(1+\nu)$, cette expression vaut $\beta = 9G^2/\big((3G + H')\,\sigma_{eq}^2\big)$ : c'est la tangente élasto-plastique \emph{continue} classique.
La tangente algorithmique d'un retour radial dépend en outre de l'incrément (facteur $1 - 3G\Delta p/\sigma_{eq}^{trial}$ sur la partie déviatorique) ; \op{KTAN} ne peut pas la produire.

Trois conséquences.
Au-dessus de la surface, \op{KTAN} s'arrête sur l'erreur 275 au lieu de revenir à l'élastique, contrairement à ce qu'annonce sa notice.
Les modèles sans tangente codée reçoivent la matrice élastique.
UNPAS passe $\mathrm{PREC} = \src{PRECISINTER} = 10^{-8}$ ($5\cdot10^{-7}$ en visqueux) : la tangente dépend donc de la précision avec laquelle \op{COMP} a ramené $\sig$ sur la surface (constat~\ref{c:ktan8}).
Le source d'\pr{UNPAS} écrit d'ailleurs « cohérente » entre guillemets (\lig{930}, \lig{944}).
Seuls \op{KTAN} \src{PERT} (\texttt{KTAPER}, qui dérive numériquement l'intégration autour de l'incrément courant), \op{ECFE} et \op{SSTE} fournissent une tangente cohérente avec l'intégration.
Enfin, \src{K_TANGENT} est ignoré, avec un message, si le modèle n'est ni plastique ni visqueux (\lig{203-207}).

\subsection{L'accélération ACT3}
\label{sec:act3}

\paragraph{Ce que calcule ACT3.}
L'opérateur reçoit trois champs « primaux » $\vc D_0, \vc D_1, \vc D_2$ et quatre champs « duaux » $\vc F_0, \dots, \vc F_3$, l'indice 0 désignant le plus récent.
Il cherche la combinaison affine des duaux de norme minimale,
\begin{equation}
\min_{x_1, x_2, x_3}\;\Big\lVert \vc F_0 + \sum_{j=1}^{3} x_j\,\big(\vc F_j - \vc F_0\big)\Big\rVert_2 ,
\end{equation}
en résolvant par élimination de Gauss les équations normales
\begin{equation}
\sum_{k=1}^{3} \big[(\vc F_j - \vc F_0)\cdot(\vc F_k - \vc F_0)\big]\,x_k = -\,\vc F_0\cdot(\vc F_j - \vc F_0),\qquad j = 1, 2, 3,
\end{equation}
puis rend la combinaison correspondante des primaux,
\begin{equation}
\vc c = (x_1 + x_2 + x_3)\,\vc D_0 + (x_2 + x_3)\,\vc D_1 + x_3\,\vc D_2 .
\end{equation}
Si les primaux sont les pas successifs, $\vc D_j = \vc P_j - \vc P_{j+1}$, on a $\vc c = -\sum_j x_j(\vc P_j - \vc P_0)$ : c'est l'extrapolation d'Anderson à mémoire 3.
Plusieurs garde-fous complètent l'algorithme : réduction à deux puis une direction si le système est singulier (message 859), renoncement si rien n'est exploitable (860), « plan B » si le résidu oscille entre deux états (858), et refus d'une extrapolation qui retomberait sur l'un des points fournis.

\paragraph{Ce que lui passe UNPAS.}
Les duaux sont les quatre derniers résidus, sans les composantes \src{FLX}, pondérés par $\vc J$ :
$\vc F_k = \vc J\,\vc r^{(i-k)}$.
La pondération $1/\sqrt{V_{nœud}}$ rend la norme euclidienne comparable à une norme $L^2$ : sans elle, les nœuds des petites mailles, qui portent de petites forces nodales, compteraient moins.
Les primaux sont ces mêmes résidus, non pondérés, diminués de la correction appliquée à l'itération précédente : $\vc D_k = \vc r^{(i-k)} - \vc c^{(i-k-1)}$ (\lig{1882-1887}).
La correction est retranchée du second membre, $\vc R \gets \vc R - \vc c$ (\lig{1914}).
Comme l'opérateur $\vc K$ est fixe, retrancher $\vc c$ du second membre revient à modifier la correction de déplacement de $-\vc K^{-1}\vc c$.

\paragraph{Garde-fous de UNPAS.}
Avec \src{STABILITE} (vrai par défaut) ou en pilotage, on compare le travail du résidu sur la dernière correction avant et après accélération, $w_0 = \vc r\cdot\delta\vc u^{-}$ et $w = (\vc r - \vc c)\cdot\delta\vc u^{-}$.
Si $w/w_0 \le 0$, l'accélération est rétrograde et on l'annule ; si $w/w_0$ dépasse une limite $L = 10^{5}$ ($L = 1$ quand $w_0 < 0$), on la multiplie par $L/(w/w_0)$ (\lig{1896-1912}).

\paragraph{Quand.}
\src{ITACC} est décrémenté au début de chaque itération et remis à 2 après chaque accélération : \op{ACT3} s'applique aux itérations 4, 6, 8…
Il est suspendu, en remettant \src{ITACC} à 4, après certains événements : souci de résolution (\lig{2027}), itération abandonnée (\lig{1507}), convergence refusée pour bascule HPP $\to$ GD, souci de comportement ou contact à actualiser (\lig{3408}), reconstruction de $\vc K$ à l'itération 2 (\lig{1792}), relaxation (\lig{3453}), et, en GD, non-convergence détectée plus de 3 itérations après la dernière reconstruction (\lig{3447}).
Selon que la remise à 4 a lieu avant ou après le décrément, la suspension dure 3 ou 4 itérations.
\op{ACT3} reste actif avec \src{K_TANGENT}, sauf en FEFP et en SSTE (\lig{210-212}).
% ============================================================================
\section{L'intégration du comportement et les configurations}
\label{sec:comportement}
% ============================================================================

\begin{glossaire}
$N$ & \src{nsoincr}, \src{nsoincrn} & Nombre de sous-incréments de l'itération (\src{SOUS_INCREMENT}, 1 par défaut), et valeur proposée pour l'itération suivante.\\
$\Delta\eps_{max}$ & \src{DEPS_MAX} & Incrément de déformation visé par sous-incrément en GD, $5\cdot10^{-3}$ ; avec \src{RECALCUL_SOUSINC} (vrai sauf si \src{SOUS_INCREMENT} est donné).\\
$\Delta\vc u_k$ & \src{ZDEPTI} & $\Delta\vc u/N$ ; \src{ZDEPTI2} et \src{ZDEPTI2n} valent $\pm\Delta\vc u_k/2$.\\
$\vc x_{k-\frac12}$, $\vc x_k$ & \src{GEOMM}, \src{GEOMF} & Configurations de mi- et de fin de sous-incrément : $\vc x_n + (k-\tfrac12)\Delta\vc u/N$ et $\vc x_n + k\,\Delta\vc u/N$.\\
$\Deps_k$, $\Deps$ & \src{DEPSTa}, \src{DEPST} & Incrément de déformation du sous-incrément, et cumul sur le pas.\\
$\sig_{k-1}$, $\vg\eps_{k-1}$, $\vc V_{k-1}$, $\vg\eps^{in}_{k-1}$ & \src{ZSIG0a}, \src{ZDEF0a}, \src{ZVAR0a}, \src{ZDEI0a} & État au début du sous-incrément.\\
$\Delta\sig_k$, $\Delta\sig$ & \src{DSIGTa}, \src{DSIGT} & Incrément de contrainte du sous-incrément, ramené sur FOR0, et cumul.\\
$t_{k-1}$, $t_k$ & \src{tsodeb}, \src{tsofin} & $t_{déb} + \alpha\,\Delta t\,(k-1)/N$ et $t_{déb} + \alpha\,\Delta t\,k/N$.\\
$T(t_{k-1})$, $T(t_k)$ & \src{TETDE}, \src{TETDF} & $T_n + \tfrac{k-1}{N}(T_{n+1} - T_n)$ et $T_n + \tfrac kN(T_{n+1} - T_n)$.\\
— & \src{LAG_TOT} & Configuration de l'intégration : 1 \src{TOTAL}, 2 \src{FIN_PAS}, 3 \src{MI_PAS} (défaut en GD), 0 hors GD et en endommagement.\\
— & \src{HYPDEF}, \src{MEPSIL} & Hypothèse de déformation (\src{HYPOTHESE_DEFORMATIONS}) ; UNPAS la recopie dans l'option globale \src{OPTI 'EPSI'} (\lig{174}), que lit \op{CONF}.\\
— & \src{ZGRDU0}, \src{D_GR_U} & Gradient du déplacement total sur FOR0, et son incrément : transmis à \op{COMP} pour les lois \src{UTILISATEUR} en GD.\\
$\Delta\eps_{max}^{pas}$ & \src{DPSMAX} & $\norminf{\Deps}$, sur les composantes de déformation.\\
\end{glossaire}

\subsection{Les sous-incréments}

Quand la loi n'est pas élastique linéaire, on découpe l'incrément de l'itération en $N$ sous-incréments égaux, $\Delta\vc u_k = \Delta\vc u/N$, intégrés l'un après l'autre (boucle \src{sousinc}, \lig{2224}).
En GD, avec \src{RECALCUL_SOUSINC}, $N$ est ajusté d'une itération à l'autre pour que chaque sous-incrément déforme d'environ \src{DEPS_MAX} (\lig{2367-2370}) :
\begin{equation}
N \gets \min\!\Big(20,\ \max\big(1,\ N - 3,\ \big\lfloor \Delta\eps_{max}^{pas}/\Delta\eps_{max} \big\rfloor + 1\big)\Big).
\end{equation}
$N$ peut donc baisser de 3 au plus d'une itération à l'autre.
Le temps et la température sont interpolés dans le sous-incrément :
\begin{equation}
t_k = t_{déb} + \alpha\,\Delta t\,\frac kN,\qquad T(t_k) = T_n + \frac kN\big(T_{n+1} - T_n\big)\qquad(\text{\lig{2467-2485}}).
\end{equation}
Le temps est mis à l'échelle par $\alpha$ (cible réduite après une relaxation), la température ne l'est pas (constat~\ref{c:details}).

\subsection{L'incrément de déformation}

\paragraph{En GD avec une hypothèse non linéaire.}
On se place sur la configuration de mi-sous-incrément $\vc x_{k-\frac12}$ et l'on calcule (\lig{2245-2299})
\begin{equation}
\label{eq:pointmilieu}
\Deps_k = \vc E_{\vc x_{k-\frac12}}\!\big(\tfrac12\Delta\vc u_k\big) - \vc E_{\vc x_{k-\frac12}}\!\big(-\tfrac12\Delta\vc u_k\big),
\qquad
\vc E(\vc a) = \tfrac12\big(\nabla\vc a + \nabla\vc a^{\mathsf T} + \nabla\vc a^{\mathsf T}\nabla\vc a\big),
\end{equation}
les gradients étant pris par rapport à $\vc x_{k-\frac12}$.
$\vc E$ est la déformation de Green-Lagrange que calcule \op{EPSI} \src{'QUAD'} : le Fortran ajoute à la partie linéaire (\texttt{BST}) le terme $\tfrac12\nabla\vc a^{\mathsf T}\nabla\vc a$ (\texttt{BST2}).
Ce terme est pair en $\vc a$ ; il s'annule exactement dans la différence, et il reste
\begin{equation}
\Deps_k = \tfrac12\big(\nabla_{\vc x_{k-\frac12}}\Delta\vc u_k + \nabla_{\vc x_{k-\frac12}}\Delta\vc u_k^{\mathsf T}\big) = \operatorname{sym}\nabla_{\vc x_{k-\frac12}}\Delta\vc u_k .
\end{equation}
C'est la règle du point milieu appliquée au taux de déformation $\vc d = \operatorname{sym}\nabla_{\vc x}\vc v$ : $\int\vc d\,\mathrm dt \approx \operatorname{sym}\nabla_{\vc x_{mid}}\Delta\vc u$, qui s'annule exactement pour une rotation rigide (les deux termes sont alors égaux) et reste du second ordre.
Le calcul est le même pour \src{QUADRATIQUE}, \src{TRUESDELL}, \src{JAUMANN} et \src{UTILISATEUR} ; seul le transport diffère (§ \ref{sec:conf}).
En endommagement, viscodommage et céramique, on calcule directement \op{EPSI} \src{'LINE'} sur $\vc x_{k-\frac12}$, ce qui revient au même pour les éléments massifs.
Si le calcul échoue (jacobien négatif, déformation supérieure à 1000), on prend la déformation linéaire sur la configuration de début de pas, avec un message ; en augmentation automatique, on annule plutôt le sous-incrément (§ \ref{sec:augmentation}).

\paragraph{Sinon.}
$\Deps_k = \operatorname{sym}\nabla\Delta\vc u_k$ sur la configuration de début de pas (\lig{2303-2320}), soit FOR0 en HPP.

\subsection{Transport des champs et configurations}
\label{sec:conf}\label{sec:configs}

\paragraph{Ce que font FORM et CONF.}
Une configuration est un jeu de coordonnées nodales ; \op{FORM} active une configuration, ou la crée en ajoutant un déplacement aux coordonnées courantes.
Chaque champ par élément mémorise la configuration sur laquelle il a été calculé.
\op{CONF} (Fortran \texttt{CONFIG}, \texttt{PIOCAP}) pousse un champ de sa configuration $\vc x$ vers la configuration courante $\vc y$, avec $\vc F = \partial\vc y/\partial\vc x$ et $J = \det\vc F$ :
\begin{equation}
\label{eq:conf}
\sig \;\mapsto\; J^{-1}\,\vc F\,\sig\,\vc F^{\mathsf T},
\qquad
\vg\eps \;\mapsto\; \vc F^{-\mathsf T}\,\vg\eps\,\vc F^{-1}
\qquad(\texttt{PIOCAF})
\end{equation}
pour les hypothèses \src{LINEAIRE}, \src{QUADRATIQUE} et \src{TRUESDELL}.
Avec \src{JAUMANN}, le champ est seulement tourné, $\sig \mapsto \vc R\,\sig\,\vc R^{\mathsf T}$, $\vc R$ venant de la décomposition polaire $\vc F = \vc R\vc U$ (\texttt{JAUCAU}).
Avec \src{UTILISATEUR}, ou pour une loi UMAT, il est recopié tel quel.
Ce n'est donc pas un simple changement d'étiquette, contrairement à ce que laisse entendre sa notice.

\paragraph{La chaîne de transport dans UNPAS.}
\begin{enumerate}
\item $\Deps_k$, calculé sur $\vc x_{k-\frac12}$, est ramené sur FOR0 par \op{CONF} (\lig{2286-2296}). Avec $\vc F_{k-\frac12} = \partial\vc x_{k-\frac12}/\partial\vc X$, la formule~\eqref{eq:conf} donne $\vc F_{k-\frac12}^{\mathsf T}\,\Deps_k\,\vc F_{k-\frac12}$ : c'est l'incrément de déformation de Green-Lagrange au point milieu.
\item $\sig$, $\vg\eps$ et $\Deps_k$ sont ensuite poussés sur la configuration de l'intégration choisie par \src{LAGRANGIEN} (\lig{2334-2352}) : FOR0 (\src{TOTAL}), $\vc x_k$ (\src{FIN_PAS}) ou $\vc x_{k-\frac12}$ (\src{MI_PAS}, défaut en GD).
\item \op{COMP} intègre la loi sur cette configuration ; $\Delta\sig_k$ est ramené sur FOR0 (\lig{2918-2928}) et cumulé.
\item Pour les forces internes, $\sig$ est poussée de FOR0 sur GEOM2 avant \op{BSIGMA} (\lig{2995-3006}).
\end{enumerate}
La contrainte rangée sur FOR0 est donc de type Piola-Kirchhoff~2, $\vc S = J\,\vc F^{-1}\sig\,\vc F^{-\mathsf T}$, et \op{BSIGMA} reçoit la contrainte de Cauchy correspondante.
En \src{LAGRANGIEN TOTAL}, \op{COMP} travaille en $(\vc S, \vc E)$ sur FOR0 : une loi élastique y devient Saint-Venant-Kirchhoff.
En \src{MI_PAS}, \op{COMP} reçoit la contrainte de Cauchy poussée sur $\vc x_{k-\frac12}$ et l'incrément $\operatorname{sym}\nabla_{\vc x_{k-\frac12}}\Delta\vc u_k$ : c'est une loi incrémentale en Cauchy, avec un transport convectif de type Truesdell entre configurations.
Avec \src{JAUMANN}, le transport se réduit à une rotation, faite en deux \op{CONF} successifs (FOR0 vers GEOM1, puis vers la cible) ; avec \src{UTILISATEUR}, il n'y en a pas, mais \op{COMP} reçoit les gradients de déplacement de début et de fin de sous-incrément, dont la loi UMAT peut tirer $\vc F$ (\lig{2505-2515}).

\begin{figure}[ht]
\centering
\begin{tikzpicture}[x=1cm,y=1cm,>={Stealth[length=2mm]},
  q/.style={draw=meca,fill=meca!8,line width=0.6pt},
  lab/.style={font=\sffamily\scriptsize,align=center,text width=3.4cm}]
  \draw[q] (0,0) -- (1.6,0) -- (1.6,1.6) -- (0,1.6) -- cycle;
  \node[lab] at (0.8,-0.75) {\textbf{FOR0} $= \vc X$\\ stockage de $\sig$ (PK2), $\Deps$, matériau ; \src{TOTAL}};
  \draw[q] (4,0) -- (5.7,0.1) -- (6.0,1.75) -- (4.3,1.65) -- cycle;
  \node[lab] at (5.0,-0.75) {\textbf{GEOM1} $= \vc x_n$\\ $\vc K$ et $\vc K_\sigma$ initiaux, $\vc F^{int}_n$, $\vc F^{*}$};
  \draw[q] (8.2,0) -- (9.95,0.15) -- (10.35,1.85) -- (8.6,1.7) -- cycle;
  \node[lab] at (9.25,-0.75) {\textbf{GEOMM} $= \vc x_{k-\frac12}$\\ $\Deps_k$ ; \op{COMP} en \src{MI_PAS}};
  \draw[q] (12.4,0) -- (14.2,0.2) -- (14.7,1.95) -- (12.9,1.75) -- cycle;
  \node[lab] at (13.5,-0.75) {\textbf{GEOM2} $= \vc X + \vc u^{(i)}$\\ \op{BSIGMA}, forces suiveuses, $\vc K$ reconstruite};
  \draw[->,gray] (1.8,0.8) -- node[above,font=\scriptsize,text=black] {$\vc u_n$} (3.9,0.8);
  \draw[->,gray] (6.1,0.9) -- node[above,font=\scriptsize,text=black] {$(k-\tfrac12)\Delta\vc u/N$} (8.1,0.9);
  \draw[->,gray] (10.45,1.0) -- node[above,font=\scriptsize,text=black] {reste de $\Delta\vc u$} (12.3,1.0);
\end{tikzpicture}
\caption{Les configurations en grands déplacements. Les flèches désignent ici la transformation d'une configuration à la suivante, et non un enchaînement de calculs. GEOMF $= \vc x_k$ (fin de sous-incrément, \src{FIN_PAS}) n'est pas dessinée.}
\label{fig:configs}
\end{figure}

\begin{table}[ht]
\centering\small
\caption{Les configurations et leur usage.}
\label{tab:configs}
\begin{tabularx}{\linewidth}{@{}>{\raggedright\arraybackslash}p{0.17\linewidth}>{\raggedright\arraybackslash}p{0.3\linewidth}>{\raggedright\arraybackslash}X@{}}
\toprule
Configuration & Définition & Utilisée pour\\
\midrule
\src{GEOREF0} = \src{FOR0} & au premier appel de \pr{PASAPAS} & matériau à $T_I$, stockage de $\sig$ et $\Deps$, lagrangien total (\src{LAG_TOT} 1, et 0 en endommagement)\\
\src{GEOM1} = \src{GE0_DEB} & HPP : FOR0. GD et FEFP : $\vc X + \vc u_n$, actualisée après chaque pas & $\vc K$ et $\vc K_\sigma$ initiaux, $\vc F^{int}_n$, pseudo-forces, thermique de la passe 1\\
\src{GEOMM} & $\vc x_n + (k-\tfrac12)\,\Delta\vc u/N$ & $\Deps_k$ ; \op{COMP} en \src{REACTUALISE} ou \src{MI_PAS} (\src{LAG_TOT} 3, défaut en GD)\\
\src{GEOMF} & $\vc x_n + k\,\Delta\vc u/N$ & \op{COMP} en \src{FIN_PAS} (\src{LAG_TOT} 2)\\
\src{GEOM2} & $\mathrm{FOR0} + \vc u^{(i)}$, à chaque itération & \op{BSIGMA}, pression suiveuse (sans raideur de chargement), reconstruction de $\vc K$, $\vc K_\sigma$, \op{KTAN}, contact ; devient \src{WTAB.FOR}\\
\bottomrule
\end{tabularx}
\end{table}

\paragraph{Le prédicteur HPP.}
Avec \src{PREDICTEUR 'HPP'}, \pr{UNPAS} converge d'abord avec les forces internes calculées sur la configuration de début de pas et des précisions divisées par 10, puis il « passe en grands déplacements » et continue (\lig{3401-3412}).
Le contact force les grands déplacements.

\subsection{L'intégration de la loi}

\paragraph{Élasticité linéaire.}
Sans plasticité ni viscosité, l'incrément de contrainte est direct (\lig{2422-2432}) :
\begin{equation}
\Delta\sig_k = \vc D_{n+1}\Big(\Deps_k - \frac{\alpha}{N}\,\Deps^{*}\Big).
\end{equation}

\paragraph{Loi non linéaire.}
On retire d'abord la part de pseudo-déformation du sous-incrément, $\Deps_k \gets \Deps_k - \tfrac{\alpha}{N}\Deps^{*}$ (\lig{2443-2447}), puis on appelle \op{COMP} avec deux états (\lig{2467-2539}) :
\begin{align}
\text{initial} &: \big(\vg\eps_{k-1},\ \sig_{k-1},\ \vc V_{k-1},\ \vg\eps^{in}_{k-1},\ T(t_{k-1}),\ t_{k-1},\ \vc D_n\big),\\
\text{final} &: \big(\vg\eps_{k-1} + \Deps_k,\ T(t_k),\ t_k,\ \vc D_{n+1}\big),
\end{align}
et \op{COMP} rend $\sig_k$, $\vc V_k$, $\vg\eps^{in}_k$.
En non local, on régularise ensuite $\vc V$ et l'on rappelle \op{COMP} (\annref{A}).

\subsection{Deux intégrateurs particuliers : FEFP et SSTE}

\paragraph{FEFP (\src{FEFP_FORMULATION}).}
Plasticité multiplicative, $\vc F = \vc F_e\vc F_p$, pour les modèles \src{VMT_FEFP}, \src{RHMC_FEFP}, \src{POWDER_FEFP} et \src{POWDERCAP_FEFP}.
\pr{PAS\_DEFA} coupe les grands déplacements « classiques » et force \src{K_SIGMA}, \src{K_TANGENT} et \src{RECARI}, et coupe \src{SUBSTEPPING} (\pr{PAS\_DEFA} \lig{1479-1485}).
L'opérateur \op{ECFE} intègre la loi directement à partir du déplacement total (ou de l'incrément en lagrangien réactualisé) (\lig{2174-2205}).
Dans le Fortran (\texttt{FEFP2}, \texttt{ADFEFP}), le retour est exponentiel : on passe en déformations logarithmiques principales du tenseur élastique $\vc b_e$, $\eps_a = \tfrac12\ln\lambda_a^2$, on y fait le retour, puis on reconstruit $\vc b_e = \sum_a e^{2(\eps_a - \Delta p_a)}\,\vc n_a\otimes\vc n_a$, avec une recherche linéaire sauf pour un des modèles de poudre.
\op{ECFE} rend $\sig$, $\vc V$, $\vc b_e$ (rangé dans \src{DEFORMATIONS_INELASTIQUES}) et la tangente cohérente : on obtient un Newton complet, sans \op{ACT3}, et le critère en déformation porte alors sur $\vc b_e$.

\paragraph{SSTE (\src{SUBSTEPPING}).}
Intégration avec sous-découpage au point de Gauss pour les modèles J2, \src{RH_COULOMB} et \src{MRS_LADE} (\lig{2456-2462}).
Dans le Fortran (\texttt{SBSTEP}), un sous-pas raté est coupé en deux, et la tangente cohérente est composée sous-pas par sous-pas par la règle de dérivation en chaîne.
Cette tangente ne sert d'opérateur qu'avec \src{K_TANGENT} ; sinon, l'opérateur reste la rigidité élastique.
\src{NMAXSUBSTEPS} vaut 1 par défaut : sans réglage, il n'y a pas de vrai sous-découpage.
SSTE repart de $\vc V_n$ à chaque sous-incrément $k$, dont seule la dernière composante est mise à jour (\lig{2458-2460}).

\begin{algobox}{5}{Intégration du comportement}{lignes d'UNPAS}
\begin{algorithmic}[1]
\Require $\Delta\vc u$, $\vc x_n$, état de début ($\sig_n$, $\vc V_n$, $\vg\eps_n$, $\vg\eps^{in}_n$), $T_n$, $T_{n+1}$, $t_{déb}$, $\alpha$
\State \textbf{si} FEFP : $(\vc K_t, \sig, \vc V, \vc b_e) \gets \op{ECFE}(\vc u$, ou $\Delta\vc u$ en réactualisé$)$ ; $\vc K \gets \vc K_t$ + blocages ; \Retourner \Comment{\lig{2174-2205}}
\State $N \gets \src{SOUS_INCREMENT}$ ; en GD, ajusté d'une itération à l'autre pour $\norminf{\Deps_k} \le \src{DEPS_MAX}$ \Comment{\lig{2217, 2367-2370}}
\For{$k = 1 \dots N$} \Comment{\src{sousinc}, \lig{2224}}
  \State $\Delta\vc u_k \gets \Delta\vc u/N$
  \If{GD et hypothèse non linéaire} \Comment{\lig{2245-2299}}
    \State $\vc x_{mid} \gets \vc x_n + (k - \tfrac12)\,\Delta\vc u/N$
    \State $\Deps_k \gets \vc E_{\vc x_{mid}}(\Delta\vc u_k/2) - \vc E_{\vc x_{mid}}(-\Delta\vc u_k/2)$ \Comment{$= \operatorname{sym}\nabla_{\vc x_{mid}}\Delta\vc u_k$}
    \State \textbf{si} jacobien négatif : $\Deps_k \gets$ déformation linéaire sur $\vc x_n$
  \Else
    \State $\Deps_k \gets \operatorname{sym}\nabla\Delta\vc u_k$ sur la configuration de début de pas \Comment{\lig{2303-2320}}
  \EndIf
  \State transporter $\Deps_k$ sur FOR0 ; en GD, $\sig$, $\vg\eps$ et $\Deps_k$ sur la configuration \src{LAGRANGIEN} \Comment{\lig{2334-2352}}
  \If{élastique linéaire}
    \State $\Delta\sig_k \gets \vc D_{n+1}\,(\Deps_k - \alpha\,\Deps^{*}/N)$ \Comment{\lig{2422-2432}}
  \Else
    \State $\Deps_k \gets \Deps_k - \alpha\,\Deps^{*}/N$ \Comment{\lig{2443-2447}}
    \If{\src{SUBSTEPPING}}
      \State $(\sig, \vc V, \vg\eps^{in}, \vc K_t) \gets \op{SSTE}(\sig, \vc V_n, \Deps_k)$ \Comment{\lig{2456-2462}}
    \Else
      \State initial $\gets (\vg\eps, \sig, \vc V, \vg\eps^{in}, T(t_{k-1}), t_{k-1}$, matériau de début$)$
      \State final $\gets (\vg\eps + \Deps_k, T(t_k), t_k$, matériau à $T_I)$
      \State $(\sig, \vc V, \vg\eps^{in}) \gets \op{COMP}(\text{initial}, \text{final})$ \Comment{\lig{2467-2539}}
      \State \textbf{si} non local : régulariser $\vc V$, puis rappeler \op{COMP} \Comment{\annref{A}}
    \EndIf
  \EndIf
  \State $\sig \gets \sig + \Delta\sig_k$, ramenée sur FOR0 \Comment{\lig{2918-2928}}
\EndFor
\State \Retourner{} $\sig$, $\vc V$, $\vg\eps^{in}$, $\Deps = \sum_k \Deps_k$
\end{algorithmic}
\end{algobox}
% ============================================================================
\section{Les critères de convergence et de non-convergence}
\label{sec:criteres}
% ============================================================================

\begin{glossaire}
$F_{ref}$, $M_{ref}$ & \src{XDENO}, \src{XDENOM} & Force et moment de référence, calculés au premier sous-pas et gardés pour tout le pas.\\
— & \src{MZDEP1M}, \src{MZFM} & $\norminf{\delta\vc u^{0}}$ sur les translations, et $\norminf{\vc F^{int}_n + \vc F^{\eps}}$ : ingrédients de $F_{ref}$.\\
$\lVert \vc r\rVert_{\infty,F}$ & \src{XNUMF} & Maximum du résidu sur les composantes de force, hors moments, fluide (\src{FP}, \src{FPQ}, \src{FTP}) et \src{FLX}.\\
— & \src{XNUMG} & Maximum de $|\vc A\cdot[\Delta\vc u^{(i)} - \Delta\vc u^{(i-1)} ; \vg\lambda^{(i)} - \vg\lambda^{(i-1)}]|$ : variation des réactions, et des valeurs de liaison sur les \src{FLX}.\\
$e_F$, $e_M$ & \src{XCONV}, \src{XCONVM} & Critères en force et en moment.\\
— & \src{TSTMOM} & Vrai si le modèle porte des moments (coques, poutres, joints, poreux) : le critère en moment s'ajoute.\\
— & \src{FTOL}, \src{MTOL} & Tolérances absolues optionnelles ; elles remplacent alors $F_{ref}$ ou $M_{ref}$ par 1.\\
$e_\eps$ & \src{DEPSTDM} & $\norminf{\Deps^{(i)} - \Deps^{(i-1)}}$ sur les composantes membranaires et axiales (\src{MLDEFOR}), courbures exclues.\\
$\epsilon_F$, $\epsilon_\eps$, $\epsilon_M$ & \src{ZPREC}, \src{ZPRECEPS}, \src{ZPRECM} & \src{PRECISION}, \src{PRECDEFO}, \src{PRECFLEX} : $10^{-4}$ chacun ; divisés par 10 en prédicteur HPP.\\
— & \src{PASTEST} & Vrai si le test est sauté à cette itération.\\
— & \src{PASUNIL} & Vrai si \pr{UNILATER} n'a pas trouvé d'ensemble actif stable.\\
— & \src{iCONV}, \src{iRECA}, \src{isoucomp} & Critères satisfaits ; contact à actualiser ; souci signalé par \op{COMP}.\\
— & \src{NONCONV} & Non-convergence détectée ; on cherche alors à la déclarer.\\
$e_F^{(i)}$ & \src{TABCONV.i} & Historique de $e_F$, mis à 1 en tête d'itération.\\
$d$ & \src{DLTAIT} & Fenêtre de détection, \src{DELTAITER} (10) au départ.\\
$e_{ref}$ & \src{XCONVREF} & $0{,}99\,e_F^{(i-d)}$ : valeur que $e_F$ doit battre.\\
$I_{max}$ & \src{ZMAXIT} & Plafond courant, \src{MAXITERATION} au départ ; non réinitialisé entre sous-pas.\\
$\epsilon_{nc}$ & \src{zprecnc} & Seuil sur $\norminf{\Deps}$ pour déclarer la non-convergence ; part de $10^{-5}$.\\
— & \src{DPSMAX}, \src{DPSMAXP} & $\norminf{\Deps}$ à l'itération courante et à la précédente.\\
— & \src{ZNITE} & Nombre d'itérations des sous-pas convergés ; devient \src{NOMBRE_ITERATIONS}.\\
\end{glossaire}

\subsection{Les références}
\label{sec:fref}

Un critère relatif a besoin d'une échelle de force.
\pr{UNPAS} la calcule au premier sous-pas, à partir de la première solution $(\delta\vc u^{0}, \vg\lambda^{0})$, et la garde pour tout le pas (\lig{1334-1355}) :
\begin{equation}
\label{eq:fref}
F_{ref} = \max\left(
\frac{\big|\delta\vc u^{0}\cdot\vc F^{ext}_{n+1} - \vg\lambda^{0}\cdot\Delta\vc u^{imp}\big|}{\norminf{\delta\vc u^{0}}} + 2\,\norminf{\vc F^{int}_n + \vc F^{\eps}},\ \
\norminf{\vc F^{ext}_{n+1}},\ \
\norminf{\vc K\,\vc u_n}^{\text{(GD)}}\right).
\end{equation}
Le premier terme est le travail du chargement dans le premier incrément, divisé par sa plus grande composante : une force caractéristique du pas, qui reste significative quand le chargement est un déplacement imposé ($\vc F^{ext} = 0$).
Le deuxième ajoute deux fois les forces déjà présentes, $\vc F^{\eps}$ regroupant les forces des déformations imposées totales en fin de pas (thermique et \src{DEFI}) et la pseudo-force $\vc F^{*}$ (\lig{1065-1071}, \lig{1340}).
En GD, la force élastique de la configuration atteinte, $\vc K\vc u_n$, sert de plancher.
$F_{ref}$ est remplacée par 1 si elle est quasi nulle ou indéfinie.
La référence en moment $M_{ref}$ divise le même terme de travail par le maximum des rotations de $\delta\vc u^{0}$.

\subsection{Les grandeurs testées}

Après l'itération $i$, on calcule (\lig{3229-3258}, \lig{3321-3322})
\begin{align}
e_F &= \frac{\max\Big(\lVert \vc r^{(i)}\rVert_{\infty,F},\ \norminf{\vc A^{\mathsf T}\big(\vg\lambda^{(i)} - \vg\lambda^{(i-1)}\big)}\Big)}{F_{ref}},\\[2pt]
e_M &= \frac{\lVert \vc r^{(i)}\rVert_{\infty,M}}{M_{ref}},\qquad
e_\eps = \norminf{\Deps^{(i)} - \Deps^{(i-1)}} .
\end{align}
\begin{itemize}
\item Les normes sont des maximums sur tous les nœuds et toutes les composantes. $\lVert \cdot\rVert_{\infty,F}$ exclut les moments, le fluide et les \src{FLX}.
\item Le second terme de $e_F$ exige que les réactions se stabilisent ; sur les \src{FLX}, il porte aussi sur la variation de $\vc A\vc u$, nulle une fois les liaisons satisfaites (\lig{3230}).
\item Si aucun moment n'est appliqué, $e_M = 0$ ; si les forces appliquées sont négligeables devant les moments (rapport inférieur à $10^{-12}$), $e_F = 0$.
\item $e_\eps$ est un écart \emph{absolu}, sans dimension puisque c'est une déformation. À l'itération 1, $\Deps^{(0)} = 0$ : $e_\eps$ vaut tout l'incrément de déformation, si bien qu'un pas qui déforme de plus de $10^{-4}$ ne peut pas converger en une itération, même en élasticité linéaire (constat~\ref{c:deuxit}).
\end{itemize}

\subsection{La convergence}

Si le test n'est pas sauté, la convergence est acquise quand
\begin{equation}
e_F < \epsilon_F,\qquad e_\eps < \epsilon_\eps,\qquad \text{et}\quad e_M < \epsilon_M \ \text{si le modèle porte des moments}
\qquad(\text{\lig{3330-3337}}).
\end{equation}
Le test est sauté à l'itération 1 après le prédicteur $\vc F^{NL}$, avec frottement tant que $i < 2$, quand le comportement ou la résolution signale un souci, et dans certains cas de pilotage.
Même satisfaits, les critères ne valent convergence qu'après quatre vérifications, dans cet ordre :
\begin{enumerate}
\item le contact doit être résolu, sinon échec « contact » ;
\item hors pilotage, $\alpha$ doit valoir 1 : sinon la cible réduite est atteinte, et l'on sort en échec « relaxation » pour repartir de là en sous-pas ;
\item l'augmentation automatique et la réduction du chargement doivent être inactives (échec « augmentation », § \ref{sec:augmentation}) ;
\item il ne doit pas rester de prédicteur HPP, de souci de comportement, ni de réactions de contact qui changent de plus de \src{PRECISION} quand on recalcule les relations sur GEOM2 ; sinon on bascule en GD ou l'on reconstruit $\vc K$, et l'on continue (\lig{3401-3412}).
\end{enumerate}

\subsection{La non-convergence : détection, puis déclaration}

La non-convergence se traite en deux temps.
On la \emph{détecte} quand le critère stagne ; on la \emph{déclare} ensuite, quand l'incrément de déformation est assez petit pour qu'un nouvel état de départ ait un sens.

\paragraph{Détection.}
On compare le critère courant à celui d'il y a $d$ itérations (\lig{3429-3432}) :
\begin{equation}
e_{ref} = \begin{cases}0{,}99\;e_F^{(i-d)} & \text{si } i > d,\\ +\infty & \text{sinon.}\end{cases}
\end{equation}
Si le test n'est pas sauté, que la non-convergence n'est pas déjà détectée, que $i \ge 2$, hors pilotage et contact résolu, on la détecte quand $i \ge I_{max}$ ou $e_F > e_{ref}$ (\lig{3437-3488}).
Autrement dit, le critère doit baisser d'au moins 1\,\% toutes les $d$ itérations.
On pose alors $I_{max} \gets 3\,i$, et $\alpha \gets \src{RELAXATION_NONCONV}$ ; si $\alpha < 1$, on réduit la cible à $t_{déb} + \alpha\,\Delta t$ et l'on réévalue chargement et matériau.

La fenêtre $d$ vaut d'abord \src{DELTAITER}.
Une bascule HPP $\to$ GD, un souci de comportement ou un contact à actualiser la portent à $\src{DELTAITER} + i$ (\lig{3410}), sans retour en arrière, même aux sous-pas suivants (\lig{123}).
Deux autres chemins mettent \src{NONCONV} à vrai : plus de 20 reconstructions de $\vc K$ dans un sous-pas en GD (\lig{1552-1553}), et, en augmentation automatique, la quatrième annulation d'itération dans un sous-pas, avec $I_{max} \gets i + 5$ (\lig{2390-2396}) ; à partir de la cinquième, $I_{max} \gets i - 1$, ce qui force la déclaration (\lig{2407}).

\paragraph{Déclaration.}
Une fois la non-convergence détectée, le seuil $\epsilon_{nc}$ double à chaque itération au-delà de la quinzième (\lig{3491-3496}).
On déclare l'échec, « NON CONVERGENCE », si le test n'est pas sauté, si $\norminf{\Deps} < \epsilon_{nc}$, et si
\begin{equation}
i \ge I_{max}\quad\text{ou}\quad\big(e_F > e_{ref}\ \text{et}\ \norminf{\Deps^{(i)}} \ge \tfrac14\norminf{\Deps^{(i-1)}}\big)
\qquad(\text{\lig{3498-3506}}).
\end{equation}
Le seuil qui double laisse le temps aux itérations de converger tant que l'incrément de déformation est grand ; il finit toujours par le dépasser.

\begin{remarque}[Exemple]
Prenons un pas avec $\norminf{\Deps} = 10^{-3}$ dont le résidu stagne dès le départ, et $d = 10$.
La non-convergence est détectée à l'itération 11 ($e_F > 0{,}99\,e_F^{(1)}$).
Le seuil $\epsilon_{nc}$ double ensuite à chaque itération à partir de la seizième et dépasse $10^{-3}$ à l'itération 22 ($10^{-5}\times 2^{7} = 1{,}28\cdot10^{-3}$) : c'est là que la non-convergence est déclarée, et que commence le sous-pas suivant.
\end{remarque}

\begin{remarque}[MAXITERATION n'est pas un plafond]
L'atteindre déclenche la détection, qui triple la limite.
Un résidu qui baisse lentement mais régulièrement (plus de 1\,\% toutes les $d$ itérations) n'est jamais déclaré non convergé avant $3 \times 250 = 750$ itérations.
La limite modifiée n'est pas réinitialisée au sous-pas suivant.
\end{remarque}

\begin{algobox}{6}{Tests après l'itération $i$}{lignes d'UNPAS}
\begin{algorithmic}[1]
\State $e_F \gets \max(\lVert \vc r\rVert_{\infty,F},\ \norminf{\vc A^{\mathsf T}(\vg\lambda^{(i)} - \vg\lambda^{(i-1)})}) / F_{ref}$ \Comment{\lig{3229-3242}}
\State $e_M \gets \lVert \vc r\rVert_{\infty,M} / M_{ref}$, si coques, poutres, joints ou poreux \Comment{\lig{3245-3258}}
\State aucun moment appliqué : $e_M \gets 0$ ; forces appliquées négligeables devant les moments : $e_F \gets 0$
\State $e_\eps \gets \norminf{\Deps^{(i)} - \Deps^{(i-1)}}$ \Comment{\lig{3321-3322}}
\Statex \textit{// convergence}
\If{le test n'est pas sauté} \Comment{sauté : $i = 1$ après prédicteur, frottement, souci}
  \State $p \gets \src{PRECISION}$ ; $p_\eps \gets \src{PRECDEFO}$ ; divisées par 10 en prédicteur HPP
  \If{$e_F < p$ et $e_\eps < p_\eps$ et ($e_M < \src{PRECFLEX}$ si moments)} \Comment{\lig{3330-3337}}
    \State \textbf{si} le contact n'est pas résolu : \Retourner{} échec \Comment{« contact »}
    \State \textbf{si} $\alpha < 1$ : \Retourner{} échec \Comment{« relaxation » : cible réduite atteinte}
    \State \textbf{si} augmentation ou réduction du chargement active : \Retourner{} échec \Comment{« augmentation »}
    \If{prédicteur HPP, souci de \op{COMP}, ou réactions de contact changées sur GEOM2}
      \State basculer en GD ou reconstruire $\vc K$ ; \Continuer \Comment{\lig{3401-3412}}
    \EndIf
    \State \Retourner{} convergé
  \EndIf
\EndIf
\Statex \textit{// détection de non-convergence}
\State $e_{ref} \gets 0{,}99\;e_F^{(i-d)}$ si $i > d$, sinon $\infty$ \Comment{$d = \src{DELTAITER}$ (10), \lig{3429-3432}}
\If{test non sauté, non détectée, $i \ge 2$, pas de pilotage, contact résolu et ($i \ge I_{max}$ ou $e_F > e_{ref}$)} \Comment{\lig{3437-3488}}
  \State NONCONV $\gets$ vrai ; $I_{max} \gets 3\,i$
  \State $\alpha \gets \src{RELAXATION_NONCONV}$ ; si $\alpha < 1$ : cible $\gets t_{déb} + \alpha\,\Delta t$ ; chargement et matériau réévalués
\EndIf
\State \textbf{si} NONCONV et $i > 15$ : $\epsilon_{nc} \gets 2\,\epsilon_{nc}$ \Comment{$\epsilon_{nc}$ part de $10^{-5}$}
\Statex \textit{// déclaration}
\If{test non sauté et $\norminf{\Deps} < \epsilon_{nc}$ et ($i \ge I_{max}$ ou (NONCONV et $e_F > e_{ref}$ et $\norminf{\Deps} \ge \tfrac14\norminf{\Deps_{prec}}$))} \Comment{\lig{3498-3506}}
  \State \Retourner{} échec \Comment{« NON CONVERGENCE »}
\EndIf
\State \Continuer
\end{algorithmic}
\end{algobox}
% ============================================================================
\section{Options avancées}
\label{sec:annexes}
% ============================================================================

Les annexes qui suivent décrivent des options qui s'insèrent dans les algorithmes précédents.
Chacune a son glossaire.

% ----------------------------------------------------------------------------
\subsection{Annexe A — Non local : moyenne et Helmholtz}
% ----------------------------------------------------------------------------

\begin{glossaire}
— & \src{NON_LOCAL} & Type de régularisation : \src{MOYE} (moyenne pondérée), \src{SB} (moyenne pondérée par l'état de contrainte), \src{HELM} (gradient implicite de type Helmholtz).\\
— & \src{CONN} & Champ de « connectivité non locale » : pour chaque point de Gauss, ses voisins et leurs poids.\\
$\tilde{\vc V}$ & \src{ZVARN} & Variables internes régularisées.\\
— & \src{ISTEP} & Rang de l'appel à \op{COMP} : 1 premier appel local, 3 appel intermédiaire avec $\tilde{\vc V}$, 2 appel final.\\
$j$ & \src{BIHL} & Boucle des sous-itérations non locales, au plus \src{MAXSOUSITERATION} (500).\\
$\epsilon_{nl}$ & \src{PRECSOUSITERATION} & Seuil des sous-itérations, $10^{-6}$.\\
$\vc H$ & \src{TAHELM} (\pr{PAS\_HELM}) & Matrices de Helmholtz, une par variable délocalisée : conduction et capacité d'un modèle « thermique » auxiliaire.\\
— & \src{ZNACCEHL}, \src{ITDEPHL} & Accélération \op{ACT3} des sous-itérations de Helmholtz : toutes les 3, après la 4\textsuperscript{e}.\\
— & \src{LOGHLIN} & Vrai si toutes les formulations de Helmholtz sont linéaires.\\
\end{glossaire}

Les lois adoucissantes perdent l'unicité de la solution et localisent sur une bande de largeur liée au maillage.
La régularisation non locale remplace, dans la loi, certaines variables internes par une version lissée sur une longueur interne.
\pr{UNPAS} la fait par alternance, à chaque sous-incrément : la loi locale, puis la régularisation, puis la loi relue avec la variable régularisée.

\paragraph{Moyenne (\src{MOYE}, \src{SB}).}
\op{NLOC} remplace $v$ par sa moyenne pondérée sur le voisinage décrit par \src{CONN}, de la forme classique
\begin{equation}
\tilde v(\vc x) = \frac{\int_\Omega w(\vc x, \vc y)\,v(\vc y)\,\mathrm d\vc y}{\int_\Omega w(\vc x, \vc y)\,\mathrm d\vc y}.
\end{equation}
En \src{SB}, le Fortran (\texttt{NLOCA1}) module les poids par les contraintes principales (version de Giry et Dufour) ; \pr{UNPAS} ajoute pour cela les contraintes principales aux champs moyennés (\lig{2560-2570}).
La moyenne est linéaire : un seul aller-retour suffit.

\paragraph{Helmholtz (\src{HELM}).}
Pour chaque variable délocalisée, on résout une équation de gradient implicite, $\tilde v - \ell^2\Delta\tilde v = v$, discrétisée comme un problème thermique stationnaire :
\begin{equation}
\vc H\,\tilde{\vc v} = \vc f,\qquad \vc H = \vc K_{\ell} + \vc C,\qquad \vc f = \op{SOUR}(v),
\end{equation}
où $\vc K_\ell$ joue le rôle d'une conduction ($\ell^2\int\nabla\vc N\nabla\vc N^{\mathsf T}$), $\vc C$ celui d'une capacité ($\int\vc N\vc N^{\mathsf T}$), et \pr{PAS\_HELM} les assemble avec le matériau courant (\lig{2588-2796}).
À la première sous-itération, $\tilde{\vc v} = \vc H^{-1}\vc f$ ; si la formulation est non linéaire, on itère sur le résidu $\vc r_H = \vc f - \vc H\tilde{\vc v}$, accéléré par \op{ACT3} aux sous-itérations 6, 9, 12…, et l'on mesure l'erreur $\lVert\vc C^{-1}\vc r_H\rVert_\infty / \lVert\vc C^{-1}\vc f\rVert_\infty$.

\paragraph{Arrêt.}
On passe à l'appel final (\src{ISTEP} = 2) à la dernière sous-itération permise, si tout est linéaire, ou si l'erreur maximale passe sous $\epsilon_{nl}$ après la première sous-itération (\lig{2810-2825}).
\op{KTAN} \src{PERT} est désactivé en non local (\lig{208}).

\begin{annexbox}{A}{Non local : moyenne et Helmholtz}{UNPAS \lig{2541-2876} · PAS\_HELM}
\begin{algorithmic}[1]
\Statex \textit{// après le premier \op{COMP} local (\src{ISTEP} = 1), dans \algoref{5}}
\For{$j = 1 \dots \src{MAXSOUSITERATION}$ (500)} \Comment{\src{BIHL}, \lig{2558}}
  \If{$\src{NON_LOCAL} = \src{MOYE}$ ou \src{SB}}
    \State $\tilde{\vc V} \gets \op{NLOC}(\vc V)$ : moyenne sur le voisinage \src{CONN} ; \src{SB} pondère par les contraintes \Comment{\lig{2560-2570}}
  \Else \Comment{\src{HELM}}
    \ForAll{variable délocalisée $v$} \Comment{\lig{2588-2796}}
      \State $\vc H \gets \vc K_\ell + \vc C$ sur le modèle de Helmholtz \Comment{\pr{PAS\_HELM}}
      \State $\vc f \gets \op{SOUR}(v$ aux nœuds$)$
      \State \textbf{si} $j = 1$ : $\tilde{\vc v} \gets \vc H^{-1}\vc f$
      \State \textbf{sinon si} formulation non linéaire : $\vc r_H \gets \vc f - \vc H\tilde{\vc v}$ ($-\op{ACT3}$ aux sous-itérations 6, 9, 12…) ; $\tilde{\vc v} \gets \tilde{\vc v} + \vc H^{-1}\vc r_H$
      \State $\mathrm{erreur}_v \gets \lVert\vc C^{-1}\vc r_H\rVert_\infty / \lVert\vc C^{-1}\vc f\rVert_\infty$ ; remplacer $v$ par $\tilde v$ dans $\vc V$
    \EndFor
  \EndIf
  \State $\src{ISTEP} \gets 2$ si $j = \src{MAXSOUSITERATION}$, si tout est linéaire, ou si $j > 1$ et $\max\mathrm{erreur}_v < \src{PRECSOUSITERATION}$ ; sinon 3 \Comment{\lig{2810-2825}}
  \State $(\sig, \vc V) \gets \op{COMP}$(état initial avec $\tilde{\vc V}$, état final, \src{ISTEP}) \Comment{\lig{2834-2852}}
  \State \textbf{si} $\src{ISTEP} = 2$ : \Sortir
\EndFor
\end{algorithmic}
\end{annexbox}

% ----------------------------------------------------------------------------
\subsection{Annexe B — MAN : méthode asymptotique numérique}
% ----------------------------------------------------------------------------

\begin{glossaire}
$N$ & \src{ORDRE} & Ordre de la série, 12 par défaut.\\
$\vc U_p$, $\lambda_p$, $\sig_p$ & \src{tadep.p}, \src{talam.p}, \src{tasig.p} & Termes d'ordre $p$ de la série en déplacement, en chargement et en contrainte.\\
$\vc K_t$ & — & \op{RIGI} + $\vc K_\sigma(\sig_0)$, plus $\tfrac4{\Delta t^2}\vc M$ en dynamique.\\
$\vc f$ & \src{ff} & Chargement restant à appliquer dans le segment ; au départ, le second membre de la première itération.\\
$a_{max}$ & \src{aserie} & Rayon de validité de la série.\\
$\delta$ & \src{prec} & Tolérance du rayon de validité, $10^{-8}$.\\
— & \src{ncontmax} & Nombre maximal de segments (continuations), 250.\\
$\Delta\vc U$ & \src{deltatot} & Déplacement cumulé sur les segments ; remplace $\delta\vc u$ à la première itération.\\
— & \src{IOUT} & 1 si la MAN a abouti, 0 si elle a renoncé (poutres, coques \src{COQ2}, trop de segments).\\
\end{glossaire}

Avec \src{MAN}, la correction de la première itération est remplacée par une prédiction asymptotique calculée par \pr{CORMAN} (\lig{1973-1980}).
La méthode porte sur l'élasticité géométriquement non linéaire (Green-Lagrange, \op{SIGMA} \src{'LINE'}) : elle prédit le chemin d'équilibre, pas la plasticité.
On cherche le chemin sous forme de séries en un paramètre $a$ :
\begin{equation}
\vc U(a) = \vc U_0 + \sum_{p=1}^{N} a^p\,\vc U_p,\qquad \lambda(a) = \sum_{p=1}^{N} a^p\,\lambda_p,\qquad \sig(a) = \sig_0 + \sum_{p=1}^{N} a^p\,\sig_p,
\end{equation}
avec l'équilibre $\vc F^{int}(\vc U(a)) = \lambda(a)\,\vc f$ et la définition de type longueur d'arc $a = (\vc U - \vc U_0)\cdot\vc U_1 + \lambda\,\lambda_1$.
À l'ordre 1, $\vc K_t\vc U_1 = \lambda_1\vc f$ et $\lVert\vc U_1\rVert^2 + \lambda_1^2 = 1$, d'où
\begin{equation}
\lambda_1 = \frac1{\sqrt{\lVert\vc K_t^{-1}\vc f\rVert^2 + 1}},\qquad \vc U_1 = \lambda_1\,\vc K_t^{-1}\vc f .
\end{equation}
Le source écrit \texttt{psc2 = psc + 1.D0 ** 0.5D0}, qui vaut bien $(\mathrm{psc} + 1)^{1/2}$ par l'évaluation de gauche à droite (§ \ref{sec:gibiane}).
Aux ordres $p \ge 2$, on résout $\vc K_t\vc U_p = \lambda_p\vc f + \vc F_p$, où $\vc F_p$ ne dépend que des ordres inférieurs,
\begin{equation}
\vc F_p = -\sum_{r=1}^{p-1}\vc K_\sigma(\sig_r)\,\vc U_{p-r} - \vc B^{\mathsf T}\vc D\,\tfrac12\sum_{r=1}^{p-1}\nabla\vc U_r^{\mathsf T}\nabla\vc U_{p-r},
\end{equation}
et la condition d'orthogonalité $\vc U_p\cdot\vc U_1 + \lambda_p\lambda_1 = 0$ donne
\begin{equation}
\lambda_p = -\lambda_1\,\big(\vc K_t^{-1}\vc F_p\big)\cdot\vc U_1,\qquad \vc U_p = \vc K_t^{-1}\vc F_p + \frac{\lambda_p}{\lambda_1}\,\vc U_1 .
\end{equation}
La série est valable tant que le dernier terme reste négligeable :
\begin{equation}
a_{max} = \left(\delta\,\frac{\lVert\vc U_1\rVert}{\lVert\vc U_N\rVert}\right)^{1/(N-1)} .
\end{equation}
Si $\lambda(a_{max}) > 1$, on cherche $a$ tel que $\lambda(a) = 1$ et l'on termine ; sinon on avance de $a_{max}$, on actualise contraintes et géométrie, et l'on recommence un segment avec le chargement restant $\vc f\,(1 - \lambda(a))$.
Les approximants de Padé sont présents dans le code mais commentés.

\begin{annexbox}{B}{MAN : méthode asymptotique numérique}{CORMAN · UNPAS \lig{1973-1980}}
\begin{algorithmic}[1]
\Statex \textit{// à la première itération de \algoref{4}, si \src{MAN} ; procédure \pr{CORMAN}}
\State \textbf{si} le modèle contient \src{POUT} ou \src{COQ2} : abandon, on garde la prédiction de Newton
\State $\Delta\vc U \gets 0$
\ForAll{segment, 250 au plus} \Comment{continuation}
  \State $\vc K_t \gets \op{RIGI} + \vc K_\sigma(\sig_0)$, $+\,(4/\Delta t^2)\,\vc M$ en dynamique
  \State $\vc U_1 \gets \vc K_t^{-1}\vc f / \sqrt{\lVert\vc K_t^{-1}\vc f\rVert^2 + 1}$ ; $\lambda_1 \gets 1/\sqrt{\lVert\vc K_t^{-1}\vc f\rVert^2 + 1}$
  \For{$p = 2 \dots N$ (12)}
    \State $\vc F_p \gets -\sum\vc K_\sigma(\sig_r)\,\vc U_{p-r} - \vc B^{\mathsf T}\vc D\,\tfrac12\sum\nabla\vc U_r^{\mathsf T}\nabla\vc U_{p-r}$
    \State $\lambda_p \gets -\lambda_1\,(\vc K_t^{-1}\vc F_p\cdot\vc U_1)$ ; $\vc U_p \gets \vc K_t^{-1}\vc F_p + (\lambda_p/\lambda_1)\,\vc U_1$
  \EndFor
  \State $a_{max} \gets$ rayon de validité de la série
  \State \textbf{si} $\lambda(a_{max}) > 1$ : $a \gets$ racine de $\lambda(a) = 1$ sur $[0, a_{max}]$ ; fin $\gets$ vrai ; \textbf{sinon} $a \gets a_{max}$
  \State $\Delta\vc U \gets \Delta\vc U + \sum_p a^p\,\vc U_p$ ; contraintes et géométrie actualisées
  \State \textbf{si} fin : \Sortir
  \State $\vc f \gets \vc f\,(1 - \lambda(a))$
\EndFor
\State \Retourner{} $\Delta\vc U$ \Comment{remplace $\delta\vc u$ à $i = 1$}
\end{algorithmic}
\end{annexbox}

% ----------------------------------------------------------------------------
\subsection{Annexe C — Pilotage automatique et indirect}
% ----------------------------------------------------------------------------

\begin{glossaire}
$\alpha$ & \src{COEPI}, \src{AUTOCOEF} & Facteur de chargement : fraction de l'incrément de chargement appliquée ; \src{AUTOCOEF} le transmet à \pr{PASAPAS}.\\
$\epsilon_a$ & \src{AUTOCRIT} & Incrément de déformation visé, $10^{-3}$.\\
— & \src{AUTORECA} & Période (en itérations) de la renormalisation, 1.\\
— & \src{AUTOPAS} & Nombre maximal de reprises d'une même cible, 250.\\
— & \src{ZAUTOREDU}, \src{RED2} & Diviseur du critère, triplé quand \src{RED2} $< i/20$.\\
$n$ & \src{U1MA} & Norme de déformation de l'incrément, calculée par \pr{AUTOPILO}.\\
$\beta$ & \src{AL1} & Coefficient de normalisation de l'incrément.\\
$\vc R^{nom}$ & \src{RESIDNOR} & Incrément nominal de chargement.\\
$\vc F_p$ & \src{ZFPILIN} & Forces et déplacements pilotés (\src{FORCES_PILOTEES}, \src{DEPLACEMENTS_PILOTES}).\\
$\delta\vc u_{II}$ & \src{ZDEPII} & Réponse à $\vc F_p$ : $\vc K^{-1}\vc F_p$.\\
$\eta$, $\Delta\eta$ & \src{ETA}, \src{D_ETA} & Multiplicateur du chargement piloté et son incrément, calculé par \pr{PILOINDI}.\\
$\Delta\tau$ & \src{DTAU} & Incrément du paramètre de pilotage (\src{PARAMETRE_DE_PILOTAGE}).\\
\end{glossaire}

\paragraph{Pilotage automatique.}
Avec \src{AUTOMATIQUE}, le pas n'est plus une durée imposée : on limite l'incrément de déformation, et l'on choisit le chargement qui le rend compatible avec l'équilibre.
À chaque itération multiple de \src{AUTORECA}, on renormalise l'incrément (\lig{2094-2124}) :
\begin{equation}
n = \max\Big(\norminf{\vg\eps(\Delta\vc u)},\ \norminf{\vg\eps(\Delta\vc u - \alpha\,\Delta\vc u^{él})}\Big),\qquad
\beta = \min\Big(\frac{\epsilon_a/\mathrm{ZAUTOREDU}}{n},\ \frac1\alpha\Big),\qquad \Delta\vc u \gets \beta\,\Delta\vc u,
\end{equation}
où $\Delta\vc u^{él}$ est la part élastique (\pr{AUTOPILO}) ; la borne $1/\alpha$ ne s'applique que si $\beta > 1$.
Le diviseur \src{ZAUTOREDU} est multiplié par 3 chaque fois que \src{RED2} $< i/20$ (\lig{2100-2103}) ; on ne sait pas, sans l'avoir testé, si Gibiane fait ici une division entière.
Après le calcul des forces internes, on actualise le facteur de chargement (\lig{3155-3205}) :
\begin{equation}
\alpha \gets
\begin{cases}
\min(\alpha\beta, 1) & \text{loi visqueuse, thermique ou \src{DEFI}},\\[2pt]
\min\Big(\alpha - \dfrac{\vc r\cdot\Delta\vc u}{\vc R^{nom}\cdot\Delta\vc u},\ 1\Big) & \text{sinon,}
\end{cases}
\end{equation}
ce qui annule le travail du résidu dans la direction de l'incrément.
Le chargement appliqué est $\vc F_n + \alpha\,\Delta\vc F$, et les déplacements imposés $\alpha\,\Delta\vc u^{imp}$.
Il n'y a pas de détection de non-convergence (\algoref{6}).
\pr{PASAPAS} avance le temps de $\alpha\,(T_I - t)$ et rejoue la cible tant que $\alpha$ n'atteint pas 1 ; à la reprise, $\alpha$ repart de $\min(\alpha_{prec}/(1 - \alpha_{prec}), 1)$ (\lig{970-979}).

\paragraph{Pilotage indirect (\src{CHAR_PILO}).}
Le chargement s'écrit $\vc F = \vc F^{ext} + \eta\,\vc F_p$, et l'on choisit $\eta$ pour qu'une mesure de déformation dans une zone donnée progresse de $\Delta\tau$.
À la première résolution, on calcule $\delta\vc u_{II} = \vc K^{-1}\vc F_p$ (\lig{1303-1311}).
Après chaque résolution, \pr{PILOINDI} (\lig{2031-2037}) calcule, en chaque point de Gauss $g$ de \src{ZONE_DE_PILOTAGE}, la direction $\vc n_g = \vg\eps(\vc u_n)/\lVert\vg\eps(\vc u_n)\rVert$ (norme en énergie) et
\begin{equation}
a_0 = \vc n_g : \vg\eps(\Delta\vc u + \delta\vc u),\qquad a_1 = \vc n_g : \vg\eps(\delta\vc u_{II}),\qquad
\Delta\eta = \min_{g\,:\,a_1 > 0}\frac{\Delta\tau - a_0}{a_1},
\end{equation}
puis $\delta\vc u \gets \delta\vc u + \Delta\eta\,\delta\vc u_{II}$ et $\eta \gets \eta + \Delta\eta$.

\begin{annexbox}{C}{Pilotage automatique et indirect}{AUTOPILO · PILOINDI · UNPAS}
\begin{algorithmic}[1]
\Statex \textit{// pilotage automatique (\src{AUTOMATIQUE}), dans \algoref{4}}
\State $\alpha \gets 1$, ou $\min(\alpha_{prec}/(1 - \alpha_{prec}), 1)$ si l'on rejoue la même cible \Comment{\lig{970-979}}
\ForAll{itération multiple de \src{AUTORECA} (1)} \Comment{\lig{2094-2124}}
  \State \textbf{si} $\src{RED2} < i/20$ : $\src{RED2} \gets \src{RED2} + 1$ ; critère $\gets$ critère $/\,3$ \Comment{\lig{2100-2103}}
  \State $n \gets \max(\norminf{\vg\eps(\Delta\vc u)}, \norminf{\vg\eps(\Delta\vc u - \alpha\,\Delta\vc u^{él})})$ \Comment{\pr{AUTOPILO}}
  \State $\beta \gets \src{AUTOCRIT}/n$, au plus $1/\alpha$ ; $\Delta\vc u \gets \beta\,\Delta\vc u$ \Comment{$\src{AUTOCRIT} = 10^{-3}$}
  \State après $\vc F^{int}$ : \textbf{si} visqueux, thermique ou \src{DEFI} : $\alpha \gets \min(\alpha\beta, 1)$ ; \textbf{sinon} $\alpha \gets \min(\alpha - (\vc r\cdot\Delta\vc u)/(\vc R^{nom}\cdot\Delta\vc u), 1)$ \Comment{\lig{3155-3205}}
\EndFor
\State chargement $\gets \vc F_n + \alpha\,\Delta\vc F$ ; $\vc u^{imp} \gets \alpha\,\Delta\vc u^{imp}$ ; pas de détection de non-convergence
\Statex \textit{// pilotage indirect (\src{CHAR_PILO}) : $\vc F = \vc F^{ext} + \eta\,\vc F_p$}
\State à la première résolution : $\delta\vc u_{II} \gets \vc K^{-1}\vc F_p$ \Comment{\lig{1303-1311}}
\ForAll{résolution} \Comment{\pr{PILOINDI}, \lig{2031-2037}}
  \ForAll{point de Gauss $g$ de \src{ZONE_DE_PILOTAGE}}
    \State $\vc n_g \gets \vg\eps(\vc u_n)/\lVert\vg\eps(\vc u_n)\rVert$ en énergie ; $a_0 \gets \vc n_g : \vg\eps(\Delta\vc u + \delta\vc u)$ ; $a_1 \gets \vc n_g : \vg\eps(\delta\vc u_{II})$
  \EndFor
  \State $\Delta\eta \gets \min$, sur les $g$ où $a_1 > 0$, de $(\Delta\tau - a_0)/a_1$
  \State $\delta\vc u \gets \delta\vc u + \Delta\eta\,\delta\vc u_{II}$ ; $\eta \gets \eta + \Delta\eta$
\EndFor
\end{algorithmic}
\end{annexbox}

% ----------------------------------------------------------------------------
\subsection{Annexe D — Contact unilatéral}
% ----------------------------------------------------------------------------

\begin{glossaire}
— & \src{RFCO} & Opérateur qui calcule les relations de contact (et de frottement) et les jeux, sur la configuration courante.\\
— & \src{MODCONTA} & Modèle de contact.\\
— & \src{SUP} & Super-élément : la rigidité condensée sur les degrés de liberté de contact.\\
— & \src{MATAB} & Table gardée d'un appel de \pr{UNILATER} à l'autre : super-élément, statut \src{'OK'}, ensemble actif (\src{MATAB.6}).\\
— & \src{INIB} & Mot-clé : partir des statuts du dernier appel.\\
— & \src{LISEN}, \src{LITOT} & Ensemble actif courant ; historique des ensembles visités, pour détecter les cycles.\\
— & \src{OPTION} & Stratégie d'activation : \src{RAPIDE}, \src{MOYEN}, \src{LENT} ; \src{MOSTAB} \src{'STAB'} pour l'opérateur stabilisé.\\
— & \src{PASUNIL} & Côté \pr{UNPAS} : vrai si \pr{UNILATER} a échoué (\src{'OK'} faux).\\
— & \src{CRR}, \src{CRRN}, \src{RAP4} & Côté \pr{UNPAS} : relations utilisées, relations recalculées sur GEOM2, variation relative des réactions.\\
— & \src{CAFROTTE} & Vrai s'il y a du frottement.\\
\end{glossaire}

Le contact se traite par des relations unilatérales $\vc A_c\vc u \ge \vc g$ (jeux $\vc g$), dont les multiplicateurs doivent être de signe convenable et nuls sur les relations inactives.
\op{RESO} passe la main à \pr{UNILATER} dès que la rigidité contient de telles relations (Fortran \texttt{RESOU}).

\paragraph{La condensation.}
Au premier appel, \pr{UNILATER} condense la rigidité sur les degrés de liberté de contact (super-élément, \pr{UNILATER} \lig{68-135}) :
\begin{equation}
\vc K_c^{*} = \vc K_{cc} - \vc K_{ci}\,\vc K_{ii}^{-1}\,\vc K_{ic},
\end{equation}
ce qui rend chaque essai d'ensemble actif peu coûteux.

\paragraph{L'ensemble actif.}
On part des statuts du dernier appel (\src{INIB}), puis on répète : résoudre avec les relations actives ; \op{EXCITER} active les relations violées et relâche celles dont le multiplicateur a le mauvais signe (avec frottement, il utilise les forces limites) ; s'arrêter quand l'ensemble ne change plus (boucle \src{BOUCL1}, \lig{184-273}).
Contre les cycles, la stratégie passe de \src{RAPIDE} à \src{MOYEN} à la 5\textsuperscript{e} itération, à \src{LENT} quand un ensemble déjà visité revient (ou toutes les 100 itérations), et en dernier recours à une stratégie alternative avec un opérateur stabilisé.
En cas d'échec, \src{'OK'} est faux, et \pr{UNPAS} conclut à une non-convergence de contact (\pr{UNPAS} \lig{3312-3315}).

\paragraph{Dans UNPAS.}
Les relations \op{RFCO} sont recalculées à chaque reconstruction de $\vc K$ en GD.
Au moment de conclure, on les recalcule sur GEOM2 ; si les réactions changent de plus de \src{PRECISION} en relatif, on refuse la convergence et l'on continue avec les nouvelles relations (\lig{3364-3398}).
Avec frottement, il faut au moins 2 itérations.

\begin{annexbox}{D}{Contact unilatéral}{UNILATER · RFCO · EXCITER}
\begin{algorithmic}[1]
\Statex \textit{// \op{RESO} appelle \pr{UNILATER} dès que $\vc K$ contient des conditions unilatérales}
\State condenser $\vc K$ sur les ddl de contact (super-élément), au premier appel \Comment{\pr{UNILATER} \lig{68-135}}
\State actives $\gets$ statuts du dernier appel (\src{INIB}), sinon initialisation
\Loop \Comment{\src{BOUCL1}, \lig{184-273}}
  \State résoudre avec les conditions actives
  \State \op{EXCITER} : activer les conditions violées, relâcher celles de mauvais signe ; frottement : forces limites
  \State \textbf{si} les conditions actives n'ont pas changé : \Sortir
  \State anti-cyclage : \src{RAPIDE}, puis \src{MOYEN} à la 5\textsuperscript{e} itération, puis \src{LENT}, puis opérateur stabilisé
\EndLoop
\State \src{OK} $\gets$ faux en cas d'échec $\to$ \pr{UNPAS} : non-convergence de contact \Comment{\pr{UNPAS} \lig{3312-3315}}
\end{algorithmic}
\end{annexbox}

% ----------------------------------------------------------------------------
\subsection{Annexe E — Rayonnement thermique}
% ----------------------------------------------------------------------------

\begin{glossaire}
— & \src{IAPPEL} & Contexte de l'appel à \pr{PAS\_RAYO} : 1 \pr{DUPONT2} ; 2 \pr{TRANSNON}, avant la boucle en temps ; 3 \pr{TRANSNON}, matrices.\\
$\sigma_{SB}$ & \src{CTE_SB} & \src{CTE_STEFAN_BOLTZMANN}, $5{,}673\cdot10^{-8}$ (SI) par défaut.\\
— & \src{PRRA}, \src{WTAB.'RAYONNEMENT'} & Table des conditions de rayonnement, construite par l'opérateur \op{PRRA} à partir des sous-modèles \src{RAYONNEMENT}.\\
$\vc F_{vue}$ & \src{FAC_FORM} & Facteurs de forme de chaque cavité (\op{FFOR}).\\
— & \src{REA_GEOM} & Recalculer les facteurs de forme.\\
$T$ & \src{U_CAK2} & Température courante en kelvins (\src{THER_COURANT} + \src{TREF}).\\
$\varepsilon$ & \src{CH_EMI} & Émissivité \src{EMIS} du matériau.\\
$T_e$, $\varepsilon_e$ & \src{TERA} ou \src{T_IN}, \src{E_IN} & Température et émissivité « à l'infini ».\\
$T_{rad}$ & \src{TCAV} & Température de rayonnement équivalente vue par chaque face d'une cavité (\op{RAYE}).\\
$h$ & \src{HRAD}, \src{MAT_RAD} & Coefficient d'échange linéarisé.\\
$\vc K_r$, $\vc F_r$ & \src{RIG_RAD}, \src{ADDI_MATRICE}, \src{ADDI_SECOND} & Matrice et second membre rendus aux procédures thermiques.\\
$\vc R_{ray}$ & \src{MAT_RAYE} & Matrice de rayonnement de la méthode \src{TRAYO} : $\Phi = \vc R_{ray}T^4$.\\
\end{glossaire}

\paragraph{La linéarisation.}
\pr{PAS\_RAYO} traite le rayonnement comme une convection dont le coefficient dépend de la température.
Pour un échange entre une paroi à $T$ et un milieu à $T_e$ (\pr{HRAYO}) :
\begin{gather}
h = \sigma_{SB}\,\varepsilon_{éq}\,\big(T^2 + T_e^2\big)\big(T + T_e\big)
\quad\Longrightarrow\quad
h\,(T - T_e) = \sigma_{SB}\,\varepsilon_{éq}\,\big(T^4 - T_e^4\big),\\
\varepsilon_{éq} = \frac{\varepsilon_1\varepsilon_2}{1 - (1-\varepsilon_1)(1-\varepsilon_2)} .
\end{gather}
La linéarisation est sécante : à température donnée, $h(T - T_e)$ redonne exactement le flux rayonné.
Vers l'infini avec $\varepsilon_e = 1$, $\varepsilon_{éq}$ se réduit à l'émissivité de la paroi.
Les températures sont converties en kelvins ($T_{ref} = 273{,}15$ avec \src{CELSIUS}).

\paragraph{Les trois types de conditions.}
\begin{itemize}
\item \textbf{\src{INFINI}} : échange avec un milieu à $T_e$ (chargement \src{TERA}, sinon \src{T_IN} du matériau) ; $\vc K_r = h\int\vc N\vc N^{\mathsf T}$ (\op{CONDUCTIVITE}) et $\vc F_r = h\int\vc N\,T_e$ (\op{CONVECTION}) (\pr{PAS\_RAYO} \lig{134-161}).
\item \textbf{\src{FAC_A_FAC}} : échange direct entre deux faces en vis-à-vis, dont on projette la température et l'émissivité (\op{KPRO}) ; $\vc K_r$ couple les deux faces, il n'y a pas de second membre (\lig{117-131}).
\item \textbf{\src{CAVITE}} : échange entre toutes les faces d'une cavité, pondéré par les facteurs de forme. Sans l'option \src{TRAYO}, \op{RAYE} résout les radiosités $\vc J$,
\begin{equation}
\big(\vc I - (1 - \varepsilon)\,\vc F_{vue}\big)\,\vc J = \varepsilon\,\sigma\,T^4,
\end{equation}
par BiCGSTAB à $10^{-7}$ près (Fortran \texttt{RAYT1}), et en déduit pour chaque face la température équivalente $T_{rad}$ telle que $\Phi = \varepsilon\,\sigma\,(T^4 - T_{rad}^4)$ ; \pr{HRCAV} donne alors $h$ avec $T_e = T_{rad}$ (\lig{173-185}). Avec \src{TRAYO}, \op{RAYE} rend la matrice $\vc R_{ray}$ telle que $\Phi = \vc R_{ray}T^4$, et \op{RAYN} la matrice sécante $\vc K_r \propto \sigma\,\vc R_{ray}T^3$, sans second membre (\lig{186-191}).
\end{itemize}

\paragraph{L'usage par les procédures thermiques.}
\pr{TRANSNON} appelle \pr{PAS\_RAYO} avant la boucle en temps pour les facteurs de forme (appel 2, \pr{TRANSNON} \lig{456-462}), puis avant les itérations à $\vc T_0$ : $\theta\vc K_r$ entre dans $\vc S$ et $\vc F_r$ dans le second membre (appel 3, \lig{966-982}), et à chaque itération à $\vc T_\theta$ pour le résidu (\lig{1424-1437}).
La matrice garde $h(\vc T_0)$, mais le résidu utilise $h(\vc T_\theta)$ : à convergence, le flux rayonné vaut $\sigma\varepsilon(T^4 - T_e^4)$ à la précision du critère près.
Comme $h$ varie en $T^3$, le point fixe converge d'autant plus lentement que la température change dans le pas.
\pr{DUPONT2} l'appelle une fois par sous-pas, à $\vc T^{*}$ (appel 1, \pr{DUPONT2} \lig{381-391}) ; faute d'itération, $h$ n'est jamais réévalué à la température obtenue.
\pr{TRANSLIN} ne l'appelle pas.

\paragraph{Deux pièges.}
Pour les cavités sans \src{TRAYO}, \pr{PAS\_RAYO} appelle \pr{HRCAV} sans lui passer $\sigma_{SB}$ (\lig{178}) : \pr{HRCAV} prend sa valeur SI, $5{,}67\cdot10^{-8}$, si bien qu'un calcul dans d'autres unités obtient un $h$ faux pour ces cavités (\texttt{RAYT1} code aussi cette valeur en dur, mais $T_{rad}$ n'en dépend pas).
Et les facteurs de forme ne sont recalculés qu'au premier appel ou si \src{REA_GEOM} est vrai, ce qui n'arrive qu'aux passes $\ge 2$ de \src{CONVERGENCE_MEC_THE} en GD : avec une seule passe, une cavité qui se déforme garde ses facteurs de forme initiaux.

\begin{annexbox}{E}{Rayonnement thermique}{PAS\_RAYO 25/05/21 · HRAYO · HRCAV}
\begin{algorithmic}[1]
\Statex \textit{// \pr{PAS\_RAYO}$(t, \text{appel})$ : appel 1 = \pr{DUPONT2} ; 2 = \pr{TRANSNON} avant la boucle en temps ; 3 = \pr{TRANSNON}, matrices}
\State \textbf{si} le modèle n'a pas de rayonnement : \Retourner{} rien \Comment{\lig{23-26}}
\If{appel $\le 2$} \Comment{facteurs de forme, \lig{57-86}}
  \State $\sigma_{SB} \gets \src{CTE_STEFAN_BOLTZMANN}$, $5{,}673\cdot10^{-8}$ par défaut (unités SI)
  \State \textbf{si} \src{REA_GEOM} ou premier appel : pour chaque zone \src{CAVITE}, $\vc F_{vue} \gets \op{FFOR}(\text{zone})$
\EndIf
\If{appel = 1 ou 3} \Comment{matrices et seconds membres, \lig{91-211}}
  \State $T \gets \src{THER_COURANT} + T_{ref}$ \Comment{kelvins}
  \ForAll{condition de la table \op{PRRA}} \Comment{\lig{98-202}}
    \State $\varepsilon \gets$ émissivité \src{EMIS} du matériau courant
    \If{\src{INFINI}} \Comment{\lig{134-161}}
      \State $T_e \gets \src{TERA}(t)$, sinon \src{T_IN} ; $\varepsilon_e \gets \src{E_IN}$ (1)
      \State $h \gets \sigma_{SB}\,\varepsilon_{éq}\,(T^2 + T_e^2)(T + T_e)$ ; $\vc K_r \gets h\int\vc N\vc N^{\mathsf T}$ ; $\vc F_r \gets h\int\vc N\,T_e$ \Comment{\pr{HRAYO}}
    \ElsIf{\src{FAC_A_FAC}} \Comment{\lig{117-131}}
      \State $T_2, \varepsilon_2 \gets$ face en vis-à-vis (\op{KPRO}) ; $h \gets \sigma_{SB}\,\varepsilon_{éq}\,(T^2 + T_2^2)(T + T_2)$ ; $\vc K_r \gets$ échange ; pas de $\vc F_r$
    \ElsIf{\src{CAVITE} sans \src{TRAYO}} \Comment{\lig{173-185}}
      \State $T_{rad} \gets \op{RAYE}(\vc F_{vue}, \varepsilon, T)$ : radiosité par BiCGSTAB à $10^{-7}$ près
      \State $h \gets 5{,}67\cdot10^{-8}\,\varepsilon\,(T^2 + T_{rad}^2)(T + T_{rad})$ \Comment{\pr{HRCAV} : $\sigma_{SB}$ n'est pas transmise}
      \State $\vc K_r \gets h\int\vc N\vc N^{\mathsf T}$ ; $\vc F_r \gets h\int\vc N\,T_{rad}$
    \Else \Comment{\src{CAVITE} avec \src{TRAYO}, \lig{186-191}}
      \State $\vc R_{ray} \gets \op{RAYE}(\vc F_{vue}, \varepsilon)$ ; $\vc K_r \gets \op{RAYN}(\vc R_{ray}, T, \sigma_{SB}) \propto \sigma_{SB}\vc R_{ray}T^3$ ; pas de $\vc F_r$
    \EndIf
  \EndFor
  \State \Retourner{} $\vc K = \sum\vc K_r$ (\src{ADDI_MATRICE}) et $\vc F = \sum\vc F_r$ (\src{ADDI_SECOND}) \Comment{\lig{196-209}}
\EndIf
\end{algorithmic}
\end{annexbox}

% ----------------------------------------------------------------------------
\subsection{Rigidité augmentée}
\label{sec:augmentation}
% ----------------------------------------------------------------------------

\begin{glossaire}
$\vc M_a$ & \src{RIG_AUG} & Matrice d'augmentation : celle de l'utilisateur (\src{RIGIDITE_AUGMENTEE}), ou la matrice de masse en mode automatique.\\
— & \src{AUTAUG}, \src{IRAUGLU}, \src{IRAUG} & Mode automatique (\src{AUGMENTATION_AUTOMATIQUE}) ; matrice fournie ; augmentation active.\\
$c$ & \src{augmult} & Coefficient adaptatif : 0,3 au départ, borné à $[10^{-3}, 10^{3}]$.\\
$a_M$, $a_K$ & \src{augm}, \src{augk} & Échelles de la masse et de la rigidité ajoutées.\\
— & \src{resmul} & Facteur de réduction du chargement appliqué.\\
— & \src{RED_URG}, \src{ref} & Nombre d'annulations d'itération ; seuil d'annulation ($1/\mathrm{ref}$).\\
— & \src{LOG_CNTRL}, \src{DE_CNTRL} & Contrôle par relations de déformation (\pr{DEFO\_IMP}) : code mort, \src{LOG_CNTRL} valant faux en dur (\lig{26}).\\
\end{glossaire}

C'est une stabilisation de type relaxation dynamique : on raidit l'opérateur pour que les corrections restent petites, puis on retire la raideur ajoutée.

\paragraph{Matrice de l'utilisateur.}
Sans mode automatique, la matrice \src{RIGIDITE_AUGMENTEE} est ajoutée telle quelle à $\vc K$, sans mise à l'échelle, et la rigidité est reconstruite à chaque pas (\lig{111-116}, \lig{743-746}).
Rien ne la retire : elle modifie l'opérateur, pas l'équilibre, et ne change donc que la vitesse de convergence.

\paragraph{Mode automatique.}
On ajoute à l'opérateur
\begin{equation}
c\,\big(a_M\,\vc M + a_K\,\vc K_e\big),
\end{equation}
où $\vc M$ est la matrice de masse (ou celle de l'utilisateur).
En GD, à chaque reconstruction de $\vc K$, $a_M = 10\,\vc x^{\mathsf T}\vc K_e\vc x / \vc x^{\mathsf T}\vc M\vc x$ et $a_K = 0{,}1$, $\vc x$ étant la dernière correction ; sinon $a_M = a_K = 1$ (\lig{606-608}, \lig{1745-1762}).
Le coefficient $c$ s'adapte : en début de sous-pas, il est remis à 0,3 si le précédent a convergé sans difficulté, divisé par 1,8 sinon, et multiplié par 1,5 s'il y a eu des annulations ; il est multiplié par 1,02 à chaque reconstruction et par 0,55 en fin de pas (\lig{650-680}, \lig{1759}, \lig{3604}).
Une convergence obtenue avec l'augmentation active est refusée si $c > 0{,}06$ ou si $\norminf{\Deps} > \src{PRECDEFO}$ : $c$ est divisé par 10 et un nouveau sous-pas repart (\lig{3353-3360}).
En GD, une itération dont $\norminf{\Deps}$ dépasse $1/\mathrm{ref}$ ($\mathrm{ref} = 20$ ; 50 si $\Deps$ a triplé ; 10 si la charge est presque annulée) est annulée, et $c$ est multiplié par 1,5. À partir de la quatrième annulation dans le sous-pas (le compteur est testé avant d'être incrémenté), la charge appliquée est réduite (\src{resmul}) et la non-convergence détectée, avec $I_{max} \gets i + 5$ ; à partir de la cinquième, $I_{max} \gets i - 1$ (\lig{2362-2409}).
Le facteur \src{resmul} réduit le résidu appliqué (« Réduction du chargement ») ; tant qu'il est inférieur à 0,99, la convergence est refusée.
% ============================================================================
\section{Valeurs par défaut utiles}
\label{sec:defauts}
% ============================================================================

Les valeurs sont celles de \pr{PAS\_DEFA} ; la dernière colonne signale les écrasements qui ne se voient pas dans la documentation.

\begin{center}
\small
\renewcommand{\arraystretch}{1.15}
\begin{longtable}{@{}>{\raggedright\arraybackslash}p{0.3\linewidth}>{\raggedright\arraybackslash}p{0.14\linewidth}>{\raggedright\arraybackslash}p{0.5\linewidth}@{}}
\toprule
Indice & Défaut & Remarque\\
\midrule\endhead
\bottomrule\endfoot
\src{PRECISION} / \src{PRECDEFO} / \src{PRECFLEX} & $10^{-4}$ & \src{PRECFLEX} prend la valeur de \src{PRECISION} si celle-ci est donnée. \src{PRECISION} sert aussi à la thermique et à \src{CRITERE_COHERENCE}.\\
\src{MAXITERATION} & 49 $\to$ 250 & Mis à 250 avec la convergence forcée ou \src{RENORMALISATION}, \emph{même si l'utilisateur l'a donné} (\lig{1001-1011}) ; 250 aussi en pilotage ou avec \src{PAS_AJUSTE} s'il n'est pas donné. Ce n'est pas un plafond strict (§ \ref{sec:criteres}).\\
\src{DELTAITER} / \src{MAXSOUSPAS} & 10 / 100 & \src{DELTAITER} vaut \src{MAXITERATION} en pilotage.\\
\src{CONVERGENCE_FORCEE} / \src{RELAXATION_NONCONV} & VRAI / 1 & La convergence forcée n'est fausse que sur demande ou avec \src{CONVERGENCE_MEC_THE}.\\
\src{K_TANGENT} / \src{NITER_KTANGENT} & FAUX / 1 & Rigidité élastique par défaut. \src{K_TANG_PERT_C1} $= 10^{-3}$, \src{C2} $= C_1/100$.\\
\src{K_SIGMA} & FAUX & Vrai en GD avec \src{QUADRATIQUE} (défaut des GD) ou si \src{STABILITE} est faux (\lig{1428-1439}).\\
\src{HYPOTHESE_DEFORMATIONS} & \src{LINEAIRE} & \src{QUADRATIQUE} en GD ; aussi \src{TRUESDELL}, \src{JAUMANN}, \src{UTILISATEUR}.\\
\src{LAGRANGIEN} & \src{MI_PAS} & \src{LAG_TOT} = 3 en GD ; aussi \src{TOTAL}, \src{FIN_PAS} ; 0 en endommagement.\\
\src{SOUS_INCREMENT} / \src{DEPS_MAX} & 1 / $5\cdot10^{-3}$ & Recalcul automatique en GD (20 au plus), sauf si \src{SOUS_INCREMENT} est donné.\\
\src{INITIALISATION} & VRAI & Prédicteur $\vc F^{NL}$.\\
\src{STABILITE} & VRAI & Faux dès que l'indice \src{AUTOMATIQUE} est fourni, même à faux (\lig{738}) ; revient à vrai avec \src{AUGMENTATION_AUTOMATIQUE} (\lig{1549}) ; l'indice \src{STABILITE} donné prime (\lig{971}).\\
\src{PRECISINTER} / \src{NITERINTER_MAX} / \src{NMAXSUBSTEPS} & $10^{-8}$ / 15 / 1 & Intégration locale (\op{KTAN}, \op{SSTE}, \op{ECFE}) ; $5\cdot10^{-7}$ en visqueux, \src{PRECISION} en endommagement.\\
\src{NB_BOTH} / \src{RELAXATION_THETA} & 1 / 1 & \src{NB_BOTH} passe à 10 avec \src{CONVERGENCE_MEC_THE}.\\
\src{PROCEDURE_THERMIQUE} / \src{RELAXATION_DUPONT} & \src{NONLINEAIRE} / 0,25 & \pr{TRANSNON} ; \src{LINEAIRE} appelle \pr{TRANSLIN}, \src{DUPONT} appelle \pr{DUPONT2}. Forcé à \src{NONLINEAIRE} avec changement de phase, solubilité, diffusion ou THM (\lig{1203-1206}). \src{SOUS-RELAXATION} et \src{SOUS_RELAXATION} ne sont jamais transmis.\\
\src{MAXSOUSITERATION} / \src{PRECSOUSITERATION} & 500 / $10^{-6}$ & Non local.\\
\src{ORDRE} · \src{AUTOCRIT} · \src{AUTOPAS} · \src{AUTORECA} & 12 · $10^{-3}$ · 250 · 1 & MAN et pilotage automatique.\\
\src{CTE_STEFAN_BOLTZMANN} & $5{,}673\cdot10^{-8}$ & Unités SI ; ignorée par \pr{HRCAV} (constat~\ref{c:stefan}).\\
\end{longtable}
\end{center}

% ============================================================================
\section{Constats de lecture}
\label{sec:constats}
% ============================================================================

Ces points découlent du source tel qu'il est écrit.
Aucun n'a été confirmé par un calcul : chacun mérite un jeu de données d'essai.

\subsection*{À risque}
\begin{enumerate}[label=\textbf{\arabic*.},leftmargin=2em]
\item\label{c:ktangent} \textbf{\color{risk}\src{K_TANGENT} efface l'inertie et la perméabilité.}
L'opérateur devient « blocages + \op{KTAN} » au début du sous-pas et à chaque itération. Les termes $\tfrac4{\Delta t^2}\vc M + \tfrac2{\Delta t}\vc C$ et $-\theta_c\Delta t\,\vc H$ disparaissent alors qu'ils restent dans le résidu ; c'est aussi le cas de \src{RIGIDITE_CONSTANTE} et de l'\src{ADDI_MATRICE}. En GD, la reconstruction de $\vc K$ perd aussi $\vc C$ et $\vc H$ ; l'auteur le note en commentaire (\lig{1721-1722}). Le résidu restant juste, une solution convergée n'est pas fausse ; mais l'opérateur, plus souple que la vraie jacobienne, peut faire osciller ou diverger les itérations quand le terme de masse domine (analyse, non testée).
\hfill{\footnotesize\pr{UNPAS} \lig{840-855}, 918-951, 1714-1724, 1836}

\item\label{c:visqueux} \textbf{\color{risk}Convergence forcée et lois visqueuses.}
Avec \src{RELAXATION_NONCONV} = 1, il ne reste plus de temps à intégrer après un échec. Pour \src{VISCOPLASTIQUE}, \src{FLUAGE}, \src{VISCODOMMAGE} et \src{VISCO_EXTERNE}, en dynamique ou en consolidation, cela donne l'erreur 1155. Une loi visqueuse déclarée \src{NON_LINEAIRE UTILISATEUR} ne lève pas d'erreur, mais reçoit $\Delta t = 0$.
\hfill{\footnotesize\pr{UNPAS} \lig{3670-3674} ; \pr{PAS\_MODL} \lig{179-222}}

\item\label{c:dupont} \textbf{\color{risk}\pr{DUPONT2} dégénère d'un pas à l'autre.}
\src{TET0} reçoit la température de fin de pas au lieu de $\vc T_{n-1}$. Le schéma à trois niveaux devient un $\theta$-schéma à $\theta = \tfrac34$, avec le matériau figé au début du pas (§ \ref{sec:dupont}). Après une reprise, le premier pas passe par \pr{TRANSNON}, car \src{VIEUXPAS} repart à 0.
\hfill{\footnotesize\pr{PASAPAS} \lig{509-514}, 598 ; \pr{DUPONT2} \lig{186-197}, 238-240, 485 ; \pr{PAS\_DEFA} \lig{572} ; \pr{PAS\_INIT} \lig{84}}

\item\label{c:sources} \textbf{\color{risk}Sources de chaleur ignorées hors \pr{TRANSNON}.}
Seul \pr{TRANSNON} lit le modèle \src{SOURCE}. \pr{PAS\_DEFA} ne bascule pas vers \pr{TRANSNON} quand il y a une source : avec \src{LINEAIRE} ou \src{DUPONT}, la source est ignorée sans message. \pr{TRANSLIN} ignore aussi le rayonnement.
\hfill{\footnotesize\pr{PAS\_MODL} \lig{86-89} ; \pr{PAS\_DEFA} \lig{1203-1206}}

\item\label{c:stefan} \textbf{\color{risk}Constante de Stefan-Boltzmann ignorée dans les cavités.}
Dans la méthode des cavités par température équivalente (sans \src{TRAYO}), \pr{PAS\_RAYO} appelle \pr{HRCAV} sans lui passer $\sigma_{SB}$ ; \pr{HRCAV} prend sa valeur SI. Un calcul en mm ou en ms qui fournit \src{CTE_STEFAN_BOLTZMANN} obtient un $h$ faux pour ces cavités. \src{INFINI}, \src{FAC_A_FAC} et \src{TRAYO} utilisent bien la valeur fournie.
\hfill{\footnotesize\pr{PAS\_RAYO} \lig{92}, 129-131, 158, 178, 190 ; \pr{HRCAV} \lig{6}, 20-22}
\end{enumerate}

\subsection*{À tester}
\begin{enumerate}[label=\textbf{\arabic*.},leftmargin=2em,resume]
\item\label{c:rrrr} \textbf{\color{know}Raideur avec relations : \src{RRRR} relue sans avoir été rangée.}
Si la raideur rendue par \op{RIGI} contient des multiplicateurs (\lig{731}), \pr{UNPAS} les ajoute aux blocages mais ne range pas la raideur dans \src{WTAB.RRRR} (\lig{765-767}) ; la ligne 771 la relit pourtant. Au premier pas d'un appel, l'indice n'existe pas, et Gibiane lève l'erreur 535 (§ \ref{sec:gibiane}). Les modèles qui produisent de telles relations restent à identifier.
\hfill{\footnotesize\pr{UNPAS} \lig{711}, 731-734, 765-771 ; \pr{PASAPAS} \lig{644-645}}
\end{enumerate}

\subsection*{À connaître}
\begin{enumerate}[label=\textbf{\arabic*.},leftmargin=2em,resume]
\item\label{c:maxiter} \textbf{\src{MAXITERATION} donné est écrasé.}
Dès que la convergence forcée est active (c'est le défaut), \pr{PAS\_DEFA} le met à 250, même si l'utilisateur a fourni une autre valeur.
\hfill{\footnotesize\pr{PAS\_DEFA} \lig{1001-1011}}

\item\label{c:plafond} \textbf{\src{MAXITERATION} n'est pas un plafond.}
L'atteindre déclenche la détection, qui triple la limite : un résidu qui baisse lentement mais régulièrement peut aller jusqu'à 750 itérations. La limite modifiée n'est pas réinitialisée au sous-pas suivant.
\hfill{\footnotesize\pr{UNPAS} \lig{147}, 1467-1472, 3439-3442, 3498-3506}

\item\label{c:deuxit} \textbf{Au moins deux itérations dès que le pas déforme.}
À l'itération 1, le critère en déformation vaut tout l'incrément : pas de convergence possible à l'itération 1 si $\norminf{\Deps} \ge 10^{-4}$, même en élasticité linéaire.
\hfill{\footnotesize\pr{UNPAS} \lig{688}, 1479, 2213, 3288-3334}

\item\label{c:ajuste} \textbf{\src{PAS_AJUSTE} ne voit pas les sous-pas ratés.}
\src{NOMBRE_ITERATIONS} ne compte que les itérations des sous-pas convergés : un pas sauvé par la convergence forcée, dont le dernier sous-pas converge vite, peut voir son $\Delta t$ augmenter.
\hfill{\footnotesize\pr{UNPAS} \lig{599}, 3420, 3814 ; \pr{PASAPAS} \lig{123-158}}

\item\label{c:forcee} \textbf{La convergence forcée est presque toujours active.}
La mise à faux par \src{PAS_AJUSTE} (\pr{PAS\_DEFA} \lig{732}) est écrasée à la ligne 978, et le pilotage garde la valeur vraie de la ligne 194. Les branches de \pr{PASAPAS} qui traitent un pas non convergé ne servent que si l'utilisateur désactive explicitement l'option ; dans ce cas, \pr{PAS\_RESU} propage l'état non convergé au pas suivant.
\hfill{\footnotesize\pr{PAS\_DEFA} \lig{194}, 729-734, 974-980 ; \pr{PAS\_RESU} \lig{19-42}}

\item\label{c:ktan8} \textbf{\op{KTAN} à $10^{-8}$ près.}
\pr{UNPAS} appelle \op{KTAN} avec $\mathrm{PREC} = \src{PRECISINTER}$. Pour les modèles de plasticité de \texttt{PENT1}, un point dont la contrainte équivalente dépasse le seuil de plus de $10^{-8}$ en relatif arrête \op{KTAN} sur l'erreur 275 ; un point plastique resté sous le seuil de plus de $10^{-8}$ reçoit la matrice élastique.
\hfill{\footnotesize\pr{UNPAS} \lig{149}, 162-163, 950, 1827 ; Fortran \texttt{KTANGA}, \texttt{DOHOT1}, \texttt{PENT1}}

\item\label{c:fdf} \textbf{Facteurs de forme figés en grands déplacements.}
\pr{PAS\_RAYO} ne recalcule les facteurs de forme des cavités qu'au premier appel ou si \src{REA_GEOM} est vrai, ce qui n'arrive qu'aux passes $\ge 2$ de \src{CONVERGENCE_MEC_THE} en GD. Avec une seule passe, une cavité qui se déforme garde ses facteurs de forme initiaux ; les matrices, elles, sont recalculées sur la géométrie courante.
\hfill{\footnotesize\pr{PAS\_RAYO} \lig{71-84} ; \pr{PASAPAS} \lig{363-369} ; \pr{PAS\_DEFA} \lig{503}}

\item\label{c:mat1} \textbf{Couplage itéré : le matériau « de début » change à la passe 2.}
\pr{UNPAS} écrit le matériau de fin de pas dans \src{WTAB.MAT1} (\lig{3792}) et le relit comme matériau de début (\lig{283}). Avec \src{CONVERGENCE_MEC_THE}, cela touche la correction « matériau variable », $\vc F^{int}_n$ des éléments de structure et l'état initial passé à \op{COMP}.
\end{enumerate}

\subsection*{Détails}
\begin{enumerate}[label=\textbf{\arabic*.},leftmargin=2em,resume]
\item\label{c:details} \textbf{Petites incohérences.}
\begin{itemize}
\item Avec $\alpha < 1$, la température de fin de sous-incrément n'est pas mise à l'échelle, contrairement au temps et à $\Delta\vg\eps^{th}$ (\pr{UNPAS} \lig{2484-2485}).
\item \op{SSTE} repart de $\vc V_n$ à chaque sous-incrément, dont seule la dernière composante est mise à jour (\lig{2458-2460}).
\item Le chemin \pr{DEFO\_IMP} est inactif : \src{LOG_CNTRL} vaut faux en dur (\lig{26}).
\item Thermique : $\theta = 1$ et le maximum d'itérations de la mécanique, au lieu des défauts annoncés par \pr{TRANSNON}.
\item \src{SOUS-RELAXATION} (\pr{TRANSLIN}) et \src{SOUS_RELAXATION} (\pr{DUPONT2}) ne sont jamais recopiés par \pr{PAS\_DEFA} : $\lambda_2$ vaut toujours 1 dans \pr{TRANSLIN} et $\tfrac12$ dans \pr{DUPONT2}, qui refuse donc tout matériau thermique dépendant de champs mécaniques (\pr{DUPONT2} \lig{280-287}).
\item Les commentaires de \pr{PAS\_RAYO} se contredisent sur les cavités : « méthode avec calcul de la matrice de rayonnement » est écrit au-dessus de la branche qui calcule une température équivalente, et « méthode par défaut » dans la branche \src{TRAYO} (\lig{172-187}). Seule la présence de \src{TRAYO} décide.
\item Dans \pr{TRANSNON}, si \src{CONDUCTIVITE_CONSTANTE} est la seule conduction, la ligne 806 écrit \src{KDETO = VRAI} au lieu de \src{IKDETO = VRAI} : la matrice est écrasée par un logique. Le cas suppose un modèle thermique sans conduction, donc il est sans doute théorique.
\end{itemize}
\end{enumerate}

\subsection*{Résolu}
\begin{enumerate}[label=\textbf{\arabic*.},leftmargin=2em,resume]
\item\label{c:translin} \textbf{\pr{TRANSLIN} lit le \src{WTAB} de \pr{PASAPAS}.}
Les lignes 87-88 lisent \src{WTAB.'CAPACONST'} alors que seule \src{ETAB} est définie. Par la portée dynamique (\texttt{INIPRO}), c'est le \src{WTAB} de \pr{PASAPAS} (\lig{27}), qui est la même table : ce n'est pas une erreur.
\end{enumerate}

% ============================================================================
\section{Comment ce texte a été vérifié}
% ============================================================================

\begin{itemize}
\item \textbf{Source Gibiane.} Les procédures du tableau~\ref{tab:versions} ont été lues ligne à ligne, quatre fois. Chaque expression traduite en formule l'a été en respectant l'évaluation de gauche à droite.
\item \textbf{Fortran des opérateurs.} Tout ce qui venait d'une notice a été revérifié dans le Fortran : interpréteur (\texttt{ANASYN}, \texttt{PILOT}, \texttt{LIRABJ}, \texttt{ECPI}, \texttt{PROCED}, \texttt{INIPRO}, \texttt{REMPIL}) ; \op{ACT3} ; \op{KTAN} (\texttt{KTANGE}, \texttt{KTANGA}, \texttt{DOHOT1}, \texttt{PENT1}, \texttt{KTAPER}) ; \op{CONF} (\texttt{CONFIG}, \texttt{PIOCAP}, \texttt{PIOCAF}, \texttt{JAUCAU}, \texttt{OPTION}) ; \op{EPSI} (\texttt{EPSI1}, \texttt{EPSI2}, \texttt{BST2}) ; \op{ECFE} (\texttt{FEFP}, \texttt{FEFP2}, \texttt{ADFEFP}) ; \op{SSTE} (\texttt{SBSTEP}) ; \op{NLOC} (\texttt{NLOCAL}, \texttt{NLOCA1}) ; \op{RAYE} (\texttt{PRRAYE}, \texttt{RAYEN}, \texttt{RAYT1}), \op{RAYN}, \op{FFOR} ; \op{RESO} (\texttt{RESOU}) ; \op{HOOK} ; \op{FORM} ; \op{RFCO}.
\item \textbf{Écarts avec la documentation.} La notice de \op{KTAN} annonce un retour à l'élastique en cas d'échec, que le Fortran ne fait pas ; celle de \op{CONF} décrit un transport, sans dire que c'est une poussée $J^{-1}\vc F\sig\vc F^{\mathsf T}$.
\item \textbf{Non lu.} La correction des vitesses de contact en dynamique (\pr{VITEUNIL}, \pr{VITETFOR}), la parallélisation (\op{ASSIS}), et l'intérieur de \op{COMP} loi par loi.
\item \textbf{Non testé.} Aucune affirmation n'a été confirmée par un calcul.
\end{itemize}

\end{document}
